\documentclass[10pt,a4paper]{article}
\usepackage{fontspec,xeCJK}
\setmainfont{Times New Roman}
\setCJKmainfont[AutoFakeBold=2,AutoFakeSlant=.2]{SimSun}
\setCJKmonofont{SimSun}
\usepackage[margin=20mm,top=22mm,bottom=22mm]{geometry}
\usepackage{paracol,graphicx,hyperref,xurl,fancyhdr,needspace,xcolor}
\hypersetup{unicode=true,hidelinks}\urlstyle{same}
\setlength{\columnsep}{8mm}\setlength{\parindent}{0pt}\setlength{\parskip}{5pt}
\setlength{\emergencystretch}{3em}\tolerance=2000
\widowpenalty=5000\clubpenalty=5000
\pagestyle{fancy}\fancyhf{}\fancyhead[L]{\small Lacchei \& Reyhani · From Borders to Courts}\fancyfoot[C]{\thepage}\setlength{\headheight}{14pt}
\begin{document}

{\Large\bfseries 从边界到法院：在充满阻碍的跨国环境中组装推回案件\par}\medskip
From borders to courts: assembling pushback cases in a hostile transnational setting\par
Alice Lacchei：意大利博洛尼亚，博洛尼亚大学政治与社会科学系。Adel-Naim Reyhani：奥地利维也纳，路德维希·玻尔兹曼基本人权研究所。通讯作者：Alice Lacchei；alice.lacchei2@unibo.it。《法律与社会评论》（2026），第1—27页。收稿：2025年10月29日（UTC）；修订：2026年3月26日（UTC）；录用：2026年6月26日（UTC）。在线发表：2026年9月18日。\par
据 Cambridge University Press 出版方 HTML 全文翻译，并以同文原版 PDF 核对结构及段落起始页码。原文27页，保留摘要、全部正文、数据与资助声明、9项访谈清单、91条文献、7项判例、3项法规及两位作者简介。无编号脚注或图表。书目照录并附中文题名或条目译释。© The Author(s), 2026，Cambridge University Press 代表 Law and Society Association 出版；原文与本译文采用 CC BY 4.0。翻译、译释及编者说明为本版新增。AI辅助工作译稿，不代表作者或出版社认可。\par
\url{https://doi.org/10.1017/lsr.2026.10117}\par
\url{https://creativecommons.org/licenses/by/4.0/}\par
关键词：推回；中介者；法院；难民；正义\par\bigskip
\Needspace{10\baselineskip}\section*{摘要 / Abstract}
\begin{paracol}{2}
{\scriptsize\color{gray} [p. 1]}\ 本文考察法律中介者将案件从边界带入法院的工作，聚焦从意大利、奥地利向斯洛文尼亚，以及沿巴尔干通道发生的推回。文章表明，这些做法发生在一种充满阻碍的跨国环境中；这种环境从一开始就使潜在诉请人、痕迹和诉讼路径分散开来，并使其处于不稳定状态。文章指出，面对这种初始的错位状态，中介者开展了跨越国界的构成性与协作性工作。本文将这种工作概念化为“案件组装”，并从比较视角加以分析。\par
\switchcolumn
{\scriptsize\color{gray} [p. 1]}\ This article examines the work of legal intermediaries in moving cases from the border to the court, focusing on pushbacks from Italy and Austria to Slovenia and along the Balkan corridor. It demonstrates that these practices occurred within a hostile transnational setting that, from the outset, dispersed and destabilized potential claimants, traces, and the litigation path. Confronted with this immediate condition of misalignment, it shows that intermediaries engaged in constitutive and collaborative work across borders. The article conceptualizes this work as “case assembly” and analyzes it from a comparative perspective.\par
\switchcolumn*
\end{paracol}
\Needspace{10\baselineskip}\section*{引言 / Introduction}
\begin{paracol}{2}
{\scriptsize\color{gray} [p. 1]}\ 法院程序只是法律纠纷中可见的冰山一角（Felstiner et al. 1981）。法律与社会研究长期考察为何大多数潜在诉讼当事人始终未能进入司法阶段，指出有必要同时研究一个问题成为法院案件之前所发生的事情，尤其包括政治、文化与法律环境，以及参与其中的行动者（Olesen and Hammerslev 2023）。\par
\switchcolumn
{\scriptsize\color{gray} [p. 1]}\ Court proceedings represent only the visible tip of the iceberg in legal disputes (Felstiner et al. 1981). Socio-legal scholarship has largely investigated the reasons why most potential litigants never reach the judicial stage, suggesting that it is necessary to investigate also what occurs before a problem becomes a court case, including particularly the political, cultural, and legal environment, as well as the actors involved (Olesen and Hammerslev 2023).\par
\switchcolumn*
{\scriptsize\color{gray} [p. 1]}\ 这一问题在移民情境中更加突出（Conant 2016），在关乎能否获得庇护的案件中尤为尖锐：司法问责日益受到限制，法院只能处理极小部分伤害（Costello and Mann 2020）。事实上，当代移民治理中，世界各国为进入本国领土和庇护程序设置了各种障碍（Lacchei and Moraru 2026），以规避法律、逃避问责（Gammeltoft-hansen and Tan 2021; Moreno-lax 2024）。为此，它们将边界扩展成碎片化且往往无形的空间（Tazzioli 2021）；“移动的边界”现象（Shachar 2020）及外部化政策的兴起就是例证。这些政策操纵国家控制的方式，甚至可能使人陷入无权利状态（Reyhani 2021）。\par
\switchcolumn
{\scriptsize\color{gray} [p. 1]}\ This issue is sharpened in the context of migration (Conant 2016) and becomes even more acute in cases related to access to asylum, in which judicial accountability is increasingly constrained, and courts can only address a very small fraction of harm (Costello and Mann 2020). Indeed, in contemporary migration governance, states across the globe are erecting diverse barriers to accessing their territories and asylum procedures (Lacchei and Moraru 2026) with a view to circumventing the law and avoiding accountability (Gammeltoft-hansen and Tan 2021; Moreno-lax 2024). For this purpose, they have expanded borders into fragmented, often intangible spaces (Tazzioli 2021), as illustrated by the phenomenon of shifting borders (Shachar 2020) and the rise of externalization policies, which manipulate state control in ways that can even produce rightlessness (Reyhani 2021).\par
\switchcolumn*
{\scriptsize\color{gray} [p. 2]}\ 法院是对这些做法提出挑战的问责场域之一。此外，即使选择了法律行动，行动者也并非总认为法律能够提供充分或持续的救济（Costello and Mann 2020; Cummings 2020; Kawar 2015; Sterett 2016）。事实上，法院介入移民管控领域时，国家可能重新调整管控做法，并将其进一步外部化，使法律与国家的适应行为之间形成猫鼠博弈（Lavenex 2022; Mann 2013）。\par
\switchcolumn
{\scriptsize\color{gray} [p. 2]}\ Courts are one accountability venue among others to contest such practices. Moreover, even when legal action is chosen, actors do not always assume that the law will provide sufficient or sustained remedy (Costello and Mann 2020; Cummings 2020; Kawar 2015; Sterett 2016). Indeed, where courts intervene in the field of migration control, states may rework control practices and externalize them even further in ways producing a cat-and-mouse dynamic between the law and state adaptation (Lavenex 2022; Mann 2013).\par
\switchcolumn*
{\scriptsize\color{gray} [p. 2]}\ 从法律与社会研究的视角看，研究庇护准入诉讼尤其有价值。移民是法院与法律程序在包容和排斥的治理中居于核心位置的领域，即便形式上得到保障的权利，在实践中仍难以实现（Sterett 2016）。这一领域中诉诸法院的机会，还受到制度、程序路径及资源不平等的塑造（Campbell 2020; Passalacqua 2022）。这一情境可以揭示，为抵达司法机构所必须开展的工作有多么密集、复杂和交错（Crooke 2024; Jacobs and Maryns 2022）。学者指出，移民案件偶尔得以进入法院时，通常是律师、行动者和公民社会组织等多方努力的结果（Riva and Hoffstaedter 2021; Sterett 2016）。研究对这些行动者有不同称呼。本文使用“法律中介者”一词，指那些将个人与司法机构联系起来、努力推动具体诉请进入处理程序的法律及非法律专业人员（Lejeune and Spire 2021）。\par
\switchcolumn
{\scriptsize\color{gray} [p. 2]}\ From a socio-legal perspective, studying litigation over access to asylum is particularly valuable. Migration is a field in which courts and legal processes centrally govern inclusion and exclusion, even where formally guaranteed rights remain difficult to realize in practice (Sterett 2016). Access to courts in this field is further shaped by institutions, procedural routes, and unequal resources (Campbell 2020; Passalacqua 2022). This context can elucidate the intensity and complex or intricate nature of the work required to reach judicial institutions (Crooke 2024; Jacobs and Maryns 2022). Scholars note that, when migration cases occasionally do reach the courts, it is usually the result of efforts by a variety of actors, including lawyers, activists, and civil society organizations (Riva and Hoffstaedter 2021; Sterett 2016). They have referred to these actors in various terms. Here, we use the term legal intermediaries – i.e., those legal and nonlegal professionals who connect individuals with judicial institutions, seeking to move a concrete claim toward processing (Lejeune and Spire 2021).\par
\switchcolumn*
{\scriptsize\color{gray} [p. 2]}\ 在此背景下，我们通过研究诉讼条件异常不利的领域——巴尔干路线沿线的推回——来增进对这些中介者工作的理解。推回指未经事先评估风险或保护需求，就将个人强行移出国家领土；按其定义，它从一开始就是跨国的。此外，在巴尔干路线沿线，推回往往形成跨境链条，以非正式警务合作和重新接收安排为基础，削弱问责并利用法律的不透明性（Cantor et al. 2022; Casas-cortes et al. 2016; Moraru and Reyhani 2026）。作为一种诉讼环境，这一领域还具有不透明、难以追溯及证据反复遭到抹除的特点（Alpes and Baranowska 2025; Baranowska et al. 2025）。最后，移民的刑事化，以及有时对声援行为本身的刑事化，也塑造着这一领域（Trucco 2023）。\par
\switchcolumn
{\scriptsize\color{gray} [p. 2]}\ Against this backdrop, we seek to advance understanding of the work of these intermediaries by studying a field marked by unusually adverse litigation conditions: pushbacks along the Balkan route. Understood as the forced removal of individuals from State territory without a prior assessment of risks or protection needs, pushbacks are by definition and from the start transnational. Moreover, along the Balkan route, they often come with cross-border chains based on informal police cooperation and readmission arrangements that undermine accountability and leverage legal opacity (Cantor et al. 2022; Casas-cortes et al. 2016; Moraru and Reyhani 2026). As a litigation setting, the field has moreover been characterized by opacity, weak traceability, and recurrent problems of evidentiary erasure (Alpes and Baranowska 2025; Baranowska et al. 2025). Finally, it has been shaped by the criminalization of migration and, at times, of solidarity itself (Trucco 2023).\par
\switchcolumn*
{\scriptsize\color{gray} [p. 2]}\ 在这些认识基础上，本研究表明，推回是一种充满阻碍的跨国环境的组成部分。这种环境从一开始，就使案件的基本要素——潜在诉请人、痕迹与诉讼路径——分散开来，并使其不稳定。本文认为，面对这种最初的\textit{案件错位}状态，中介者以跨国界的构成性与协作性工作作出回应。我们将这一工作概念化为\textit{案件组装}，即识别、连接和稳定案件要素，以推动案件向法院移动的过程。\par
\switchcolumn
{\scriptsize\color{gray} [p. 2]}\ Building on these insights, our study demonstrates that pushbacks are part of a hostile transnational setting that, from the very outset, disperses and destabilizes the basic components of cases, namely potential claimants, traces, and the litigation path. Confronted with this initial condition of \textit{case misalignment}, it argues, intermediaries respond with constitutive and collaborative work across borders. We conceptualize this work as \textit{case assembly} – i.e., the process of identifying, connecting, and stabilizing the elements of cases to advance their movement toward the court.\par
\switchcolumn*
{\scriptsize\color{gray} [p. 2]}\ 我们采取比较研究的方法，考察意大利与奥地利国内法院审理的推回案件；这些案件均涉及向斯洛文尼亚的推回，且往往随后沿巴尔干路线被进一步移送。如“案例研究”一节所述，两种情境中的法律中介者都曾开展战略诉讼，促使法院谴责相关做法。此外，两国虽然共享同一跨国环境，却采用不同做法和理由，制度与程序配置也各不相同。\par
\switchcolumn
{\scriptsize\color{gray} [p. 2]}\ Our approach is comparative. We investigate cases of pushbacks litigated before domestic courts in Italy and Austria, each concerning pushbacks to Slovenia and often followed by onward removals further along the Balkan route. As elaborated in the “Case studies” “section, in both contexts, legal intermediaries had pursued strategic litigation that led courts to condemn the relevant practices. Moreover, while sharing the same transnational setting, the two countries relied on different practices and justifications and are characterized by distinct institutional and procedural configurations.\par
\switchcolumn*
{\scriptsize\color{gray} [p. 3]}\ 既有学术研究尚未探讨中介者的这种具体工作。不过，我们可以借助已有研究，因为它们已阐明了这一领域的若干部分。尤其是，对组织及法律—行政环境中法律中介活动的研究，说明中介者如何撮合准入、塑造法律意义，并引导人们通过法律—行政路径（Bonizzoni and Odasso 2024; Pélisse 2019, 2019; Talesh 2015）。此外，关于纠纷金字塔及纠纷前阶段的研究，揭示被识别出的冲突如何走向、或未能走向法律求助，以及中介者如何在这一过程中工作（Lejeune and Spire 2021; Olesen and Hammerslev 2023）。有关移民权益倡导与庇护行政的研究更进一步，说明叙事与证据材料如何经过转译，成为法律领域可识别的记录（Cabot 2012; Coutin 2025; Coutin and Fortin 2023; Lakhani 2014; Pörtner 2021）。\par
\switchcolumn
{\scriptsize\color{gray} [p. 3]}\ Scholarship has not yet engaged with this specific work of intermediaries. Nevertheless, our study can draw on existing work that already illuminates several parts of the terrain. In particular, research on legal intermediation in organizational and legal-administrative settings shows how intermediaries broker access, shape legal meaning, and route people through legal-administrative pathways (Bonizzoni and Odasso 2024; Pélisse 2019, 2019; Talesh 2015). Moreover, work on the dispute pyramid and the pre-dispute phase shows how recognized conflicts move, or fail to move, toward legal recourse and how intermediaries work in that progression (Lejeune and Spire 2021; Olesen and Hammerslev 2023). Work on immigration advocacy and asylum administration goes further by showing how narratives and evidentiary materials are translated so that they become records cognizable in the legal sphere (Cabot 2012; Coutin 2025; Coutin and Fortin 2023; Lakhani 2014; Pörtner 2021).\par
\switchcolumn*
{\scriptsize\color{gray} [p. 3]}\ 本文结构如下：“充满阻碍的跨国环境”一节解释这一跨国环境；“法律中介者及其工作”一节对法律中介者的工作加以概念化，并将本文置于现有文献中；“案例研究”一节介绍案例及其选择理由；“研究方法”一节说明方法、范围与限制；“经验材料”一节呈现有关中介者工作的经验数据；最后，在结论之前，“讨论”一节讨论所得认识，以及它们对理解法律中介活动的意义。\par
\switchcolumn
{\scriptsize\color{gray} [p. 3]}\ The article is structured as follows. “A hostile transnational setting” section explains the hostile transnational setting. “Legal intermediaries and their work” section conceptualizes the work of legal intermediaries and situates the article in the available literature. “Case studies” “Case studies”section introduces the case studies and their rationale. “Methods” section outlines the methodology, its scope, and limitations. “Empirical data” section presents empirical data on intermediaries’ work. Finally, before concluding, the “Discussion” section discusses the insights gained and their implications for understanding legal intermediation.\par
\switchcolumn*
\end{paracol}
\Needspace{10\baselineskip}\section*{充满阻碍的跨国环境 / A hostile transnational setting}
\begin{paracol}{2}
{\scriptsize\color{gray} [p. 3]}\ 按定义，推回是跨国的：它涉及将人移送过国际边界，而不允许其进入庇护程序，也不对其保护需求作个别评估（UN Special Rapporteur on the human rights of migrants 2021）。推回至少连接着实施移送的国家与人们被送往的国家。在当代移民治理中，涉及的范围往往更加广泛。边境管控日益延伸至领土之外，推回的跨国范围也随之扩大。包括非正式跨境警务合作在内的外部化措施，沿迁徙路线形成一条推回链（Giuffré 2020; Legomsky 2003）。巴尔干路线就是这一现象的典型场域。沿线各国的合作，促成了跨越多个国家的反复向外移送，公民社会与监测行动者对此已有大量记录（Aulsebrook et al. 2021; Border Violence Monitoring Network 2024）。在此环境中，难民斗争、声援行动、监测网络、新闻调查，以及多样的记录和反档案实践，都对推回提出挑战；它们试图为不同受众和政治目的，让边境暴力变得可见，其中只有部分后来进入诉讼（Cantat 2020; Davies et al. 2023; El-shaarawi and Razsa 2019; Hameršak et al. 2020）。\par
\switchcolumn
{\scriptsize\color{gray} [p. 3]}\ Pushbacks are transnational by definition. They involve removal across an international border without access to asylum procedures or an individual assessment of protection needs (UN Special Rapporteur on the human rights of migrants 2021). At a minimum, a pushback links the state carrying out the removal and the state to which people are sent. In contemporary migration governance, the setting is often more extended. As border control becomes increasingly extraterritorial, the transnational scope of pushbacks expands accordingly. Externalization efforts, including informal cross-border police cooperation, are producing a chain of pushbacks along migratory routes (Giuffré 2020; Legomsky 2003). The Balkan route is a typical terrain for this phenomenon. Cooperation among states along the route has contributed to repeated onward removals across several states, as documented extensively by civil society and monitoring actors (Aulsebrook et al. 2021; Border Violence Monitoring Network 2024). In this setting, pushbacks are contested through refugee struggles, solidarity, monitoring networks, journalistic investigation, and diverse practices of documentation and counter-archiving, which seek to make border violence visible for multiple audiences and political purposes, only some of which later travel into litigation (Cantat 2020; Davies et al. 2023; El-shaarawi and Razsa 2019; Hameršak et al. 2020).\par
\switchcolumn*
{\scriptsize\color{gray} [p. 3]}\ 推回诉讼是这一更广泛抗争生态的组成部分，处于范围更广、充满阻碍的跨国环境中。而这种环境，确实影响着其中的法律与诉讼动态。\par
\switchcolumn
{\scriptsize\color{gray} [p. 3]}\ The litigation over pushbacks, as part of this wider ecosystem of contestation, is situated within a broader, hostile, and transnational environment. And, indeed, this setting has consequences for the legal and litigation dynamics at play.\par
\switchcolumn*
{\scriptsize\color{gray} [p. 4]}\ 有关外部化的研究表明，流动中的人与国家相遇的位置发生了转移，由此提出一个问题：当管控被外包到领土边界之外时，法律义务是否仍随移民治理一道延伸（Gammeltoft-Hansen 2011）？问责可能在不同法律制度与行动地点之间碎裂，使责任更难界定，也更难追究（Tan and Gammeltoft-Hansen 2020）。在向解体中的国家进行更严重外部化的情况下，法律有时可能完全无法使用或不适用（Reyhani 2021）。法院介入时，国家可能重新调整管控做法，将其进一步外部化，以维持遵法的表象，同时削弱法律约束的实际效力，从而在法律判决与国家的适应行为之间形成猫鼠博弈（Lavenex 2023; Mann 2013）。\par
\switchcolumn
{\scriptsize\color{gray} [p. 4]}\ Scholarship on externalization has shown how the encounter of people on the move with states is displaced, raising the question of whether legal obligations continue to travel with migration governance as control is outsourced beyond the territorial border (Gammeltoft-Hansen 2011). Accountability can become fractured across legal regimes and operational sites, making responsibility harder to map and harder to claim (Tan and Gammeltoft-Hansen 2020). Under more severe forms of externalization to disintegrating states, the law may at times become entirely unavailable or inapplicable (Reyhani 2021). Where courts intervene, states may rework control practices and externalize them even further in ways that preserve the appearance of compliance while weakening the practical force of legal constraint, producing a cat-and-mouse dynamic between legal judgment and state adaptation (Lavenex 2023; Mann 2013).\par
\switchcolumn*
{\scriptsize\color{gray} [p. 4]}\ 在诉讼中，推回的跨国性质还带来要素分散这一实际问题。在推回案件里，诉请人立即与事件发生地，以及可能确立问责的国内司法场域断开联系。同样，相关事件可能横跨数国，而痕迹或证据又可能位于另一个司法管辖区。更进一步，侵权行为必须在不透明、暴力与制度非正式性突出的边境空间中被发现、记录和追究（Fortarezza 2026; Hameršak et al. 2020）。最后，移民的刑事化，以及有时对声援行为本身的刑事化，使这些努力进一步复杂化，限制了接触流动中人群的机会，也提高了监测、记录与支持的成本（Kogovšek Šalamon 2020; Trucco 2023）。\par
\switchcolumn
{\scriptsize\color{gray} [p. 4]}\ In litigation, the transnationality of pushbacks is also a practical concern for dispersal. In pushback cases, the claimant is immediately disconnected from the site of the event and the domestic judicial forum that could establish accountability. Similarly, the relevant event may stretch across several states, and traces or evidence may yet be located in another jurisdiction. What is more, violations have to be noticed, recorded, and pursued across border spaces marked by opacity, violence, and institutional informality (Fortarezza 2026; Hameršak et al. 2020). Finally, efforts are further complicated by the criminalization of migration and, at times, of solidarity itself, limiting access to people on the move and raising the costs of monitoring, documentation, and support (Kogovšek Šalamon 2020; Trucco 2023).\par
\switchcolumn*
{\scriptsize\color{gray} [p. 4]}\ 尚在兴起的有关推回诉讼的法律与社会研究，迄今主要聚焦国际司法救济；在这类文献中，上述特征以程序形式出现：证据不对称、难以追溯，以及否认或模糊事件本身、压缩公民社会行动空间、妨碍司法挑战的国家做法。Alpes（2026）表明，恰恰在国家既不记录推回、也不允许记录推回之时，欧洲人权法院（ECtHR）的证据制度往往优先采信国家生产的事实。近期分析进一步表明，“不留纸面记录的推回”、没收手机、限制进入边境地区，以及策略性的不合作，都可能使关键事实无法为欧洲人权法院所见；该院本身也已承认这一点（\textit{A.R.E. and G.R.J. v. Greece}）（Baranowska et al. 2025）。更广泛地说，这些文献把诉讼放在其所处的更大政治环境中，揭示该环境如何塑造能够抵达法院的内容（Alpes and Baranowska 2025; Cusumano and Villa 2021; Pijnenburg and Van der Pas 2022; Trucco 2023）。它们追问，为何欧洲边境暴力的大量记录，只有如此少的一部分反映在欧洲人权法院的判决中，由此也表明，通过监测网络在实地汇集的反档案，很少呈现在法院面前（Baranowska et al. 2025）。由此可得一个关键区别：可见性不等于法律上的可用性。边境暴力可能得到大量记录，并为公众所见，却仍无法以程序上可使用、或证据上足够充分的形式出现在法院面前（Alpes and Baranowska 2025; Baranowska et al. 2025; Davies et al. 2023）。\par
\switchcolumn
{\scriptsize\color{gray} [p. 4]}\ In the still emerging socio-legal literature on pushback litigation, which has so far focused mainly on international judicial remedies, these features appear in procedural form: evidentiary asymmetries, weak traceability, and state practices that deny or blur the event itself, narrow the space for civil society action, and impede judicial contestation. Alpes (2026) shows that the evidentiary regime of the European Court of Human Rights (ECtHR) often privileges state-produced facts precisely where states neither record nor allow the recording of pushbacks. Recent analysis further shows that “non-paper pushbacks,” phone confiscation, restricted access to border zones, and strategic non-cooperation can render key facts invisible to the ECtHR, as the Court has itself recognized (\textit{A.R.E. and G.R.J. v. Greece}) (Baranowska et al. 2025). More broadly, this literature places litigation within its wider political environment and shows how that environment shapes what can reach the court (Alpes and Baranowska 2025; Cusumano and Villa 2021; Pijnenburg and Van der Pas 2022; Trucco 2023). By asking why so little of the extensive documentation of European border violence is reflected in ECtHR judgments, it also shows that the counter-archive assembled on the ground through monitoring networks rarely arrives before the court (Baranowska et al. 2025). What follows is a central distinction: visibility is not the same as legal usability. Border violence may be extensively documented and made publicly visible, yet still fail to appear before a court in forms that are procedurally usable or evidentially sufficient (Alpes and Baranowska 2025; Baranowska et al. 2025; Davies et al. 2023).\par
\switchcolumn*
{\scriptsize\color{gray} [p. 4]}\ 诉请人的位置、证据与程序路径之间存在某种不匹配，在诉讼中并不少见。伤害经历本身阻碍后续问责，也并不罕见。然而，这里显现的是：在充满阻碍的跨国环境中，案件基本要素跨越国界分散开来，并变得不稳定。这种状态，即案件错位，正是由受到挑战的那套制度本身造成的。\par
\switchcolumn
{\scriptsize\color{gray} [p. 4]}\ Some degree of mismatch between the position of claimants, the evidence, and the procedural route is not unusual in litigation. Nor is it unusual that the experience of harm itself can obstruct later accountability. What becomes visible here, however, is the dispersal and destabilization of basic case elements across borders within a hostile transnational setting. This condition, namely case misalignment, is generated by the very regime being contested.\par
\switchcolumn*
\end{paracol}
\Needspace{10\baselineskip}\section*{法律中介者及其工作 / Legal intermediaries and their work}
\begin{paracol}{2}
{\scriptsize\color{gray} [p. 5]}\ 如果充满阻碍的跨国条件使推回案件赖以成立的要素分散、削弱，那么，究竟由谁、通过何种工作，使这些案件首先成为可能，并得以向前推进？我们将其视为法律动员早期、提出诉请之前阶段的法律中介问题。不过，本文关注点的特殊性，要求与其他若干文献领域建立联系，并使它们展开对话。\par
\switchcolumn
{\scriptsize\color{gray} [p. 5]}\ If hostile transnational conditions scatter and weaken the elements on which a pushback case depends, by whom and through what work do such cases first become possible and moveable at all? We see this as a question of legal intermediation in the early pre-claim phase of legal mobilization. However, the specificity of this article’s focus requires engagement with and dialogue between a set of further bodies of literature.\par
\switchcolumn*
{\scriptsize\color{gray} [p. 5]}\ 第一个出发点，是有关纠纷产生与转变的研究。学术研究早已表明，纠纷通过具有选择性、反复进行的命名、归责与索求过程形成；其轨迹不断分叉，而非遵循单一线性路径（Albiston et al. 2014; Felstiner et al. 1981; Miller and Sarat 1980–1981; Olesen and Hammerslev 2023）。有关语言与受众的研究同样指出，纠纷在不同机构与公众面前被框定时，会被重新塑造（Mather and Yngvesson 1981）。这些文献共同确立了诉请提出之前阶段的重要性，既包括必须被启动的法律机构，也包括介入这一阶段的第三方；后者被分别称为经纪者、守门人、转译者和中介者（Olesen and Hammerslev 2023）。这些研究还强调，此类介入具有权力不对称和矛盾效应：它们可能促进人们进入法律机构，同时也可能再生产排斥机制（Albiston et al. 2014; Michelson 2006）。不过，文献尚未精确说明的是，当一个问题的基本要素分散在不同地点与机构之间时，它究竟通过何种机制获得足够的连贯性，从而可以付诸法律行动。\par
\switchcolumn
{\scriptsize\color{gray} [p. 5]}\ A first point of departure is work on the emergence and transformation of disputes. Scholarship has long shown that disputes emerge through selective and iterative processes of naming, blaming, and claiming, and that their trajectories branch rather than follow a single linear path (Albiston et al. 2014; Felstiner et al. 1981; Miller and Sarat 1980–1981; Olesen and Hammerslev 2023). Work on language and audience likewise shows how disputes are reshaped as they are framed before different institutions and publics (Mather and Yngvesson 1981). Together, this literature establishes the importance of the pre-claim phase, including both the legal institutions that must be activated and the third parties who intervene within this phase, variously described as brokers, gatekeepers, translators, and intermediaries (Olesen and Hammerslev 2023). It also highlights the power asymmetries and contradictory effects that mark such interventions, which may facilitate access to legal institutions while also reproducing exclusionary dynamics (Albiston et al. 2014; Michelson 2006). What it leaves less precisely specified, however, is the mechanism by which a problem acquires enough coherence to become actionable when its basic elements are dispersed across sites and institutions.\par
\switchcolumn*
{\scriptsize\color{gray} [p. 5]}\ 另一组相关研究考察法律转译，说明经历及支持材料如何经过改造，成为机构可以处理的形式（Coutin 2025; Coutin and Fortin 2023; Cunningham 1992; Lakhani 2014）。有关庇护证据与行政的研究从另一角度指出，证据效力取决于制度形式，以及经授权或获得接受的类别与核验方式（Cabot 2012; Fassin and D’halluin 2005; Pörtner 2021）。类似地，Latour（2010）表明，法律推理依赖于已经按法律用途整理好的材料，以及将文书与法律类别、论证和先例联系起来的脆弱文本连接（Latour 2010）。这些研究对本文至关重要。同时，其中许多研究只有在法律工作已经具备可行基础之后才开始考察，因此，对先于转译工作、并使转译成为可能的条件，关注仍不充分。\par
\switchcolumn
{\scriptsize\color{gray} [p. 5]}\ A related set of works examines legal translation by showing how experiences and supporting materials are reworked into forms that institutions will work with (Coutin 2025; Coutin and Fortin 2023; Cunningham 1992; Lakhani 2014). Research on asylum evidence and administration suggests, from a different angle, that evidentiary force depends on institutional forms, and authorized or accepted categories and modes of verification (Cabot 2012; Fassin and D’halluin 2005; Pörtner 2021). Similarly, Latour (2010) shows how legal reasoning depends on materials already profiled for legal use and on the fragile textual connections through which documents are related to legal categories, arguments, and precedents (Latour 2010). This body of scholarship is key to the present article. At the same time, much of it begins only once a viable basis for legal work is already in place, thus not yet sufficiently engaging with the conditions that precede and enable translation work.\par
\switchcolumn*
{\scriptsize\color{gray} [p. 5]}\ 我们的分析主要通过法律中介活动这一视角展开。由于我们关心的是：在一个高度错位的案件最初阶段，哪些行动者及实际工作能够把其要素组装起来，同时又须关注连接、路径引导和选择性案件建构的问题，因此，这里从功能而非职业角度理解法律中介活动。所谓法律中介者，是那些通过连接个人与司法机构，并且连接个人、痕迹、文书及程序上的可能性，帮助潜在案件走向法院的人。本研究在操作上将法律和非法律专业人员都纳入这一术语，因为问题在于，哪些工作与过程直接关系到潜在推回案件能否推进。在本文情境中，这不仅包括律师，也包括行动者、记者、非政府组织（NGO）和监测行动者。因此，研究仍只聚焦部分行动者与活动。照护、声援或监测实践，只有在直接影响潜在案件向法院推进时，才纳入研究。由此也可知，不能把这里考察的行动者视为整个反推回抗争的代表，更不能视为遭受推回者的取向或能力的代表。相反，他们所处的位置，使他们只能将部分伤害引向裁判（Felstiner et al. 1981; Lejeune and Spire 2021; Miaz et al. 2021; Trucco et al. 2023）。\par
\switchcolumn
{\scriptsize\color{gray} [p. 5]}\ Our analysis proceeds primarily through the lens of legal intermediation, understood here in functional rather than professional terms, given the interest in understanding the actors and practical work that can assemble the elements of an intensely misaligned case at its very initial stage, keeping in view problems of connection, routing, and selective case-making. Legal intermediaries are understood, then, as those who help move a potential case toward court by connecting not only people and judicial institutions but also people, traces, and documents, as well as procedural possibilities. The term is here operationalized to include both legal and nonlegal professionals, since the question is which work and processes directly bear on whether a potential pushback case can advance. In the present setting, this includes not only lawyers but also activists, journalists, nongovernmental organizations (NGOs), and monitoring actors. The focus, therefore, remains on a subset of actors and activities. Care, solidarity, or monitoring practices are included only when they directly influence the movement of a potential case toward court. It also follows, then, that the actors examined here cannot be treated as representative of anti-pushback contestation as a whole, still less of the orientations or capacities of people subjected to pushbacks. Rather, they occupy a position where they direct only some harms toward adjudication (Felstiner et al. 1981; Lejeune and Spire 2021; Miaz et al. 2021; Trucco et al. 2023).\par
\switchcolumn*
{\scriptsize\color{gray} [p. 6]}\ 现有中介活动研究已经阐明若干上游过程，尽管它们并不处于同一分析层次：从冲突走向法律求助并进入纠纷金字塔的过程（Lejeune and Spire 2021）、居间撮合和法律—行政程序引导，以及与法律相关的类别、路径和记录的塑造（Billows et al. 2019; Pélisse 2019; Talesh 2015）。这些文献共同尚未充分说明的，仍是这样一个连接问题：在跨国碎片化条件下，一个可能的案件仅以分散形式存在，此时应如何将其联结起来。\par
\switchcolumn
{\scriptsize\color{gray} [p. 6]}\ Existing scholarship on intermediation already illuminates several upstream processes, though not at the same analytical level: movement from conflict toward legal recourse and into the dispute pyramid (Lejeune and Spire 2021), brokerage and legal-administrative navigation, and the shaping of legally relevant categories, routes, and records (Billows et al. 2019; Pélisse 2019; Talesh 2015). What remains less clearly specified across these literatures, again, is the linkage problem that arises when a possible case exists only in disaggregated form under conditions of transnational fragmentation.\par
\switchcolumn*
{\scriptsize\color{gray} [p. 6]}\ 移民研究补充了这幅图景，表明进入移民法律与行政体系通常需要中介；中介有时打开路径，有时再生产排斥（Bonizzoni and Odasso 2024; Bonjour 2025; Johan et al. 2012; Könönen 2018; Miaz et al. 2021; Tuckett 2018）。这些研究的价值在于，它们确立了移民领域由一系列中介链条组织起来这一点。不过，其重心通常仍在其他方面：居间撮合、行政程序引导、法律—行政中介活动、移民治理的日常生产，以及承担中介工作者充满矛盾的政治能动性。而中介者将流动中的人与法院连接起来的作用，在此类研究中尚缺乏考察。\par
\switchcolumn
{\scriptsize\color{gray} [p. 6]}\ Migration scholarship complements this picture by showing that access to migration law and administration is routinely mediated, sometimes opening routes and sometimes reproducing exclusion (Bonizzoni and Odasso 2024; Bonjour 2025; Johan et al. 2012; Könönen 2018; Miaz et al. 2021; Tuckett 2018). This scholarship is valuable as it establishes that the migration field is structured by chains of mediation. Yet, again, its center of gravity usually lies elsewhere, in brokerage, administrative navigation, legal-administrative intermediation, the everyday production of migration governance, and the ambivalent political agency of those who mediate it. The role of intermediaries in connecting people on the move with courts, however, is understudied here.\par
\switchcolumn*
{\scriptsize\color{gray} [p. 6]}\ 另一组相关研究探讨，为何即使法律行动者已经在场，走向法院仍然困难；这有助于理解特定环境中的中介功能。有关跨国法律动员与社会运动律师工作的研究表明，法律抗争依赖法律、社会运动与知识生产行动者之间跨地点的关系（Cummings 2018, 2020; Kawar 2015, 2021）。有关支持结构与法律机会结构的研究进一步指出，法律行动取决于组织能力以及国内程序所提供的入口（Alagna 2025; Epp 1998, 2011; Southworth 2000; Van der Pas 2023; Vanhala 2012, 2018）。这些文献共同解释了，为何边境抗争如此困难，为何分散在各处的能力如此重要，以及国内法律环境为何塑造最终可以追求的目标。就本文而言，它们有助于解释：在不同程序环境中，类似伤害的发生，为何可能需要不同组织形式或安排下的中介工作。\par
\switchcolumn
{\scriptsize\color{gray} [p. 6]}\ A related body of work studying why the movement toward the court is difficult, even when legal actors are present, helps to understand this intermediation function within a specific setting. Scholarship on transnational legal mobilization and movement lawyering shows that legal contestation depends on cross-site relations among legal, movement, and knowledge-producing actors (Cummings 2018, 2020; Kawar 2015, 2021). Research on support structures and legal opportunity structures further shows that legal action depends on organizational capacity and domestic procedural openings (Alagna 2025; Epp 1998, 2011; Southworth 2000; Van der Pas 2023; Vanhala 2012, 2018). Taken together, these literatures explain why contestation at the border is difficult, why distributed capacities are relevant, and why domestic legal settings shape what can ultimately be pursued. For our purposes, then, they help explain why a similar harm occurring may require different organizational forms or arrangements of intermediary work in different procedural settings.\par
\switchcolumn*
{\scriptsize\color{gray} [p. 6]}\ 在此背景下，转向法律不应被理解为对法院的天真信仰。有关支持结构、法律动员与社会运动律师工作的研究反而指出，当组织能力等因素使法律行动成为诸多可用行动方式之一，而且行动者追求程序与法律学说变革等超出即时救济的效果时，法律行动才变得可以设想并能够持续（Cummings 2020; Epp 1998; Kawar 2015; Southworth 2000; Vanhala 2012, 2018）。反过来，在疲惫、风险、不稳定的流动与等待负担突出的领域，不能假定人们会持续参与诉讼，因为他们可能优先考虑继续迁移、眼前安全或其他迫切需求，而非法律追索（Riva and Hoffstaedter 2021）。\par
\switchcolumn
{\scriptsize\color{gray} [p. 6]}\ Against this backdrop, a turn to law should not be read as naïve faith in courts. Research on support structures, legal mobilization, and movement lawyering suggests instead that legal action becomes thinkable and sustainable when organizational capacity, among other factors, makes it a usable repertoire among others, and when actors pursue effects that exceed immediate remedy, such as procedural and doctrinal change (Cummings 2020; Epp 1998; Kawar 2015; Southworth 2000; Vanhala 2012, 2018). Conversely, in a field marked by exhaustion, risk, precarious mobility, and the burdens of waiting, continued participation in litigation cannot be assumed, since people may instead prioritize onward movement, immediate safety, or other urgent needs over legal pursuit (Riva and Hoffstaedter 2021).\par
\switchcolumn*
{\scriptsize\color{gray} [p. 7]}\ 最后，有关巴尔干路线的研究，将推回与诉讼置于不断演变的欧洲边境制度及更广阔的抗争场域中（Hameršak et al. 2020; El-Shaarawi and Razsa 2019; Cantat 2020; Jovanović 2020; Mitrović and Vilenica 2019）。与之密切相关的研究考察基层记录工作，强调为对抗不透明与抹除，必须努力将边境知识转化为法律上可使用的事实（Alpes and Baranowska 2025; Baranowska et al. 2025; Davies et al. 2023; Fortarezza 2026; Hameršak 2021）。然而，这些文献虽然详细描述了推回如何被记录、变得可见，却很少解释，最终把这些材料转变为法律案件需要何种实际中介工作。\par
\switchcolumn
{\scriptsize\color{gray} [p. 7]}\ Finally, scholarship on the Balkan route situates pushbacks and litigation within the evolving European border regime and broader fields of contestation (Hameršak et al. 2020; El-Shaarawi and Razsa 2019; Cantat 2020; Jovanović 2020; Mitrović and Vilenica 2019). Closely related work examines grassroots documentation, highlighting the struggle to convert border knowledge into legally usable facts to combat opacity and erasure (Alpes and Baranowska 2025; Baranowska et al. 2025; Davies et al. 2023; Fortarezza 2026; Hameršak 2021). However, while this literature details how pushbacks are documented and made visible, it rarely explains the practical intermediation work required to ultimately turn them into legal cases.\par
\switchcolumn*
{\scriptsize\color{gray} [p. 7]}\ 因此，总体而言，这些文献从许多方面阐明了拼图的重要部分，提供了本研究作出贡献所需的基本类别和词汇。我们的贡献是辨识并解释：在转译开始之前，中介者开展何种初始工作，使处于错位状态的推回案件成为可处理的案件。为作此分析，本文将“案件组装”概念化为识别和连接案件初始要素，并使其在时间上保持足够长的稳定，以便向司法处理推进的工作。在此基础上，我们可以探索案件组装中更加细致的过程，也能够考察这些工作采取的协作组织形式。如前所述，当要素分散在不同行动者、地点和司法管辖区之间时，此类安排可能成为必要条件。\par
\switchcolumn
{\scriptsize\color{gray} [p. 7]}\ Overall, then, this literature in many respects illuminates important pieces of the puzzle. It provides us with the basic categories and vocabulary required for the contribution of our study: to isolate and explain the initial intermediary work, before translation sets in, through which pushback cases, characterized by misalignment, become processable. For this analysis, the article conceptualizes the term “case assembly” as the work of identifying and connecting initial elements of cases, as well as stabilizing them in time and long enough, so that they can move toward judicial processing. On this basis, we can explore a set of more granular processes of case assembly. Moreover, we can engage with the collaborative organizational form of their engagement. As noted, when elements are dispersed across actors, places, and jurisdictions, such arrangements can become a necessity.\par
\switchcolumn*
\end{paracol}
\Needspace{10\baselineskip}\section*{案例研究 / Case studies}
\begin{paracol}{2}
{\scriptsize\color{gray} [p. 7]}\ 为理解中介者如何组装案件，本文在更广泛的巴尔干路线移民管控场域内，比较意大利与奥地利的诉讼。在两个司法管辖区，国内法院都被要求评估促成向斯洛文尼亚推回、且大多进一步向克罗地亚及波斯尼亚和黑塞哥维那推回的做法。在两种情境中，寻求国际保护的人都遭受向外移送、暴力及被拒绝进入庇护程序，有时还反复遭受这些待遇。因此，比较保持了更大范围内充满阻碍的跨国环境这一条件不变。此外，在两个案例中，法院最终都谴责了相关做法，从而有助于终止向斯洛文尼亚的系统性推回（Association for Legal Studies on Immigration 2025; Fortarezza 2026; Gahleitner-gertz 2025; Moraru and Reyhani 2026）。不过，这并未终结更广泛的围堵与外部化动态（Bergesio and Bialasiewicz 2023; Fortarezza 2026; Zoppi and Puleri 2022）。而且，如下文所述，比较中的变量包括国家采用的具体做法，以及国内法律体系与程序。可获得的救济、程序路径及法律代理条件的差异，决定案件在哪些环节难以前进。因此，这一比较使本研究能够分析不同法律机会与约束对诉讼成败的影响。\par
\switchcolumn
{\scriptsize\color{gray} [p. 7]}\ To understand how intermediaries assemble cases, this article compares litigation in Italy and Austria within the broader Balkan-route field of migration control. In both jurisdictions, domestic courts were asked to assess practices that contributed to pushbacks to Slovenia, and mostly also further to Croatia and Bosnia and Herzegovina. In both, people seeking international protection were exposed to, sometimes repeatedly, onward removals, violence, and denial of access to asylum. The comparison thus holds constant the wider hostile transnational setting. Moreover, in both cases, courts came to condemn the practice, thereby contributing to an end of systematic pushbacks toward Slovenia (Association for Legal Studies on Immigration 2025; Fortarezza 2026; Gahleitner-gertz 2025; Moraru and Reyhani 2026). However, they did not end the broader dynamics of containment and externalization (Bergesio and Bialasiewicz 2023; Fortarezza 2026; Zoppi and Puleri 2022). Moreover, as elaborated below, the comparison varies the specific practices states deployed, as well as the domestic legal systems and procedures. Differences in the available remedies, procedural routes, and conditions of legal representation shaped where cases became difficult to advance. The comparison thus allows our study to analyze the influence of distinct legal opportunities and constraints on the success or failure of litigation.\par
\switchcolumn*
{\scriptsize\color{gray} [p. 7]}\ 第一项主要差异，是实施推回所采用的法律形式。在意大利，推回依据与斯洛文尼亚1996年的双边协议，以非正式重新接收的方式实施；这种做法绕过国内法与欧盟法规定的正式庇护程序，2020年至少影响1,240人（Facchini 2020）。内政部在书面声明中宣称，此类重新接收也包括寻求庇护者（Italian Ministry of the Interior 2020）；不过，后来它在面向媒体及议会的口头陈述中改变说法，称这些做法只适用于未在边境提出庇护申请的人，而这也是各国普遍采用的辩解（Italian Ministry of the Interior 2021）。通过与斯洛文尼亚的双边协议实施推回，体现了边境管控中正式与非正式做法的混合（Fortarezza 2026）。国家以重新接收协议这一正式政策工具，为推回做法提供正当化理由，但其执行却发生在法律灰区与不透明地带，例如，国家不提供任何证明此次重新接收的书面决定。\par
\switchcolumn
{\scriptsize\color{gray} [p. 7]}\ The first major point of variation concerns the legal form through which pushbacks were carried out. In Italy, pushbacks were executed as informal readmissions under a 1996 bilateral agreement with Slovenia, a practice that bypassed formal asylum procedures as regulated under national and EU law and affected at least 1240 people in 2020 (Facchini 2020). In a written statement, the Ministry declared that these readmissions also included asylum seekers (Italian Ministry of the Interior 2020), although it later changed its version in oral statements to the media and in front of the parliament, stating that they were applied only to those who had not claimed asylum at the border, a justification largely adopted by states (Italian Ministry of the Interior 2021). Pushbacks through the bilateral agreement with Slovenia are an example of the mix of formal and informal practices for border control (Fortarezza 2026). The state justified the pushback practice through a formal policy tool, the readmission agreement, whose implementation occurs within gray zones of law and opacity, for instance, by not providing any written decision by the state attesting to this readmission.\par
\switchcolumn*
{\scriptsize\color{gray} [p. 8]}\ 值得注意的是，罗马法院于2021年、又于2023年作出不利于政府的裁决（Tribunale di Roma 2021a; 2023），质疑与斯洛文尼亚的重新接收协议及其后果的合法性，包括向克罗地亚、波斯尼亚和黑塞哥维那的连锁推回，认为这些行为违反国内法、欧盟法及国际人权义务。第三项相关裁决源于针对罗马法院2021年决定的上诉（Tribunale di Roma 2021a）。在该案中（Tribunale di Roma 2021b），政府挑战最初决定，主张重新接收合法，且没有证据证明申请人本人曾进入意大利并被非正式重新接收至斯洛文尼亚。合议庭未反驳此前有关依双边协议实施重新接收属违法行为的认定，却回避讨论连锁推回和不推回原则等更广泛的问题。它转而以申请人未提供个体化证据为由，推翻了原先对其有利的决定。这一裁决表明，此类案件中提供个人证据的困难，以及涉及推回做法的法律程序的复杂性；后文将进一步讨论这些问题。\par
\switchcolumn
{\scriptsize\color{gray} [p. 8]}\ Notably, the Tribunal of Rome ruled against the government in 2021 and again in 2023 (Tribunale di Roma 2021a; 2023), questioning the legality of the readmission agreement with Slovenia and its consequences, including chain pushbacks to Croatia and Bosnia and Herzegovina, contrary to national and EU law and to international human rights obligations. The third relevant ruling arose from an appeal against the Tribunal of Rome’s 2021 decision (Tribunale di Roma 2021a). In this case (Tribunale di Roma 2021b), the government contested the initial decision, arguing the legitimacy of the readmissions and the lack of evidence that the appellant had personally entered Italy and was informally readmitted to Slovenia. The panel of judges did not refute the earlier findings on the illegality of readmissions under the bilateral agreement, but avoided discussing broader issues of chain pushbacks and non-refoulement. Instead, it overturned the favorable decision, stating that the appellant provided no individualized evidence. This decision demonstrated the challenges of presenting individual evidence in such cases and the complexities of legal proceedings involving pushback practices, which will be further explored in the following sections.\par
\switchcolumn*
{\scriptsize\color{gray} [p. 8]}\ 相比之下，在奥地利，当局将相关措施表述为依据《申根边境法典》及《奥地利外国人警察法》合法拒绝入境。这呈现了正式性与非正式性的另一种组合：国家采用正式机制，也留下文书记录；然而，边境上的非正式做法却导致庇护被系统性拒绝。更具体地说，当局声称，当事人未在边境寻求庇护，因此可以按普通边境执法处理。不过，奥地利法律并未为推回寻求庇护者提供法律依据。相反，法律要求受理并处理在境内或边境提出的庇护申请。在施泰尔马克州行政法院（LVwG Steiermark）的相关判决（2021; 2022）中，奥地利法院认定，当事人已经表达的保护请求被忽略，或未得到适当核实；当局尚未澄清他们为何进入奥地利、是否寻求庇护，就将其送返。相关认定后来在行政法院（VwGH）2022年5月的判决中获维持。州行政法院还认为，案卷显示一种更广泛模式，即此类“推回”被系统性使用；而行政法院维持了违法认定，但未将这一更广泛的定性视为裁决案件所必需的依据（LVwG Steiermark 2021, 2022, VwGH 2022a, 2022b; Gahleitner-gertz 2025）。这一差异与欧盟内部日益增长的趋势相关：庇护法日益程序化、统一化，但各国仍通过主要为非正式的单边及双边做法，规避法律，限制庇护准入（Moraru and Reyhani 2026）。\par
\switchcolumn
{\scriptsize\color{gray} [p. 8]}\ In Austria, by contrast, the relevant measures were framed by the authorities as lawful refusals of entry under the Schengen Borders Code and the Austrian Aliens Police Act. It presents a different mix of formality and informality. The state uses a formal mechanism, also leaving a paper trail behind. However, informal practices at the border led to the systematic denial of asylum. More specifically, the authorities claimed that the people concerned had not sought asylum at the border and could therefore be treated as subject to ordinary border enforcement. However, Austrian law contains no legal basis for the pushback of asylum seekers. On the contrary, it requires the acceptance and processing of asylum applications made on the territory or at the border. In decisions such as LVwG Steiermark (2021; 2022), later upheld in the related VwGH decisions of May 2022, Austrian courts found that articulated requests for protection had been ignored or not properly verified, and that the persons concerned were returned without the authorities first clarifying why they had entered Austria or whether they sought asylum. The LVwG further described the file as suggesting a broader pattern in which such “pushbacks” were systematic in their use, while the VwGH upheld the unlawfulness findings without treating that broader characterization as necessary to decide the case (LVwG Steiermark 2021, 2022, VwGH 2022a, 2022b; Gahleitner-gertz 2025). This variation relates to a growing trend within the European Union, where asylum law is increasingly proceduralized and harmonized. Yet, countries still circumvent the law to restrict access to asylum by mostly informal uni- and bilateral approaches (Moraru and Reyhani 2026).\par
\switchcolumn*
{\scriptsize\color{gray} [p. 9]}\ 第二项主要差异涉及审理机关、救济与程序。在意大利，诉讼主要在民事法院展开。诉请人可以依据《意大利民事诉讼法典》第700条启动紧急简易程序，也可以提出损害赔偿请求。相应地，可获得的救济范围更广，包括撤销重新接收决定、发放人道主义签证，以及对基本权利遭受侵犯给予赔偿。在意大利第一起案件中，诉请人在连锁推回后仍滞留波斯尼亚和黑塞哥维那，紧急程序尤为重要，因为它提供了可能的入境途径。在第二起案件中，罗马法院2023年的判决，对非法重新接收造成的非财产性损害判给赔偿。相比之下，奥地利的相关案件以针对行政执行措施的申诉，提交州行政法院审理。救济较为有限，主要包括宣告违法和金钱赔偿，并未赋予诉请人合法返回的权利。申诉还必须在短短六周的期限内提出。因此，核心法律问题是：当事人是否被剥夺提出庇护申请的机会，并因此被非法拒绝入境。\par
\switchcolumn
{\scriptsize\color{gray} [p. 9]}\ A second major point of variation concerns forum, remedy, and procedure. In Italy, litigation unfolded primarily before civil courts. Claimants could proceed through urgent summary proceedings under Article 700 of the Italian Code of Civil Procedure or through damages claims. The available remedies were correspondingly broader and could include annulment of readmissions, humanitarian visas, and compensation for violations of fundamental rights. In the first Italian case, where the claimant remained stranded in Bosnia and Herzegovina after a chain pushback, urgent proceedings were especially significant because they offered a possible route to entry. In the second, the Tribunale di Roma judgment of 2023 awarded compensation for non-pecuniary harm caused by the unlawful readmission. In Austria, by contrast, the relevant cases were brought before provincial administrative courts through complaints against executive measures. The remedies were more limited, consisting chiefly of declarations of illegality and monetary compensation, and claimants were not granted a right to return lawfully. The complaint had to be lodged within a short six-week deadline. The central legal question was therefore whether people had been denied the chance to lodge asylum claims and had accordingly been unlawfully refused entry.\par
\switchcolumn*
{\scriptsize\color{gray} [p. 9]}\ 第三项差异涉及诉请人的处境。在意大利，两起关键案件的诉请人处境明显不同：一位是处于非正规身份的巴基斯坦公民，在从意大利向斯洛文尼亚、再向克罗地亚的连锁推回后，滞留波斯尼亚和黑塞哥维那；另一位也是巴基斯坦公民，但在此前遭遇推回尝试之后，已经在意大利取得难民身份。相比之下，在奥地利，施泰尔马克州行政法院两案（2021, 2022）的当事人，是在奥地利边境被拒绝进入庇护程序、随后被移送至斯洛文尼亚的非正规身份者，其中包括一名未成年人。\par
\switchcolumn
{\scriptsize\color{gray} [p. 9]}\ A third point of variation concerns claimant position. In Italy, the two key cases involved claimants in markedly different situations: one was a Pakistani citizen in irregular status, stranded in Bosnia and Herzegovina after a chain pushback from Italy to Slovenia and onward to Croatia, while the other was a Pakistani citizen who had already obtained refugee status in Italy after a prior pushback attempt. In Austria, by contrast, the litigants in LVwG Steiermark (2021, 2022) were people in irregular status who had been removed to Slovenia after being denied access to asylum at the Austrian border, including a minor.\par
\switchcolumn*
{\scriptsize\color{gray} [p. 9]}\ 最后，法律代理规则也有差别：奥地利的推回案件缺乏免费法律援助，依赖公益法律代理（Gahleitner-gertz 2025）；意大利虽有免费法律援助，但诉请人位于司法管辖区之外时，严格的程序规则使法律代理关系难以正式确立（Massimi 2026）。\par
\switchcolumn
{\scriptsize\color{gray} [p. 9]}\ Finally, rules for legal representation vary too: Austria, for pushback cases, lacks free legal aid and relies on pro bono legal representation (Gahleitner-gertz 2025), and Italy has free legal aid available but encounters difficulties in formalizing legal representation when the claimant is outside the jurisdiction due to strict procedural rules (Massimi 2026).\par
\switchcolumn*
\end{paracol}
\Needspace{10\baselineskip}\section*{研究方法 / Methods}
\begin{paracol}{2}
{\scriptsize\color{gray} [p. 9]}\ 本研究结合案头研究与半结构式访谈，访谈对象是参与在意大利及奥地利法院对巴尔干路线沿线推回提出司法挑战的关键法律中介者。有关资料的可获取性，参见“数据可用性声明”。\par
\switchcolumn
{\scriptsize\color{gray} [p. 9]}\ The research combines desk research with semi-structured interviews with key legal intermediaries involved in judicial contestation of pushbacks along the Balkan route before Italian and Austrian courts. For information on the availability of the sources, please see the “Data availability statement.”\par
\switchcolumn*
{\scriptsize\color{gray} [p. 9]}\ 首先，我们对涉及意大利、奥地利向斯洛文尼亚推回的七项判决、政府官方声明，以及媒体、网站和报告中的补充资料，进行了定性分析。这有助于理解参与其中的法律中介者的总体构成、两种情境中他们的异同，并据此选择要访谈的关键相关人士。此外，这也使我们能够突出可能与案例相关的组织及制度因素，交代背景，从而理解中介者为进入国内法院所采取的策略与行动。\par
\switchcolumn
{\scriptsize\color{gray} [p. 9]}\ First, we qualitatively analyzed seven judgments on pushbacks from Italy and Austria to Slovenia, governments’ official statements, and additional sources from the media, websites, and reports. This allowed for understanding the overall constellation of legal intermediaries involved, the similarities and differences between them in the two contexts, and the selection of key stakeholders to be interviewed. Moreover, it enabled us to set the context by highlighting the organizational and institutional factors that may be relevant to our case studies, thereby facilitating an understanding of the strategies and actions intermediaries take to access national courts.\par
\switchcolumn*
{\scriptsize\color{gray} [p. 9]}\ 第二阶段，是对与已诉案件最直接相关的行动者进行半结构式专家访谈（见附录）。意大利的受访者包括两名活跃于庇护与移民诉讼的律师、一名具有跨境实地经验且曾参与社会工作的行动者，以及一名曾帮助识别和记录案件的记者。奥地利的访谈覆盖两起主要法院案件中的关键行动者：两名建立热线、为在奥地利边境遭遇推回者提供首次联系的行动者，一名充当联系窗口及组织者的律师兼非政府组织协调员，以及在相关案件中代理诉请人的律师。鉴于诉讼的跨国结构，我们还访谈了一名斯洛文尼亚非政府组织律师，其工作与意大利、奥地利两国案卷均有关联。在奥地利，我们进一步访谈了对两案作出一审裁决的法官，以理解裁判如何回应中介者的努力。我们未能访谈意大利法官，但其判决提供了了解司法采纳情况的替代资料。我们原本也计划访谈诉请人，但由于相关裁决极少，由此存在匿名性风险，最终未实施。总体而言，样本旨在捕捉面向诉讼的中介者如何组装案件并将其带向法院，而非重建所有受到推回影响者的经历。访谈采用半结构形式，既确保该领域专家能够开放作答，也收集可比较的数据（Ahlin 2019）。问题涉及诉诸法院的策略，包括涉及向斯洛文尼亚推回的案件中，从监测、构建案件到进入法庭各阶段的诉讼选择。此外，我们还询问这些策略背后的原因，以及与个人、组织和制度动态有关的挑战与优势。访谈也关注受访者与其他参与行动者之间的互动，以及这些互动对诉讼和策略选择的影响。\par
\switchcolumn
{\scriptsize\color{gray} [p. 9]}\ The second stage consisted of semi-structured expert interviews with the actors most directly involved in the litigated cases (Appendix). In Italy, this included two lawyers active in asylum and migration litigation, one activist with cross-border field experience and prior social-work involvement, and one journalist who had helped identify and document cases. In Austria, interviews covered the key actors involved in the two main court cases: two activists who built a hotline for first contact in case of pushbacks at the Austrian border, a lawyer and NGO coordinator who functioned as a contact point and organizer, and the attorney representing claimants in the cases. Given the transnational structure of the litigation, we also interviewed a Slovenian NGO lawyer whose work was relevant to both the Italian and Austrian files. In Austria, we further interviewed the judge who decided both cases in the first instance, in order to understand how intermediary efforts were received in adjudication. We were unable to interview the Italian judges, but their judgments provided a surrogate source on judicial uptake. We initially planned to interview claimants as well, but given the very few relevant court rulings and the resulting anonymity risks, we did not pursue this. Overall, moreover, the sample is designed to capture how litigation-facing intermediaries assembled and carried cases toward court, rather than to reconstruct the experience of all actors affected by pushbacks. The interviews followed a semi-structured format to guarantee open-ended answers from experts in the field, while also collecting comparable data (Ahlin 2019). Interview questions addressed strategies for accessing courts, including choices related to litigation in cases involving pushbacks to Slovenia at various stages, from monitoring and constructing the case to the courtroom. These were supplemented by inquiries into the reasons behind these strategies, as well as the challenges and strengths related to individual, organizational, and institutional dynamics. Moreover, interviews focused on interactions with other actors involved and their influence on litigation and on strategic choices.\par
\switchcolumn*
{\scriptsize\color{gray} [p. 10]}\ 访谈已获博洛尼亚大学研究伦理委员会批准；所有参与者均签署了知情与隐私同意书，访谈材料已匿名化。每次访谈时长一小时，使用意大利语、德语或英语进行。访谈均被录音，并以原语言逐字转录，随后译为英语，采用溯因策略进行主题编码分析（Timmermans and Tavory 2012）。分析文献及半结构式访谈数据时，我们采用从现有文献中提取的主题，并将其与直接从访谈和文本中浮现的主题相结合。\par
\switchcolumn
{\scriptsize\color{gray} [p. 10]}\ Research ethics approval from the University of Bologna’s research ethics board has been received to conduct interviews, and all participants signed information and privacy consent forms, and the interview material was anonymized. Interviews lasted one hour and were conducted in Italian, German, and English. They were recorded and transcribed verbatim in the original language. Then, they were translated into English and analyzed through thematic coding using an abductive strategy (Timmermans and Tavory 2012). Data from documents and semi-structured interviews were analyzed using themes emerging from the existing literature, which were integrated with themes directly emerging from the interviews and texts.\par
\switchcolumn*
{\scriptsize\color{gray} [p. 10]}\ 最后，有必要对访谈样本作几点反思。第一，由于未访谈诉请人，本文中的诉请人能动性只能通过面向诉讼的中介者的讲述间接呈现。因此，当受访者描述诉请人决定推进、退出、等待、继续迁移，或仅有条件地保持参与时，本研究能够识别这些行动，却无法恢复诉请人以自己的声音表达的观点。因此，相比诉请人的选择、拒绝、策略性迁移或脱离，样本更可能突出中介者的劳动。第二，我们的访谈仅捕捉反推回抗争中的较小部分：这些行动者因其角色、组织位置及跨境网络而接近罕见的法律机会，并在实践中承担决定潜在案件能否转化为法院案卷的责任。对于已抵达或接近抵达法院的案卷，本研究的重建最为充分；对于更早退出、在访谈材料中留下较少痕迹的广泛潜在案件，重建则较为薄弱。\par
\switchcolumn
{\scriptsize\color{gray} [p. 10]}\ Finally, some reflections on the interview sample are needed. First, because claimants were not interviewed, the claimant agency appears here only indirectly, through litigation-facing intermediary accounts. The reconstruction can therefore identify claimant action where interviewees describe it, for example, in decisions to proceed, withdraw, wait, move on, or remain only conditionally engaged, but it cannot recover claimant perspectives in their own voice. The sample is thus likely to foreground intermediary labor more strongly than claimant choice, refusal, tactical movement, or disengagement. Second, our interviews capture only a smaller segment of anti-pushback contestation: actors whose roles, organizational positions, and cross-border networks placed them close to rare legal openings and made them responsible, in practice, for deciding whether a possible case could be turned into a court file. The reconstruction is strongest for files that reached, or came close to reaching, court. It is thinner for the wider universe of possible cases that dropped out earlier and left fewer traces in the interview material.\par
\switchcolumn*
\end{paracol}
\Needspace{10\baselineskip}\section*{经验材料 / Empirical data}
\begin{paracol}{2}
{\scriptsize\color{gray} [p. 11]}\ 本节追踪在推回案件进一步走向法院之前，使其成为可处理案件的工作。根据访谈数据，内容围绕中介者组装工作的三个基本领域组织：痕迹、关系及诉讼路径。此外，本节还描述裁判如何检验已组装的案件，以及之后发生的事情。\par
\switchcolumn
{\scriptsize\color{gray} [p. 11]}\ This section follows the work that made pushback cases processable before they moved further toward the courts. It is organized around the three basic areas of the assembly work of intermediaries, as suggested by interview data: on traces, relationships, and litigation paths. In addition, the section describes how adjudication tested the assembled cases and what came thereafter.\par
\switchcolumn*
{\scriptsize\color{gray} [p. 11]}\ 不过，我们首先依据现有法院判决，重建基础案件的主要事实，以交代法律中介者介入之前的背景。意大利两案涉及两名申请人，他们在沿巴尔干路线长期、反复迁移之后，抵达的里雅斯特。在后来由罗马法院作出裁决的第一案中（2021a），一名当事人在克罗地亚边境遭受暴力后，于2020年7月中旬抵达。在第二案中，当事人于2019年已到达波斯尼亚和黑塞哥维那，遭受过九次克罗地亚推回和三次斯洛文尼亚推回，又经过一次始于数周前的尝试，于2020年10月中旬在的里雅斯特首次进入意大利（Tribunale di Roma 2023）。两案中，当事人几乎抵达后立即接触意大利当局，被带到警察设施内，接受拍照或医疗筛查，被要求提供信息，并签署既未翻译、也未解释内容的文件。第一案中，手机被没收，整组人被戴上手铐（Tribunale di Roma 2021a, 2023）。他们表达了在意大利寻求庇护的愿望，却仍被非正式重新接收至斯洛文尼亚，随后继续被推回克罗地亚及波斯尼亚和黑塞哥维那；后续移送伴随着暴力（Tribunale di Roma 2021a, 2023; 访谈1; 访谈3; 访谈9）。在波斯尼亚和黑塞哥维那，第一位申请人因利帕营地没有空位而未获得接待，只能在萨拉热窝一栋遭战争破坏的废弃建筑中栖身；第二位申请人后来则在比哈奇附近的非正式聚居点被找到（Tribunale di Roma 2021a, 2023）。\par
\switchcolumn
{\scriptsize\color{gray} [p. 11]}\ First, however, we provide the context preceding the engagement of legal intermediaries by reconstructing the main facts of the underlying cases from available court judgments. The Italian cases concern two applicants who reached Trieste after lengthy, repeated movements along the Balkan route. In the first case, later decided in Tribunale di Roma (2021a), an individual arrived in mid-July 2020 after violence at the Croatian border. In the second, the individual had already reached Bosnia and Herzegovina in 2019, had suffered nine Croatian and three Slovenian pushbacks, and made his first entry into Italy at Trieste in mid-October 2020 after a further attempt begun a few weeks before (Tribunale di Roma 2023). In both files, contact with Italian authorities followed almost immediately after arrival and the individuals were taken to police facilities, photographed or medically screened, asked for information, and made to sign papers whose content was not translated or explained. In the first case, phones were confiscated, and the group was handcuffed (Tribunale di Roma 2021a, 2023). They expressed a wish to seek asylum in Italy, but were nevertheless informally readmitted to Slovenia and then pushed onward to Croatia and Bosnia and Herzegovina, where violence accompanied the onward removal (Tribunale di Roma 2021a, 2023; Interview 1; Interview 3; Interview 9). In Bosnia and Herzegovina, the first applicant was left without reception after Lipa camp had no space and found shelter in an abandoned, war-damaged building in Sarajevo, while the second was later located in the informal settlement near Bihać (Tribunale di Roma 2021a, 2023).\par
\switchcolumn*
{\scriptsize\color{gray} [p. 11]}\ 奥地利2020年的案件中，一群人从斯洛文尼亚越过未设正式口岸的边界；他们担心若在边境附近被拦截会立即遭推回，因而试图向内陆移动，最后在锡歇尔多夫附近被发现，当时浑身湿透、肮脏，瑟瑟发抖（LVwG Steiermark 2021; VwGH 2022a; 访谈8）。在运输与拘留过程中，诉请人及其他人反复说出“庇护”一词，但在锡歇尔多夫，他们被完全脱光衣服、采集指纹，虽提出请求却未获食物，也未被询问到奥地利的目的。当天下午，他们便被移交斯洛文尼亚警察，随后连锁推回继续经过克罗地亚，直至波斯尼亚和黑塞哥维那（LVwG Steiermark 2021; VwGH 2022a; 访谈7，奥地利，法官; 访谈8; 访谈9）。涉及一名未成年人的2021年案件中，一群人再次从斯洛文尼亚步行越境，并主动寻求联系奥地利警察，以请求保护（LVwG Steiermark 2022; VwGH 2022b）。但这名未成年人仍被带往锡歇尔多夫，接受新冠检测与身份核验；当局反复告诉他，程序稍后才会开始，四至五小时后，却将他交回斯洛文尼亚警员（LVwG Steiermark 2022; VwGH 2022b）。不过，与2020年的诉请人不同，他未被继续推回波斯尼亚和黑塞哥维那，因为斯洛文尼亚启动了庇护程序，他后来在那里获得庇护（LVwG Steiermark 2022; VwGH 2022b; 访谈7，奥地利，法官）。\par
\switchcolumn
{\scriptsize\color{gray} [p. 11]}\ In Austria, in the 2020 case, a group crossed the green border from Slovenia, tried to move inland because they feared immediate pushback if stopped near the border, and was eventually found near Sicheldorf, wet, dirty, and shivering (LVwG Steiermark 2021; VwGH 2022a; Interview 8). During transport and custody, the claimant and others repeatedly used the word asylum, but at Sicheldorf they were fully stripped, fingerprinted, denied food despite requests, not asked what they wanted in Austria, and then handed to the Slovenian police the same afternoon, after which the chain continued through Croatia to Bosnia and Herzegovina (LVwG Steiermark 2021; VwGH 2022a; Interview 7, Austria, judge; Interview 8; Interview 9). In the 2021 case, relating to a minor, the group again crossed on foot from Slovenia and sought contact with the Austrian police in order to ask for protection (LVwG Steiermark 2022; VwGH 2022b). He was nevertheless taken to Sicheldorf, subjected to Covid testing and identification, repeatedly told that the procedure would begin later, and after four to five hours handed back to Slovenian officers (LVwG Steiermark 2022; VwGH 2022b). Yet, unlike the 2020 claimant, he was not pushed onward to Bosnia and Herzegovina, because Slovenia opened an asylum procedure and he later obtained asylum there (LVwG Steiermark 2022; VwGH 2022b; Interview 7, Austria, judge).\par
\switchcolumn*
\end{paracol}
\Needspace{10\baselineskip}\section*{关于痕迹 / On traces}
\begin{paracol}{2}
{\scriptsize\color{gray} [p. 12]}\ 鉴于这种被本文称为充满阻碍的跨国环境的现实，推回最初以碎片形式出现：在波斯尼亚和黑塞哥维那收集的证词、边境附近的观察、紧急来电、行政记录中的矛盾，以及在一个司法管辖区产生、后来却在另一个管辖区变得重要的文件（访谈9; 访谈1; 访谈5）。此时，这些材料形成了一个潜在案件库，也逐渐让人意识到，这些并非孤立事件。\par
\switchcolumn
{\scriptsize\color{gray} [p. 12]}\ Given this reality, described here as a hostile transnational setting, pushbacks first appeared in fragments: testimonies gathered in Bosnia and Herzegovina, observations near the border, emergency calls, administrative inconsistencies, and documents produced in one jurisdiction that would later matter in another (Interview 9; Interview 1; Interview 5). At that point, the material produced a pool of possible cases and, gradually, the sense that these were not isolated events.\par
\switchcolumn*
{\scriptsize\color{gray} [p. 12]}\ 在斯洛文尼亚，最初警讯来自一名在波斯尼亚工作的斯洛文尼亚记者。随后在大克拉杜沙及比哈奇开展的访谈，收集到一些人的证词与文件；他们表示，自己被从斯洛文尼亚推回克罗地亚，又推回波斯尼亚和黑塞哥维那（访谈9）。这项工作促成了连接斯洛文尼亚、克罗地亚和意大利行动者的网络，目的是“相互通报侵权情况，并尽可能多地找到之后可以交由法院处理的案件”（访谈9）。\par
\switchcolumn
{\scriptsize\color{gray} [p. 12]}\ In Slovenia, the first alert came from a Slovenian journalist working in Bosnia. Follow-up interviews in Velika Kladuša and Bihać then produced testimonies and documents from people who said they had been pushed back from Slovenia to Croatia and then to Bosnia and Herzegovina (Interview 9). A network linking actors in Slovenia, Croatia, and Italy grew out of that work, “to notify each other about the violations and in order to find as much as we can cases that can then be processed before the courts” (Interview 9).\par
\switchcolumn*
{\scriptsize\color{gray} [p. 12]}\ 在意大利，的里雅斯特一名律师回忆，有关巴尔干路线沿线暴力与虐待的报告，使人相信可能存在从意大利领土经斯洛文尼亚、克罗地亚向波斯尼亚和黑塞哥维那的连锁推回；后来前往波斯尼亚和黑塞哥维那的实地访问，以及与当地非政府组织的联系，开始填补这幅图景（Amnesty International 2019; 访谈1）。一名驻波斯尼亚记者这样描述早期的一次转介：“有些人告诉我，他们被意大利非法拒绝入境，我就让他们与\textit{[法律协会]}取得联系。其中一人走完了全部流程［……］以核实［……］，然后，一切就从那里开始运转起来了”（访谈4）。这名记者只是通过媒体资料，间接得知该法律协会针对沿线推回开展诉讼（访谈4）。同次访谈还提到，一些出于其他目的收集的材料，后来随案卷流转：“所有照片都有日期，如果需要，元数据里还有地点”，以及“这个小型非正式营地所在地的GPS坐标”（访谈4）。在的里雅斯特一侧，从事接待工作的行动者开始注意到，边境警察的说法与实际抵达接待点的人群并不相符。随后，律师通过紧急会议、申请查阅档案、议会质询，以及在行动者和波斯尼亚和黑塞哥维那联系人帮助下进一步收集证词，将这一问题向前推进（访谈1; Association for Legal Studies on Immigration 2023）。\par
\switchcolumn
{\scriptsize\color{gray} [p. 12]}\ In Italy, a lawyer in Trieste recalls that reports on violence and abuse along the Balkan route made chain pushbacks through Slovenia, Croatia, and Bosnia and Herzegovina from Italian territory plausible, while later trips to Bosnia and Herzegovina and contacts with local NGOs began to fill in that picture (Amnesty International 2019; Interview 1). One Bosnia-based journalist described an early handoff like this: “I put some of these people who told me they had been illegally rejected by Italy in touch with \textit{[the legal association]}, one person went through the whole process […] to verify […] and then from there everything started moving” (Interview 4). The journalist knew of the legal association’s litigation against pushbacks along the route only indirectly through media sources (Interview 4). The same interview points to the material that, collected for other reasons, later traveled with the file: “the photos all have the date and, if desired, the location in the metadata” and “the GPS coordinates of where this small informal camp is located” (Interview 4). On the Trieste side, reception actors began noticing a mismatch between what border police said and who actually arrived at the reception point. Lawyers then took the issue into urgent meetings, access-to-files requests, a parliamentary inquiry, and further testimony-gathering with the help of activists and contacts in Bosnia and Herzegovina (Interview 1; Association for Legal Studies on Immigration 2023).\par
\switchcolumn*
{\scriptsize\color{gray} [p. 12]}\ 在奥地利，行动者描述，往往“没有姓名，没有地点，什么都没有”，尤其是在只有六周起诉期限的情况下（访谈5）。当地媒体与志愿者关于施泰尔马克州南部大规模警察行动的报道，只有与流动人群网络、克罗地亚及斯洛文尼亚合作组织提供的信息连接之后，才变得有用（访谈6）。当地非政府组织为在边境注意到情况、却不知道下一步该怎么办的人提供接收线索的窗口；一个行动者网络则建立了热线，使身处奥地利境内的人可以说明，他们想寻求庇护，却正在遭受推回（访谈5; 访谈6）。记者、当地行动者和其他组织，也转交行动者自己未曾产生的联系电话及其他痕迹（访谈5; 访谈9）。在2021年的未成年人案件中，一项痕迹保存在案卷本身：这群人利用智能手机GPS抵达巴特拉德克斯堡，并主动寻找警察局以请求庇护，尽管当局后来否认这一请求，或称未听到该请求（LVwG Steiermark 2022; VwGH 2022b）。\par
\switchcolumn
{\scriptsize\color{gray} [p. 12]}\ In Austria, activists describe that there was often “no names, no location, nothing,” especially under a six-week filing deadline (Interview 5). Reports from local media and volunteers about large police operations in southern Styria only became useful once they were connected with information from people on the move networks and partner organizations in Croatia and Slovenia (Interview 6). Local NGOs worked as an intake point for those who noticed something at the border but did not know what to do next, while a network of activists set up a hotline through which people on Austrian territory could say that they wanted asylum and were being pushed back (Interview 5; Interview 6). Journalists, local actors, and other organizations also passed on contact numbers and other traces that the activists themselves had not generated (Interview 5; Interview 9). In the 2021 minor case, one trace survived in the case file itself: the group used smartphone GPS to reach Bad Radkersburg and actively sought the police station in order to ask for asylum, though the request was later denied or said not to have been heard (LVwG Steiermark 2022; VwGH 2022b).\par
\switchcolumn*
{\scriptsize\color{gray} [p. 13]}\ 奥地利的材料还表明，这些碎片要与地点相互核验。行动者返回相关边境地区，拍摄河流、边境站及当地标志，并将这些环境与诉请人所述的事件进行比较。对施泰尔马克州南部边境的地方知识，有助于理解零碎描述，并判断它们能否在一份案卷中相互衔接（访谈5）。\par
\switchcolumn
{\scriptsize\color{gray} [p. 13]}\ The Austrian material also shows that these fragments were tested against place. Activists returned to the relevant border area, photographed the river, the border station, and local signs, and compared those surroundings with what claimants said had happened there. Local knowledge on the southern Styrian border helped make sense of fragmentary descriptions and decide whether they could hold together in a file (Interview 5).\par
\switchcolumn*
{\scriptsize\color{gray} [p. 13]}\ 后续用途已经开始塑造当下收集的内容。行动者表示，他们与律师密切合作，因为律师知道哪些细节对法庭重要。他们借助翻译人员，记录地点、先后顺序、制服及事件过程，而非停留在一般叙述层面（访谈5）。正如他们所说，“始终都在于细节”，否则，一份陈述后来可能被认为不可信而遭驳回（访谈5）。\par
\switchcolumn
{\scriptsize\color{gray} [p. 13]}\ Later use was already shaping what was being collected. Activists say they worked closely with the lawyer because he knew which details would matter in court. Place, sequence, uniforms, and the course of events were recorded with the help of translators rather than left at the level of a general story (Interview 5). As they put it, “it’s always about the details,” because otherwise a statement could later be dismissed as implausible (Interview 5).\par
\switchcolumn*
{\scriptsize\color{gray} [p. 13]}\ 在两种情境中，斯洛文尼亚的文书记录往往提供了其他司法管辖区缺少的痕迹。斯洛文尼亚受访者描述，当地的重新接收会生成记录，证明当事人确曾在斯洛文尼亚领土上接受处理；这些记录后来可用于其他地方的诉讼，包括意大利和奥地利。在这些国家，国家留下的痕迹可能不存在、更为稀少，或因其他原因较不利于问责（访谈9; 访谈5）。由于斯洛文尼亚组织也可以在庇护程序中代理当事人，获取文书与针对个人的后续跟进，有时能够在当地结合起来，而其他地方未必具备同样便利的条件（访谈9）。但即便在那里，问题仍然是：是否有“具体证据证明，他们确实曾在斯洛文尼亚领土上”，以及“文件是否与他们的陈述相符”（访谈9）。\par
\switchcolumn
{\scriptsize\color{gray} [p. 13]}\ Across both settings, Slovenian documentation often supplied a trace that other jurisdictions lacked. The Slovenian interviewee describes readmissions there as producing records showing that the person had in fact been processed on Slovenian territory, records that could later be used in litigation elsewhere, including in Italy and Austria, where state traces were absent, thinner, or otherwise less useful for accountability (Interview 9; Interview 5). Because Slovenian organizations could also represent people within asylum procedures, documentary access and person-specific follow-up could sometimes be combined there in ways that were not as readily available elsewhere (Interview 9). Even there, though, the issue remained whether there was “concrete proof that they were actually on Slovenian territory,” and whether “the documentation would coincide with their statements” (Interview 9).\par
\switchcolumn*
\end{paracol}
\Needspace{10\baselineskip}\section*{关于关系 / On relationships}
\begin{paracol}{2}
{\scriptsize\color{gray} [p. 13]}\ 下一个问题，是与某个仍能被找到的人建立、或重新建立工作关系。此人可能愿意、也可能不愿继续，而且需要保持足够长时间的联系，以便随后完成授权、作证与起诉。在奥地利，最初的一组人，是通过波斯尼亚同事转交的电话号码进入视野的。以下这段话概括了之后的工作：\par
\switchcolumn
{\scriptsize\color{gray} [p. 13]}\ The next problem was to establish, or re-establish, a working relation with a person who could still be found, might or might not be willing to proceed, and could remain in contact long enough for mandate, testimony, and filing to follow. In Austria, the first group arrived as phone numbers passed on by colleagues in Bosnia. One passage captures what followed:\par
\switchcolumn*
{\scriptsize\color{gray} [p. 13]}\ 我们得投入大量工作，才能确保这些案件甚至存在［……］你必须找到那些人［……］我们投入很多资源，真正了解受影响的人，向他们解释这里的制度如何运作［……］四个人作了陈述，但只有一个人同意担任申诉人。（访谈5）\par
\switchcolumn
{\scriptsize\color{gray} [p. 13]}\ we had to put in a lot of work to ensure that these cases even existed […] you have to find the people […] we put a lot of resources into really getting to know the people affected and explaining to them how the system works here […] four made statements, but only one person agreed to act as the complainant. (Interview 5)\par
\switchcolumn*
{\scriptsize\color{gray} [p. 13]}\ 意大利材料呈现了同样的问题，只是空间距离更远、持续时间更长。中介者在意大利，诉请人在波斯尼亚和黑塞哥维那，且常处于不稳定、不断变化的条件中，因此，维持联系依靠一名律师所说的“从的里雅斯特一路往下”的整条链（访谈1）。驻波斯尼亚的行动者及当地联系人，在营地、占住的空屋或非正式聚居点寻找诉请人，保持联系，陪同他们赴约，并在电话号码不断变化、人员消失，以及等待本身就可能使案卷无法成立的环境中工作（访谈2; 访谈1）。一名行动者将这一等待期与保障食物、住宿、洗浴设施及衣物的需要联系起来；与此同时，当法律与官僚程序在别处准备时，还要使诉请人保持参与（访谈2）。同样的工作还包括解释程序、可能带来的结果，以及短期内无法承诺的事情。在意大利的紧急程序路径中，诉请人必须在波斯尼亚停留足够久，才能使这条路径具有意义，尽管更广泛的环境始终在推动其继续迁移。\par
\switchcolumn
{\scriptsize\color{gray} [p. 13]}\ The Italian material shows the same problem stretched across a greater distance and over a longer period of time. Intermediaries were in Italy while claimants were in Bosnia and Herzegovina, often in precarious and shifting conditions, so maintaining contact depended on what one lawyer called the whole chain “from Trieste down” (Interview 1). Bosnia-based activists and local contacts located claimants in camps, squats, or informal settlements, kept them in touch, accompanied them to appointments, and worked in an environment where phone numbers changed, people disappeared, and waiting itself could undo the file (Interview 2; Interview 1). One activist ties that waiting period to the need to secure food, accommodation, washing facilities, and clothing, while also keeping the claimant engaged as legal and bureaucratic steps were prepared elsewhere (Interview 2). The same work also involved explaining the proceedings, what might follow from them, and what they could not promise in the short term. In the urgent Italian route, the claimant had to remain in Bosnia long enough for that route to make sense, even though the wider environment kept pulling toward onward movement.\par
\switchcolumn*
{\scriptsize\color{gray} [p. 14]}\ 并非所有最初找到的人，后来都能成为诉至法院的当事人。有些叙述因时间顺序与路线无法一致重建而被排除；另一些则因无法取得文件、无法保持联系，或当事人未在原地停留足够久以建立案卷而退出（访谈2; 访谈4）。不过，案件流失也反映了诉请人在严峻约束下的优先事项。斯洛文尼亚受访者表示，监测正在发生的事情，往往比找到愿意追究案件的人更容易，因为许多身在波斯尼亚的人更关心抵达目的国，而非启动法律程序；他们不断移动，也使跟进困难（访谈9）。奥地利行动者同样描述，必须通过反复交谈建立信任，某些细节——包括暴力搜身与强迫脱衣——才会得到更充分披露（访谈5）。材料中不断出现同样的实际问题：当事人是否继续参与？能否维持联系？是否有足够证据，或至少有一条可信的取证路径，证明当事人曾在意大利、奥地利或斯洛文尼亚领土上，随后遭到推回（访谈5; 访谈1; 访谈9）？潜在案件库因此大幅缩小。\par
\switchcolumn
{\scriptsize\color{gray} [p. 14]}\ Not everyone first identified could later become a claimant for court. Some accounts dropped out because chronology and route could not be reconstructed consistently. Others fell away because documents could not be obtained, contact could not be maintained, or the person did not remain in place long enough for the file to be built (Interview 2; Interview 4). Yet attrition also reflected claimant priorities under severe constraint. The Slovenian interview states it was often easier to monitor what was happening than to find a person willing to pursue a case, because many people in Bosnia were more concerned with reaching the desired state than with initiating legal proceedings, and were difficult to follow as they moved (Interview 9). Austrian activists similarly describe how trust had to be built over repeated conversations before some details, including violent searches and forced stripping, were disclosed more fully (Interview 5). Across the material, the same practical questions recurred: would the person continue, could contact be maintained, and was there enough proof, or at least a plausible route to proof, that the person had been on Italian, Austrian, or Slovenian territory and then been pushed back (Interview 5; Interview 1; Interview 9). The pool thus became considerably smaller.\par
\switchcolumn*
\end{paracol}
\Needspace{10\baselineskip}\section*{关于诉讼路径 / On litigation paths}
\begin{paracol}{2}
{\scriptsize\color{gray} [p. 14]}\ 意大利一侧，一名律师这样概括授权问题：\par
\switchcolumn
{\scriptsize\color{gray} [p. 14]}\ On the Italian side, one lawyer reduced the problem of authorization to this:\par
\switchcolumn*
{\scriptsize\color{gray} [p. 14]}\ 授权委托书无疑是最难处理的要素之一［……］这份委托书必须在公证人或公职人员面前签署，由公职人员或公证人认证，经使馆认证并翻译［……］从你找到当事人、他向当地行动者讲述经历并表示自己曾被从意大利重新接收\textit{[诉请人]}次，到你能够提出诉请，可能会经过数月，或至少数周，因为你需要时间撰写诉状，可能还要申请查阅档案、取得授权委托书，等等。（访谈3）\par
\switchcolumn
{\scriptsize\color{gray} [p. 14]}\ the power of attorney was certainly one of the most complex elements to manage […] This power of attorney must be signed before a notary or a public official, authenticated by the public official or notary, legalized and translated by the embassy […] from the moment you identify the person who shares their testimony with local activists and says they have been readmitted \textit{[the claimant]} times from Italy, to the moment you are able to file the appeal, several months or at least several weeks can pass, because you need time to write the appeal, perhaps request access to the files, collect the power of attorney, etc. (Interview 3)\par
\switchcolumn*
{\scriptsize\color{gray} [p. 14]}\ 这一漫长的官僚程序还有另一重大障碍：要求委托人持有有效身份证件，而迁徙路线上的人很少具备这一条件（访谈1）。\par
\switchcolumn
{\scriptsize\color{gray} [p. 14]}\ This long bureaucratic process has another major obstacle: the requirement for a valid client ID, which is rare among people on the migratory route (Interview 1).\par
\switchcolumn*
{\scriptsize\color{gray} [p. 14]}\ 在波斯尼亚，同样的困难表现为一条物质性链条。一名行动者这样描述：\par
\switchcolumn
{\scriptsize\color{gray} [p. 14]}\ In Bosnia, the same difficulty appears as a material chain. One activist describes it like this:\par
\switchcolumn*
{\scriptsize\color{gray} [p. 14]}\ 策略是一起去波斯尼亚Biaz的一位公证人那里，由他用公证文书确认此人确实在他面前［……］我为这个人作担保；要使公证文书被认可为可用于法院程序，它还必须有附加证明书［……］这个年轻人就在那里，而且有经过公证的文件证明他在那里，这使此前收集信息、确认他9月曾在的里雅斯特出现的整个过程有了意义。（访谈2）\par
\switchcolumn
{\scriptsize\color{gray} [p. 14]}\ the strategy was to go together to a notary in Bosnia in Biaz who would confirm with a notarial deed the physical presence of this person in front of him […] I acted as guarantor for this person in order for a notarial deed to be recognized as usable in court proceedings, it must also have the apostille […] the fact that this guy was there with a notarized document proving that he was there, made the whole previous process of gathering information that saw him in Trieste in September make sense. (Interview 2)\par
\switchcolumn*
{\scriptsize\color{gray} [p. 15]}\ 在同一过程中，文件经萨拉热窝与比哈奇流转，每一步仍需当地人员陪同。行动者与律师都强调，没有这些当地行动者，提出诉请在实践中就不可能实现（访谈2; 访谈3）。后来的罗马判决记载，授权委托书已“在比哈奇完成翻译并加具附加证明书［……］”（Tribunale di Roma 2023, p. 4）。\par
\switchcolumn
{\scriptsize\color{gray} [p. 15]}\ In the same trajectory, documents moved through Sarajevo and Bihać, local accompaniment remained necessary at each step, and both activists and lawyers stress that without those local actors, the appeal would not have been possible in practice (Interview 2; Interview 3). The later Rome judgment recorded the mandate as “translated and apostilled […] in Bihać” (Tribunale di Roma 2023, p. 4).\par
\switchcolumn*
{\scriptsize\color{gray} [p. 15]}\ 意大利案卷中反复出现的问题，是步骤的先后安排。律师有时必须决定，是否在取得文书查阅权之前先行起诉，尽管这一选择后来可能削弱案件。访谈得出的实际教训是，正式手续与证据未必能以相同速度组装起来。有些案卷中，尽早起诉保留了紧急性；另一些案件则只有先取得针对诉请人的特定记录并相互核对，才变得更有力（访谈1; 访谈3）。\par
\switchcolumn
{\scriptsize\color{gray} [p. 15]}\ In the Italian files, a recurring problem was sequencing. Lawyers sometimes had to decide whether to file before documentary access had been secured, even though that choice could later weaken the case. The practical lesson drawn from the interviews was that formalization and proof could not always be assembled at the same speed. In some files, filing early preserved urgency. In others, the case became stronger only once claimant-specific records could first be obtained and cross-referenced (Interview 1; Interview 3).\par
\switchcolumn*
{\scriptsize\color{gray} [p. 15]}\ 案卷仍在组装时，路径本身也可能改变。在一宗意大利案件中，诉请人必须留在波斯尼亚，因为紧急程序路径取决于这一位置。正如行动者所说：“我们把这些文件交给另一个人，让他带到的里雅斯特，再送给律师。与此同时，他必须留在那里，因为只有你留在那个国家，这件事才行得通”（访谈2）。这种程序逻辑也制造了访谈本身记录下的张力。在波斯尼亚等待，并非中性的法律要求，而是加诸个人的一项要求；此人自己的优先事项，可能恰恰是继续迁移（访谈2）。一名的里雅斯特律师描述，这如何使诉讼策略发生变化：\par
\switchcolumn
{\scriptsize\color{gray} [p. 15]}\ The route itself could also change while the file was still being assembled. In one Italian case, the claimant had to remain in Bosnia because the urgent route depended on that location. As the activist put it, “we send these documents through another person who takes them to Trieste, we get them to the lawyer. In the meantime, he had to stay there because this thing works, if you stay in the country” (Interview 2). That procedural logic also created a tension that the interviews themselves register. Waiting in Bosnia was not a neutral legal requirement, but a demand placed on a person whose own priority might instead be to continue the route (Interview 2). A Trieste lawyer describes how this could create shifts in the litigation strategy:\par
\switchcolumn*
{\scriptsize\color{gray} [p. 15]}\ 就在我们准备起诉时，他离开了［……］我们实际上已经拿到授权委托书，就要提交诉状时，他不再回应了。几天后，他透露自己已在[意大利北部城市]，于是我们取得另一份授权委托书。这时，就直接请求损害赔偿了。所以，他先去委员会面前，随后获准难民身份，然后我们提起诉请。（访谈1）\par
\switchcolumn
{\scriptsize\color{gray} [p. 15]}\ While we were about to file the appeal, he left […] we were practically there with the power of attorney already taken, when we were about to file the appeal, he was no longer responding. After a few days, he revealed that he was already in [northern Italian city] and we collected a different power of attorney. At that point, for compensation for damages directly. So, first he went before the Commission, then he was granted refugee status, and then we appealed. (Interview 1)\par
\switchcolumn*
{\scriptsize\color{gray} [p. 15]}\ 许多意大利案卷根本未能走到起诉。律师开始准备多起诉请，随后决定大部分不予提交。一名律师这样概括这一流失过程：\par
\switchcolumn
{\scriptsize\color{gray} [p. 15]}\ Many Italian files never reached filing at all. Lawyers began several appeals and then decided not to file most of them. One lawyer summarizes that attrition this way:\par
\switchcolumn*
{\scriptsize\color{gray} [p. 15]}\ 有了查阅文件的机会，我们才得以开始。通过与巴尔干路线上行动者合作，我们开始收集大量证词。遗憾的是，其中许多未能转化为诉请，因为仍需要取得正式授权委托书。（访谈1）\par
\switchcolumn
{\scriptsize\color{gray} [p. 15]}\ thanks to this access to documents, we were able to get started. We began to collect a lot of testimonies thanks to the collaboration with activists who were on the Balkan route. Unfortunately, many of them could not be turned into appeals because it was still necessary to obtain a formal power of attorney. (Interview 1)\par
\switchcolumn*
{\scriptsize\color{gray} [p. 15]}\ 这种筛选也由策略考量推动：\par
\switchcolumn
{\scriptsize\color{gray} [p. 15]}\ The selection was also strategically driven:\par
\switchcolumn*
{\scriptsize\color{gray} [p. 15]}\ 我们不想提出一项最终毫无结果的诉请。我们想提出有力的诉请，一项能够起决定作用的诉请。所以，我们着手处理了至少15起案件，准备了至少五份诉状，但因没有把握，没有提交。最后，我们只提交了两份。（访谈1）\par
\switchcolumn
{\scriptsize\color{gray} [p. 15]}\ We didn’t want to file an appeal that would end up going nowhere. We wanted to file a strong appeal, one that would be definitive. So, we started at least 15 cases, we prepared at least five appeals, but we didn’t file them because we weren’t sure. In the end, we only filed two. (Interview 1)\par
\switchcolumn*
{\scriptsize\color{gray} [p. 16]}\ 奥地利将同样的转化工作压缩进短得多的程序时间窗口。一旦找到一名申诉人并保持联系，就必须足够迅速地稳定其叙述，以便在六周内提交申诉。这意味着与律师一起预想可能的问题，通过口译人员澄清时间顺序与地点，并把反复交谈转化为足够详尽的陈述，使其能够在面对由警察塑造的记录时经受检验（访谈5; 访谈8）。\par
\switchcolumn
{\scriptsize\color{gray} [p. 16]}\ Austria compressed the same conversion work into a far shorter procedural window. Once one complainant had been identified and remained in contact, the task was to stabilize the account quickly enough for filing within six weeks. This meant anticipating likely questions with the lawyer, clarifying chronology and place through interpreters, and turning repeated conversations into a statement detailed enough to survive against a police-shaped record (Interview 5; Interview 8).\par
\switchcolumn*
{\scriptsize\color{gray} [p. 16]}\ 斯洛文尼亚的访谈从通道另一侧呈现了同样的工作。受访者说：\par
\switchcolumn
{\scriptsize\color{gray} [p. 16]}\ The Slovenian interview shows the same work from the other side of the corridor. As the interviewee put it,\par
\switchcolumn*
{\scriptsize\color{gray} [p. 16]}\ 我们还把[协调员]及其组织纳入跨团队小组，目的是找到能够提交国内法院的案件。我们基本上能够为意大利人取得大量文件，也能为[协调员]找到可向奥地利法院推进的案件。我们会去庇护接待中心，试着找到那些人，以取得授权委托书并提起诉讼。（访谈9）\par
\switchcolumn
{\scriptsize\color{gray} [p. 16]}\ we also included [coordinator] and his organization in our cross-team group with the purpose of finding cases to put before the national courts. We were basically able to get a lot of documentation for the Italians, and we were able to find cases for [coordinator] to push to the Austrian court. We would go to the asylum home and try to find them in order to get powers of attorney to start lawsuits. (Interview 9)\par
\switchcolumn*
{\scriptsize\color{gray} [p. 16]}\ 这一接力并不对称。记者可能掌握姓名和图像，却没有授权；驻波斯尼亚行动者可能接触得到诉请人，也能在当地行动，却无法进入意大利或奥地利的司法场域；斯洛文尼亚行动者可能拥有文件和进入庇护接待中心的机会，却不负责最终诉状。律师决定诉讼路径及是否起诉，却依赖其他人取得联系、文件、翻译，并维持诉请人的稳定参与。口译员、公证人、当地律师，以及后来在部分意大利案卷中出现的使馆或机场相关人员，只在特定节点介入；但其中任何一步缺失，都可能如同失去诉请人本人一样，确定无疑地使整条链中断（访谈2; 访谈3; 访谈5; 访谈9）。\par
\switchcolumn
{\scriptsize\color{gray} [p. 16]}\ The relay was unequal. Journalists might have names and images but no mandate. Bosnia-based activists might have the claimant and local mobility, but not the Italian or Austrian forum. Slovenian actors might have documents and access to asylum homes, but not the final pleadings. Lawyers decided on route and filing, but depended on others for contact, documents, translation, and claimant stability. Interpreters, notaries, local counsel, and, later in some Italian files, embassy or airport actors entered only at particular points, but a missing step there could stop the chain just as surely as the loss of the claimant himself (Interview 2; Interview 3; Interview 5; Interview 9).\par
\switchcolumn*
{\scriptsize\color{gray} [p. 16]}\ 这些接力关系也十分薄弱，承受风险的程度不均。访谈指出，工作依赖可信赖的当地行动者，人员配置极为有限；能够更迅速在实地行动的个人，与必须权衡公开曝光、资金或自身运作风险的组织之间，也存在差异。尤其是驻波斯尼亚的当地行动者，似乎承担着其他地方处境更优越的行动者所不承担的实际风险与社会风险（访谈2; 访谈5; 访谈9）。\par
\switchcolumn
{\scriptsize\color{gray} [p. 16]}\ These relays were also thin and unevenly exposed. The interviews point to dependence on trusted local actors, very limited staffing, and differences between actors who could move more quickly in the field and organizations that had to weigh visibility, funding, or risks to their own operations. Bosnia-based local actors, in particular, appear to carry practical and social risks that better-positioned actors elsewhere did not (Interview 2; Interview 5; Interview 9).\par
\switchcolumn*
\end{paracol}
\Needspace{10\baselineskip}\section*{裁判 / Judgment}
\begin{paracol}{2}
{\scriptsize\color{gray} [p. 16]}\ 案件组装完毕并提交后，法院裁判将检验中介者的工作。意大利第一案仍然脆弱，因为在起诉前未取得相关记录。最终取得记录时，其中“没有这个人的痕迹”；其他在场者害怕作证，而现有视觉材料也不足以明确识别这名诉请人（访谈1）。一名律师从程序本身解释同样的弱点：在紧急程序中，“不能一边推进，一边补充证据”，因此在第一案中，案卷一旦进入法院，证据问题便成为“无法克服的事实”（访谈3）。罗马上诉判决用更正式的措辞表达了同一点。内政部以意大利、斯洛文尼亚档案中“完全没有诉请人的任何痕迹”为依据，主张“意大利与斯洛文尼亚当局都不知有此人”；合议庭则认定，仍没有充分证据证明“申请人本人遭受了非正式推回”，且证据整体“不足”（Tribunale di Roma 2021b, pp. 1, 4–6）。同一决定还认为，边境暴力监测网络报告及记者文章“不能构成形成确信的唯一依据”（Tribunale di Roma 2021b, p. 5）。\par
\switchcolumn
{\scriptsize\color{gray} [p. 16]}\ Once the case was assembled and filed, adjudication before the court would test the work of the intermediaries. The first Italian case remained vulnerable because relevant records had not been secured before filing. When those records were eventually obtained, there was “no trace of this person,” other people who had been present were afraid to testify, and available visual material did not individualize the claimant strongly enough (Interview 1). One lawyer explains the same weakness from the standpoint of the procedure itself: in the urgent proceedings, “evidence cannot be added as you go along,” so that in the first case, the evidentiary problem became “an insurmountable fact” once the file was already before the court (Interview 3). The Rome appeal judgment makes the same point in more formal terms. The Ministry relied on a “total absence of any trace at all” of the claimant in the Italian and Slovenian archives and argued that he was “unknown both to the Italian authorities and to the Slovenian authorities,” while the collegial court held that there was still no sufficient proof “that the appellant personally underwent an informal pushback” and that the evidentiary picture was “insufficient” (Tribunale di Roma 2021b, pp. 1, 4–6). The same decision also held that the Border Violence Monitoring Network report and journalist article “cannot constitute the sole source of conviction” (Tribunale di Roma 2021b, p. 5).\par
\switchcolumn*
{\scriptsize\color{gray} [p. 17]}\ 意大利第一案的轨迹，还揭示了标准在程序内部的进一步转变。在较早的紧急程序中，法院曾接受相关待遇“按照保全程序所特有的、必然有所降低的证明标准，已获得证明”（Tribunale di Roma 2021a, p. 2）。上诉时，问题变成：这一诉请人本人是否确实经历了所述非正式重新接收；答案是否定的（Tribunale di Roma 2021b, p. 6）。\par
\switchcolumn
{\scriptsize\color{gray} [p. 17]}\ The first Italian trajectory also reveals a more internal shift in standards. In the earlier urgent proceedings, the court had accepted the treatment as “proven according to the necessarily attenuated standard of proof that is typical of precautionary proceedings” (Tribunale di Roma 2021a, p. 2). On appeal, the issue became whether this claimant personally had undergone the narrated informal readmission, and the answer was negative (Tribunale di Roma 2021b, p. 6).\par
\switchcolumn*
{\scriptsize\color{gray} [p. 17]}\ 后来的意大利案件更能保持完整，因为文件序列在起诉前就已取得，并能被法院理解为一条专属于该诉请人的链条。罗马法院依据从的里雅斯特查阅档案取得的材料，包括在“表达申请国际保护的意愿”一项填写“是”的多语种身份登记表，以及同日的重新接收请求、斯洛文尼亚拘留与移交记录（Tribunale di Roma 2023, p. 3）。法院指出，“这些文件合在一起证明”，意大利将其重新接收至斯洛文尼亚，以及随后向克罗地亚的重新接收，均确有发生（Tribunale di Roma 2023, p. 3）。当这一文件序列已建立，波斯尼亚一侧的证人及针对诉请人的视觉材料，就可以强化、而非替代核心证据。因此，法院接收到的案卷，并不是关于沿线暴力的一般叙述，而是一条个体化的序列，将这名诉请人与意大利的处理、向斯洛文尼亚的重新接收及后续移送连接起来（Tribunale di Roma 2023, pp. 3–4; 访谈1）。\par
\switchcolumn
{\scriptsize\color{gray} [p. 17]}\ The later Italian case held together more effectively because the documentary sequence had been secured before filing and could be read by the court as one claimant-specific chain. The Rome court relied on the access-to-files materials from Trieste, including the multilingual identification sheet recording “Expression of the will to apply for international protection” with the answer “yes,” together with the same-day readmission request and the Slovenian custody and handover records (Tribunale di Roma 2023, p. 3). It stated that “The set of these documents demonstrates” both the readmission by Italy to Slovenia and the subsequent readmission to Croatia (Tribunale di Roma 2023, p. 3). Once that documentary sequence was in place, Bosnia-side witnesses and claimant-specific visual material could reinforce rather than substitute for the core proof. The court therefore received the file not as a general account of route violence but as an individualized sequence linking this claimant to Italian handling, Slovenian readmission, and onward removal (Tribunale di Roma 2023, pp. 3–4; Interview 1).\par
\switchcolumn*
{\scriptsize\color{gray} [p. 17]}\ 在奥地利，抵达法院的案卷已经通过反复访谈、翻译与时间顺序澄清建立起来。进入法院后，问题在于，这一诉请人特定的叙述，能否经受住一份稀少、不完整，且已由当局自身事件版本塑造的警察记录的挑战（访谈5; 访谈8）。\par
\switchcolumn
{\scriptsize\color{gray} [p. 17]}\ In Austria, the file that reached court had already been built through repeated interviewing, translation, and chronological clarification. Once before the court, the issue was whether that claimant-specific account could withstand a police record that was thin, incomplete, and already shaped by the authorities’ own version of events (Interview 5; Interview 8).\par
\switchcolumn*
{\scriptsize\color{gray} [p. 17]}\ 奥地利判决明确了法院后来关注什么。在2020年的案件中，一审法院认定，诉请人“多次清晰可闻地说出‘庇护’一词”，却没有警员询问他为何来到奥地利（LVwG Steiermark 2021, p. 8）。行政法院随后指出，关键要求是，警员必须“确认”是否提出庇护请求，并记录相关陈述，以便日后审查，而“这两项都未做到”（VwGH 2022a, p. 6 para. 12）。2021年的案件中，州行政法院同样审查手写记录本身，认为“从手写记录单中看不出”，是否有人问过这群人来奥地利的目的，或是否提出了庇护请求（LVwG Steiermark 2022, p. 7）。\par
\switchcolumn
{\scriptsize\color{gray} [p. 17]}\ The Austrian judgments specify what the court later looked for. In the 2020 case, the first-instance court found that the claimant had “used the word ‘asylum’ several times audibly,” while no officer had asked why he had come to Austria (LVwG Steiermark 2021, p. 8). The VwGH then stated that the decisive requirement was that the officer must “make sure” whether asylum had been sought and must document the relevant statements so that later review was possible, and that “both were omitted” (VwGH 2022a, p. 6 para. 12). In the 2021 case, the LVwG similarly examined the hand protocol itself and held that “it does not emerge from the hand slip” whether anyone had asked what the group wanted in Austria or whether asylum had been requested (LVwG Steiermark 2022, p. 7).\par
\switchcolumn*
{\scriptsize\color{gray} [p. 17]}\ 奥地利后来的案卷，也以更有利的条件进入法院。受访者表示，涉及未成年人的案件更容易维持，因为他留在斯洛文尼亚，可以在那里受到探访，并与斯洛文尼亚一侧的文书记录和法律代理相连接。他们还指出，该案事件结构更加清晰：这群人是为寻求庇护而主动寻找警察局；而后来取得的斯洛文尼亚保护身份，也使质疑其路线或意图的空间更小（访谈5; 访谈7，奥地利，法官; 访谈9）。即便如此，奥地利受访者也强调，存在超出案卷本身的偶然因素，包括警察证词中的矛盾、案件分配，以及具体承办法官的重要性（访谈6; 访谈7，奥地利，法官; 访谈8）。\par
\switchcolumn
{\scriptsize\color{gray} [p. 17]}\ The later Austrian file also came to court from a stronger position. Interviewees describe the case involving the minor as easier to sustain because he remained in Slovenia, could be visited there, and was linked to Slovenian-side documentation and legal representation. They also point to the cleaner event structure in that file, since the group had sought out the police station in order to ask for asylum, and to the later Slovenian protection status, which left less room to contest either route or intention (Interview 5; Interview 7, Austria, judge; Interview 9). Even so, Austrian interviewees also stress contingencies beyond the file itself, including contradictions in police testimony, judicial allocation, and the significance of the particular judge who heard the case (Interview 6; Interview 7, Austria, judge; Interview 8).\par
\switchcolumn*
\end{paracol}
\Needspace{10\baselineskip}\section*{判决之后 / Beyond judgment}
\begin{paracol}{2}
{\scriptsize\color{gray} [p. 18]}\ 判决之后发生的事情，同样具有启发意义，尤其能说明诉讼的限制。在意大利，一项有利裁决本身，并不能把诉请人带回意大利。访谈材料描述了围绕判决继续开展的工作：跟进案件者之间的电话协调、陪同前往机场，以及意大利驻波斯尼亚和黑塞哥维那大使的介入；直到这些工作完成，入境才真正得到保障（访谈2）。\par
\switchcolumn
{\scriptsize\color{gray} [p. 18]}\ What happens after judgment is equally illuminating, particularly as regards the limitations of litigation. In Italy, a favorable ruling did not itself bring the claimant back to Italy. Interview material describes further work around the judgment itself: phone coordination among those following the case, accompaniment to the airport, and intervention by the Italian ambassador in Bosnia and Herzegovina before entry was actually secured (Interview 2).\par
\switchcolumn*
{\scriptsize\color{gray} [p. 18]}\ 此外，据受访者所述，判决之后，国家行动者进行了调整。在奥地利，受访者将这些案件与后来部委的审查，以及一项澄清庇护申请也可默示提出的法令联系起来（访谈6）。同一受访者描述，内政部内部仔细研究判决，以改变未来官方行为，并减少再次在法院面临追究的风险（访谈6）。在意大利，受访者描述了管控向更难追溯的形式转变，包括依据申根规则拒绝入境、三边警务合作，以及在链条更前端进行干预。一名律师评论，非正式重新接收本来已很难证明，而要证明意大利警员曾参与克罗地亚境内的推回行动，则变得“实际上不可能”（访谈3）。\par
\switchcolumn
{\scriptsize\color{gray} [p. 18]}\ Furthermore, according to interviewees, what followed the judgments was adjustment by state actors. In Austria, interviewees link the cases to later ministry scrutiny and to a decree clarifying that asylum applications could also be made implicitly (Interview 6). The same interviewee describes the judgments as being closely studied inside the Ministry of the Interior with a view to changing future official conduct and limiting further exposure in court (Interview 6). In Italy, interviewees describe movement toward less traceable forms of control, including Schengen refusals, trilateral police cooperation, and interventions earlier in the chain. One lawyer remarks that, while informal readmissions were already difficult to prove, showing that Italian officers had participated in pushback operations in Croatia became “practically impossible” (Interview 3).\par
\switchcolumn*
{\scriptsize\color{gray} [p. 18]}\ 此外，奥地利行动者描述，裁决后他们定期交流，以追踪做法的变化并传递信息；斯洛文尼亚和意大利受访者则描述了更长期的跨境合作，通过这些合作，记录、联系人和取证策略被带入后续案件（访谈5; 访谈3; 访谈9）。\par
\switchcolumn
{\scriptsize\color{gray} [p. 18]}\ Furthermore, Austrian activists describe regular exchange after the rulings in order to follow changing practices and pass on information, while Slovenian and Italian interviewees describe longer-standing cross-border cooperation through which records, contacts, and evidentiary tactics circulated into later cases (Interview 5; Interview 3; Interview 9).\par
\switchcolumn*
{\scriptsize\color{gray} [p. 18]}\ 访谈还把诉讼置于更广泛的行动方式之中。奥地利受访者在谈及法律工作时，也谈到记录、公开报告及即时干预；非政府组织协调员明确把诉讼称为让推回变得可见的多种工具之一（访谈5; 访谈6）。意大利行动者同样将法律程序与公众认知工作并列（访谈2）。与此同时，这些案件背后的支持基础仍很薄弱，依赖志愿劳动、信任，以及可能因人员更替、疲惫或沿线条件变化而削弱的关系（访谈2; 访谈3; 访谈5; 访谈9）。\par
\switchcolumn
{\scriptsize\color{gray} [p. 18]}\ The interviews also place litigation within a wider repertoire of action. Austrian interviewees describe legal work alongside documentation, public reporting, and immediate intervention, and the NGO coordinator explicitly characterizes litigation as one vehicle among others for making pushbacks visible (Interview 5; Interview 6). The Italian activist similarly places legal procedure alongside public awareness work (Interview 2). At the same time, the infrastructure behind these cases remained thin, dependent on volunteer labor, trust, and relations that could weaken through turnover, exhaustion, or changing conditions on the route (Interview 2; Interview 3; Interview 5; Interview 9).\par
\switchcolumn*
\end{paracol}
\Needspace{10\baselineskip}\section*{讨论 / Discussion}
\begin{paracol}{2}
{\scriptsize\color{gray} [p. 18]}\ 如经验材料所示，推回诉讼始于案件错位状态，因为受到挑战的做法本身，拆散并动摇了通常必须汇合、才能使案件可诉的要素。潜在诉请人被移离司法场域，往往继续跨境迁移，面临疲惫、不安全及优先事项变化。痕迹同样分散且不稳定：它散布在不同司法管辖区、机构、手机、证人叙述、警察档案、医疗记录、重新接收文件与非正式记录中，同时往往仍然不完整、无法获取、可以被否认，或容易遭抹除。最后，诉讼路径本身也在不同地点、行动者、程序步骤与法律体系之间碎片化；失去联系、缺少授权、错过期限、证据缺口，或诉请人位置变化，都可能使其断裂。因此，案件错位指这样一种状态：诉请人、痕迹与程序路径无法以足以使案件向法院移动的形式汇合。故而，在较狭义的法律转译开始之前，中介者必须首先开展识别、重新连接并暂时稳定这些分散要素的工作。本文考察的案件也确实表明，推回案件往往因潜在案件的要素无法汇合而失败。这种困难从一开始就被制造出来，并发生在暴力、疲惫、不透明与刑事化条件之下（Alpes and Baranowska 2025; Baranowska et al. 2025; Hameršak et al. 2020; Trucco 2023; Zocchi 2024）。因此，研究发现将注意力集中于诉请提出之前及提出初期阶段（Olesen and Hammerslev 2023）。\par
\switchcolumn
{\scriptsize\color{gray} [p. 18]}\ As our empirical data demonstrate, pushback litigation began from a condition of case misalignment because the very practice being challenged separated and unsettled the elements that must normally converge for a case to become actionable. The potential claimants were displaced from the forum, often moved onward across borders, exposed to exhaustion, insecurity, and changing priorities. The trace was likewise dispersed and destabilized. It was scattered across jurisdictions, institutions, phones, witness accounts, police files, health records, readmission papers, and informal documentation, while often remaining incomplete, inaccessible, deniable, or vulnerable to erasure. The litigation path, finally, was itself fragmented across places, actors, procedural steps, and legal systems, and it could become broken by the loss of contact, the absence of a mandate, a missed deadline, an evidentiary gap, or a change in the claimant’s location. Case misalignment, then, refers to a condition in which claimant, trace, and procedural route do not coincide in a form that cases can move toward courts. Before legal translation in the narrower sense can begin, intermediaries must therefore first undertake the work of identifying, reconnecting, and temporarily stabilizing these disaggregated elements. And the cases examined here indeed show that pushback cases often failed because the elements of potential cases could not be made to coincide. That difficulty was produced from the outset and under conditions of violence, exhaustion, opacity, and criminalization (Alpes and Baranowska 2025; Baranowska et al. 2025; Hameršak et al. 2020; Trucco 2023; Zocchi 2024). The findings, therefore, focus attention on the pre-claim and early-claim phases (Olesen and Hammerslev 2023).\par
\switchcolumn*
{\scriptsize\color{gray} [p. 19]}\ 经验材料指出了中介工作的三个主要领域，对应于他们所组装案件的基本要素，即痕迹、关系与诉讼路径。综合来看，这些领域显示了一个分阶段、但并非整齐线性的过程，其中不同实践识别、连接并稳定这些要素。面对错位，中介者首先从若干起点发现、接收并筛选可能的案件：一份证词、一个电话号码、一张照片、一张重新接收单、一份医疗记录，或接待点出现的矛盾。此后，必须识别一个人、接触他、支持他，并保持沟通，使其参与稳定下来，才能尝试下一步：把此人与某项痕迹连接，通常通过文件交叉核对、地点核验或档案查阅申请，再把这种联结转化为适合特定审理机关的程序形式。当这一连接或对齐稳定下来，一份能够在一定程度上被处理的案卷才出现。\par
\switchcolumn
{\scriptsize\color{gray} [p. 19]}\ Our empirical material indicates three main areas of intermediary work that relate to the basic elements of the cases they assemble, namely work on traces, relationships, and litigation paths. When read together, they indicate a staged, though not neatly linear, process with distinct practices that identify, connect, and stabilize those elements. Confronted with misalignment, intermediaries began by detecting, receiving, and screening possible cases from those points: a testimony, a phone number, a photograph, a readmission slip, a health record, or an inconsistency at a reception point. From there, a person had to be identified, reached, supported, kept in communication, thus stabilized, so that the next step could be attempted: connecting that person to a trace, often through documentary cross-referencing, place verification, or access-to-files requests, and converting that linkage into a forum-specific procedural form. Where that connection or alignment was stabilized, a file emerged that could be processed to some extent.\par
\switchcolumn*
{\scriptsize\color{gray} [p. 19]}\ 意大利与奥地利的比较表明，同一基本机制在不同国内法律约束和机会下运作，而这些差异改变了已组装案件得以推进的节点。在意大利，案件组装跨越更远距离，通常也持续更长时间。主要压力在于授权的正式化、文书获取及证据时序。由于诉请人往往位于意大利领土之外，授权本身就成为包含公证、证明、认证、翻译及实际陪同的跨国链条。同时，较有力的意大利案卷，是那些中介者能在起诉之前建立文件序列的案件；他们通常交叉核对意大利与斯洛文尼亚记录，并利用健康筛查材料等文书切入点。因此，意大利程序既制造障碍，也提供有限机会：某些紧急路径若过早启动，就十分脆弱；而另一些案卷，却可以通过推迟、重新安排步骤，或后来从以返回为目标的救济转向赔偿请求，变得更强。奥地利则将同一机制压缩进短得多的程序时段。在那里，关键压力来自迅速恢复诉请的需要、较短的起诉期限、范围狭窄的救济，以及可预见的需要——必须挑战一份已由警察塑造的官方记录。从这个意义上说，法律机会结构很早就开始发挥作用，而且就在组装工作本身的内部。\par
\switchcolumn
{\scriptsize\color{gray} [p. 19]}\ The comparison between Italy and Austria shows that the same basic mechanism operated under different domestic legal constraints and opportunities, and that these differences shifted the point at which the assembled case could move. In Italy, case assembly was stretched across a greater distance and usually over a longer period. The central pressures lay in mandate formalization, documentary access, and evidentiary timing. Because claimants were often outside Italian territory, authorization itself became a transnational chain of notarization, certification, legalization, translation, and physical accompaniment. At the same time, the stronger Italian files were those in which intermediaries could build a documentary sequence before filing, often by cross-referencing Italian and Slovenian records and exploiting documentary levers such as health-screening material. The Italian procedure thus created both obstacles and a limited opportunity: while some urgent routes were highly vulnerable if filed too early, other files could become stronger through delay, resequencing, or a later shift from a return-oriented remedy to a compensation claim. Austria compressed the same mechanism into a much shorter procedural period. There, the decisive pressures lay in the need to recover claims rapidly, in short filing deadlines, in narrow remedies, and in the anticipated need to contest an official record already shaped by the police. In this sense, legal opportunity structures became relevant early, and inside the assembly work itself.\par
\switchcolumn*
{\scriptsize\color{gray} [p. 20]}\ 协作成为中介者为组装案件所选择的安排，因为案件本身从一开始就在地理与制度上被拆散。信息与案卷通过不对称的接力流转，每个行动者提供下一位行动者所欠缺的要素。记者与路线沿线行动者，常将最初碎片带入视野；接待与协调行动者，将其他信息碎片转化为潜在案件库。位于诉请人被移送地的行动者、记者及当地联系人，承担许多最艰难的工作，例如在营地、占住的空屋与非正式聚居点寻找诉请人，在其继续迁移及更换电话号码时维持沟通，解释程序，陪伴人们克服实际困难，有时仅靠使诉请人仍可联系，就维系住案卷，或在后来成为证人。斯洛文尼亚行动者占据结构上的关键位置，因为他们往往能够把个人重新连接到官僚记录、重新接收记录、健康表格、警察档案，以及进入庇护接待中心的机会；这些为意大利与奥地利的诉讼提供了支点。链条的另一部分由翻译人员、公证人和当地律师组成，他们使授权稳定下来。最后，律师选择程序路径，判断所组装的链条是否足够有力，并将其转化为诉状；但在联系、文书及稳定工作方面，他们依赖他人。因此，总体而言，这一接力以不对称的依赖和不均等的风险承担为特征。在意大利，链条更长，更侧重授权，延伸也更远；在奥地利，则更加紧凑，依赖热线行动者、翻译人员、非政府组织及一名律师，在短暂的申诉窗口内迅速协调。\par
\switchcolumn
{\scriptsize\color{gray} [p. 20]}\ Collaboration became the arrangement that intermediaries chose for case assembly, since the case itself was geographically and institutionally disassembled from the outset. Information and files moved through an unequal relay in which each actor supplied the element that the next actor lacked. Journalists and route-level actors often brought the first fragments to the fore. Reception and coordination actors turned other fragments of information into a pool of possible cases. Activists and journalists in locations to which claimants had been displaced, and local contacts, carried much of the hardest work, e.g., by finding claimants in camps, squats, and informal settlements, maintaining communication despite onward movement and changing numbers, explaining the process, accompanying people through practical challenges, and at times holding the file together simply by keeping a claimant reachable or later becoming witnesses. Slovenian actors occupied a structurally pivotal position because they could often reconnect individuals to bureaucratic records, readmission records, health forms, police files, and access to asylum homes, which anchored litigation in Italy and Austria. Another part of the chain consisted of translators, notaries, and local counsel, who stabilized authorizations. Finally, lawyers selected procedural routes, judged whether the assembled chain was strong enough, and turned it into pleadings, but depended on others for contact, documentation, and stabilization. In sum, the relay was therefore marked by asymmetrical dependence and uneven exposure. In Italy, it was longer and more mandate-heavy, and it stretched further. In Austria, it was more compressed and relied on rapid coordination among hotline actors, translators, NGOs, and a single lawyer, all within a short complaint window.\par
\switchcolumn*
{\scriptsize\color{gray} [p. 20]}\ 因此，在某种意义上，这一环境也使法律行动所需的能力碎片化，令案件组装必须依靠行动者协作（Cummings 2020; Davies et al. 2023; Kawar 2015）。而这种协作，又促成一个小型、脆弱的支持基础在更广泛场域内形成；该场域还包括监测、陪同、新闻报道及公开谴责等活动，诉讼由此呈现为众多行动方式之一（Alpes and Baranowska 2025; Cummings 2020; Hameršak et al. 2020; Kawar 2015; Trucco et al. 2023）。\par
\switchcolumn
{\scriptsize\color{gray} [p. 20]}\ In some ways, then, the setting had also fragmented the capacities needed for legal action such that case assembly required collaboration among actors (Cummings 2020; Davies et al. 2023; Kawar 2015). Moreover, this collaboration had contributed to the emergence of a small and fragile infrastructure within a wider field that includes, e.g., monitoring, accompaniment, journalism, and denunciation, and where litigation thus appears as one repertoire among others (Alpes and Baranowska 2025; Cummings 2020; Hameršak et al. 2020; Kawar 2015; Trucco et al. 2023).\par
\switchcolumn*
{\scriptsize\color{gray} [p. 20]}\ 本研究进一步确认，抵达冰山顶端仍属例外，并有助于理解为何推回诉讼尤其如此。除了充满阻碍的环境本身的一般运作，材料还揭示法律动员自身如何筛选案件：发现工作产生许多潜在案件，但大多数从未成为正式提交的案件。有些因无法识别诉请人而失败；另一些则因联系无法维持、诉请人不再愿意参与、文件无法取得，或案卷不再适合可用程序路径而失败。诉请人的优先事项、中介者的能力，以及对案卷能否经受早期审查的判断，共同推动了筛选（Lejeune and Spire 2021; Miaz et al. 2021）。在意大利，律师按案件可持续性与证据强度筛选，因为薄弱诉请可能使案件崩溃，并损害未来诉讼。在奥地利，中介者更倾向于把案件流失归因于组织能力有限及支持结构极其薄弱；这呼应了有关持续权利诉讼依赖能够承担追诉成本的基础设施的论述（Epp 2011; Southworth 2000）。\par
\switchcolumn
{\scriptsize\color{gray} [p. 20]}\ Our study further confirms that reaching the top of the iceberg remains exceptional, and it contributes to understanding why this is particularly so in pushback litigation. Beyond the general workings of the hostile setting itself, the material also reveals how legal mobilization itself filtered cases, since detection produced many possible cases, but most never became filed cases. Some failed because no claimant could be identified. Others failed because contact could not be maintained, claimants were no longer willing, documents could not be obtained, or the file no longer fit the procedural route available. Selection occurred through a mix of claimant priorities, the capacity of intermediaries, and judgments about whether the file could survive early scrutiny (Lejeune and Spire 2021; Miaz et al. 2021). In Italy, lawyers screened cases for sustainability and the strength of evidence, as weak claims risked collapsing the case and could undermine future litigation. In Austria, intermediaries attributed attrition more strongly to their limited organizational capacity and a very thin support structure, which recalls accounts according to which sustained rights litigation depends on infrastructures able to absorb the costs of pursuing cases (Epp 2011; Southworth 2000).\par
\switchcolumn*
{\scriptsize\color{gray} [p. 20]}\ 此外，这种诉请提出之前的筛选，将法律动员研究与有关边境记录、报告及知识层面边境工作的研究连接起来（Davies et al. 2023; Hameršak 2021; Hameršak et al. 2020）。我们的分析表明，仅有可见性并不足够。监测、报告及路线沿线的记录，有时是相关起点，但单凭它们无法转向法院。它们的重要性，在于实地行动者有时能够将其与诉请人及正式路径连接，从而转化为诉讼链条。此外，当行动者发现可信的法院机会时，现有的人道、监测与行动者支持基础，有时会被选择性地转用于新目的。\par
\switchcolumn
{\scriptsize\color{gray} [p. 20]}\ Furthermore, this pre-claim filtration connects legal mobilization research to studies on border documentation, reporting, and epistemic borderwork (Davies et al. 2023; Hameršak 2021; Hameršak et al. 2020). Our analysis suggests that visibility alone is insufficient. Monitoring, reporting, and route-level documentation were, at times, relevant starting points, but they were not enough on their own to turn to court. Their significance lay in their potential to be converted, at times, into litigation chains by actors on the ground who tied them to a claimant and a formal path. Moreover, existing humanitarian, monitoring, and activist infrastructures were sometimes selectively repurposed when actors identified a plausible court opportunity.\par
\switchcolumn*
{\scriptsize\color{gray} [p. 21]}\ 因此，访谈中可见的行动者群体，也比以诉讼为基础的样本初看所显示的范围更广。诉讼嵌入这一更广阔且不均衡的场域中。访谈还在有限程度上提供材料，说明行动者、记者及其他非律师中介者为何可能倾向于转向法律。一些人此前已有相关领域经验，又反复遭遇似乎不可能仅靠谴责而改变的虐待模式。不过，这批访谈材料无法支持任何判断，说明这种诉讼取向在更广泛的巴尔干路线声援场域中共享到何种程度，更不能说明遭受推回者本人是否具有同样取向。\par
\switchcolumn
{\scriptsize\color{gray} [p. 21]}\ The group of actors visible in the interviews was thus also wider than the litigation-based sample would initially suggest. Litigation was nested within that wider and uneven field. The interviews also include, although only to a limited extent, data on why activists, journalists, and other non-lawyer intermediaries were potentially inclined to turn to law. Some had prior experience in the field, combined with repeated exposure to patterns of abuse that seemed unlikely to shift through denunciation alone. However, the interview corpus does not permit any statements about how far such an orientation toward litigation was shared across the wider Balkan-route solidarity field, still less among people affected by pushbacks themselves.\par
\switchcolumn*
{\scriptsize\color{gray} [p. 21]}\ 访谈数据还提供了对流动人群能动性如何被看待、如何被中介的认识。由于样本未包括诉请人或潜在诉请人，材料只能间接呈现他们对诉讼与法院的态度。不过，访谈确实反复显示诉请人调整优先事项、拒绝及作出策略选择：有人继续迁移或拒绝等待，拒绝作证、退出程序，也有人选择保持可联系、取回文件并继续诉讼。因此，在程序上表现为案件流失的现象，也可被看作严峻约束下能动性的表达。同时，材料确实揭示了不对称权力关系。若干受访者明确反对强迫他人，接受退出，也承认许多人有诉讼以外的优先事项。然而，中介者仍然筛查可行性、选择程序路径、判断风险是否值得承担，并有时鼓励人们参与那些价值部分在于政治及制度层面的申诉。\par
\switchcolumn
{\scriptsize\color{gray} [p. 21]}\ The interview data also contain insights into how the agency of people on the move was viewed and mediated. Since the interview sample did not include claimants or potential claimants, the material offers only indirect access to their attitudes toward litigation and courts. But the interviews do show repeated reprioritization, refusal, and tactical choice of claimants: people moved on or declined to wait, refused to testify, withdrew from proceedings, or chose to remain reachable, retrieve documents, and pursue litigation. What appears procedurally as attrition could therefore also be seen as an expression of agency under considerable constraint. At the same time, the material does reveal asymmetrical power relationships. Several interviewees explicitly rejected forcing people, accepted withdrawal, and recognized that many had priorities other than litigation. Yet, intermediaries also screened viability, selected procedural routes, judged whether risks were worth bearing, and sometimes encouraged participation in complaints whose value was partly political and systemic.\par
\switchcolumn*
{\scriptsize\color{gray} [p. 21]}\ 因此，这批材料凸显的主要问题，并非广泛的“法律救世主义”，而是在一个中介者必然塑造哪些诉请人行为能够被法律理解的领域中，存在较为有限的代表关系不对称风险，以及偶发的法律中心主义（Bonizzoni and Odasso 2024; Bonjour 2025）。与此同时，不应把行动者反复求助法律，误解为对法院的一般性信仰。他们也将诉讼与预防、热线干预、记录、接待、陪同、媒体工作和公开报告并置，而且往往明确承认其直接收益有限、所施加的负担、诉请人参与的脆弱性，以及当局多么容易通过调整来绕开不利裁决。事实上，按法律中介者的叙述，司法干预使向斯洛文尼亚的系统性推回停止；但国家行动者也以不那么可见的移民管控调整作出回应，减少未来受到法律挑战的可能（Thomas and Hathaway 2014; Lavenex 2023; Mann 2013）。最后，法律上的胜利，并不必然使个别诉请人的整体处境得到稳定改善。\par
\switchcolumn
{\scriptsize\color{gray} [p. 21]}\ The stronger concern disclosed by the corpus, however, is thus not a broad form of “legal saviorism,” but a rather limited risk of representational asymmetry and occasional legal-centrism within a field where intermediaries necessarily shape which forms of claimant conduct become legally intelligible at all (Bonizzoni and Odasso 2024; Bonjour 2025). At the same time, the actors’ repeated resort to law should not be mistaken for a generalized faith in courts. They also situated litigation alongside prevention, hotline intervention, documentation, reception work, accompaniment, media work, and public reporting, and were often explicit about its limited direct payoff, the burdens it imposed, the fragility of claimant participation, and the ease with which authorities could adapt around adverse rulings. Indeed, in the account of legal intermediaries, judicial interventions led to a halt in systematic pushbacks toward Slovenia. But state actors also responded with less visible adaptations of migration control practices, reducing their exposure to future legal challenges (Thomas and Hathaway 2014; Lavenex 2023; Mann 2013). Finally, a legal victory did not necessarily produce a stable improvement in the individual claimant’s broader situation.\par
\switchcolumn*
\end{paracol}
\Needspace{10\baselineskip}\section*{结论 / Conclusion}
\begin{paracol}{2}
{\scriptsize\color{gray} [p. 22]}\ 本文表明，推回诉讼的困难在于：在事件或文件能够被转译为法律案卷之前，首先就必须努力使一个案件得以成立。在充满阻碍的跨国环境中，制造伤害的同一套制度，也使潜在诉请人、痕迹及可能寻求补救的程序路径分散且不稳定，因而法律抗争从所谓“案件错位”的状态起步。“案件组装”概念，捕捉了在这一极早阶段，为识别、重新连接并暂时稳定一项诉请的破碎要素，所需开展的独特中介工作。因此，推回案件中的法律中介活动，关乎构成使裁判成为可能的条件。意大利与奥地利的比较进一步说明，这一工作既受到国内程序环境的塑造，又依赖跨境协作。\par
\switchcolumn
{\scriptsize\color{gray} [p. 22]}\ This article has demonstrated that the difficulty of pushback litigation lies in the struggle to bring a case into being at all, before events or documents can be translated into a legal file. In a hostile transnational setting, the same regime that produces the harm also disperses and destabilizes the potential claimant, the trace, and the procedural path through which redress might be sought, such that legal contestation begins from a condition denoted as “case misalignment.” The concept of “case assembly” captures the distinct intermediary work required at this very early point to identify, reconnect, and temporarily stabilize the fractured elements of a claim. Legal intermediation in pushback cases is thus a matter of constituting the conditions under which adjudication becomes possible. The comparison between Italy and Austria further shows how this work is shaped by domestic procedural settings, while depending on cross-border collaboration.\par
\switchcolumn*
\end{paracol}
\Needspace{10\baselineskip}\section*{数据可用性声明 / Data availability statement}
\begin{paracol}{2}
{\scriptsize\color{gray} [p. 22]}\ 数据在欧洲研究理事会（ERC）ACCESS项目内收集。案头研究所得数据以纸面形式提供。访谈数据的可获取性受到限制，并未公开。不过，附录列出了所进行的访谈及一些相关基本信息。\par
\switchcolumn
{\scriptsize\color{gray} [p. 22]}\ Data have been collected within the ERC ACCESS project. Data collected through desk research are available on paper. Restrictions apply to the availability of data from interviews, which are not publicly available. However, Appendix contains the list of interviews conducted with some basic relevant information.\par
\switchcolumn*
\end{paracol}
\Needspace{10\baselineskip}\section*{致谢 / Acknowledgements}
\begin{paracol}{2}
{\scriptsize\color{gray} [p. 22]}\ 感谢参与者愿意为本研究作出贡献，并为我们投入时间。感谢讨论、修订本文初稿的同事，尤其感谢Madalina Moraru、Sara Mariella Lambertini、Chiara Scissa、Stefen Thierse及Anuscheh Farahat提出富有洞见的意见与想法。最后，衷心感谢匿名审稿人认真评估本研究，并提出宝贵建议，使稿件得以加强。\par
\switchcolumn
{\scriptsize\color{gray} [p. 22]}\ We are grateful to the participants for their willingness to contribute to the research and the time they dedicated to us. We are also grateful to the colleagues who discussed and revised the draft versions of this paper, and particularly to Madalina Moraru, Sara Mariella Lambertini, Chiara Scissa, Stefen Thierse, and Anuscheh Farahat for their insightful comments and ideas. Finally, we sincerely thank the anonymous reviewers for their careful evaluation of this work and for their valuable suggestions, which have strengthened the manuscript.\par
\switchcolumn*
\end{paracol}
\Needspace{10\baselineskip}\section*{资助声明 / Funding statement}
\begin{paracol}{2}
{\scriptsize\color{gray} [p. 22]}\ 本研究获欧洲研究理事会2022年启动基金项目“国际难民法的守门人——法院在塑造庇护准入中的作用”（101078683）资助。本研究全部或部分经费来自奥地利科学基金（FWF），项目标识为10.55776/PAT3567024。\par
\switchcolumn
{\scriptsize\color{gray} [p. 22]}\ This work was funded by the European Research Council 2022 Starting Grant Gatekeepers to International Refugee Law – The Role of Courts in Shaping Access to Asylum (101078683). This research was funded in whole or in part by the Austrian Science Fund (FWF) 10.55776/PAT3567024.\par
\switchcolumn*
\end{paracol}
\Needspace{10\baselineskip}\section*{利益冲突 / Conflicts of interest}
\begin{paracol}{2}
{\scriptsize\color{gray} [p. 22]}\ 不存在竞争性利益。\par
\switchcolumn
{\scriptsize\color{gray} [p. 22]}\ There are no competing interests.\par
\switchcolumn*
\end{paracol}
\Needspace{10\baselineskip}\section*{参考文献 / References}
\noindent\begin{minipage}{\textwidth}\begin{minipage}[t]{\dimexpr.5\textwidth-4mm\relax}{\scriptsize\color{gray} [p. 22]}\ Ahlin, Eileen M. 2019. “Semi-Structured Interviews with Expert Practitioners: Their Validity and Significant Contribution to Translational Research.” SAGE Research Methods Cases. DOI: 10.4135/9781526466037 〔题名或条目译释：与专业实践者的半结构式访谈：其效度及对转化研究的重要贡献。〕\par\end{minipage}\hfill\begin{minipage}[t]{\dimexpr.5\textwidth-4mm\relax}{\scriptsize\color{gray} [p. 22]}\ Ahlin, Eileen M. 2019. “Semi-Structured Interviews with Expert Practitioners: Their Validity and Significant Contribution to Translational Research.” SAGE Research Methods Cases. DOI: 10.4135/9781526466037\par\end{minipage}\end{minipage}\par
\noindent\begin{minipage}{\textwidth}\begin{minipage}[t]{\dimexpr.5\textwidth-4mm\relax}{\scriptsize\color{gray} [p. 22]}\ Alagna, Federico. 2025. “Pro-migrant Civil Society Organisations and the Law: Patterns of Legal Mobilisation at EU Borders.” \textit{International Migration Rev.} Research Note: 1–18. 〔题名或条目译释：支持移民的公民社会组织与法律：欧盟边境的法律动员模式。〕\par\end{minipage}\hfill\begin{minipage}[t]{\dimexpr.5\textwidth-4mm\relax}{\scriptsize\color{gray} [p. 22]}\ Alagna, Federico. 2025. “Pro-migrant Civil Society Organisations and the Law: Patterns of Legal Mobilisation at EU Borders.” \textit{International Migration Rev.} Research Note: 1–18.\par\end{minipage}\end{minipage}\par
\noindent\begin{minipage}{\textwidth}\begin{minipage}[t]{\dimexpr.5\textwidth-4mm\relax}{\scriptsize\color{gray} [p. 22]}\ Albiston, Catherine R., Edelman, Lauren B. and Milligan, Joy. 2014. “The Dispute Tree and the Legal Forest.” \textit{Annual Rev. of Law and Social Science} 10: 105–31. doi:10.1146/annurev-lawsocsci-110413-030826. 〔题名或条目译释：纠纷之树与法律之林。〕\par\end{minipage}\hfill\begin{minipage}[t]{\dimexpr.5\textwidth-4mm\relax}{\scriptsize\color{gray} [p. 22]}\ Albiston, Catherine R., Edelman, Lauren B. and Milligan, Joy. 2014. “The Dispute Tree and the Legal Forest.” \textit{Annual Rev. of Law and Social Science} 10: 105–31. doi:10.1146/annurev-lawsocsci-110413-030826.\par\end{minipage}\end{minipage}\par
\noindent\begin{minipage}{\textwidth}\begin{minipage}[t]{\dimexpr.5\textwidth-4mm\relax}{\scriptsize\color{gray} [p. 22]}\ Alpes, Maybritt Jill. 2026. “Truth-Claim Capital at Borders: The Role of Trust in the State at a Regional Human Rights Court.” \textit{Social \& Legal Studies} 35 (3) 313–330 doi:10.1177/09646639251364483. 〔题名或条目译释：边境上的真相主张资本：地区人权法院中对国家的信任所起的作用。〕\par\end{minipage}\hfill\begin{minipage}[t]{\dimexpr.5\textwidth-4mm\relax}{\scriptsize\color{gray} [p. 22]}\ Alpes, Maybritt Jill. 2026. “Truth-Claim Capital at Borders: The Role of Trust in the State at a Regional Human Rights Court.” \textit{Social \& Legal Studies} 35 (3) 313–330 doi:10.1177/09646639251364483.\par\end{minipage}\end{minipage}\par
\noindent\begin{minipage}{\textwidth}\begin{minipage}[t]{\dimexpr.5\textwidth-4mm\relax}{\scriptsize\color{gray} [p. 22]}\ Alpes, Maybritt Jill and Baranowska, Grażyna. 2025. “The Politics of Legal Facts: The Erasure of Pushback Evidence from the European Court of Human Rights.” \textit{Law \& Social Inquiry} 50(1): 225–48. doi:10.1017/lsi.2024.27. 〔题名或条目译释：法律事实的政治：欧洲人权法院中推回证据的抹除。〕\par\end{minipage}\hfill\begin{minipage}[t]{\dimexpr.5\textwidth-4mm\relax}{\scriptsize\color{gray} [p. 22]}\ Alpes, Maybritt Jill and Baranowska, Grażyna. 2025. “The Politics of Legal Facts: The Erasure of Pushback Evidence from the European Court of Human Rights.” \textit{Law \& Social Inquiry} 50(1): 225–48. doi:10.1017/lsi.2024.27.\par\end{minipage}\end{minipage}\par
\noindent\begin{minipage}{\textwidth}\begin{minipage}[t]{\dimexpr.5\textwidth-4mm\relax}{\scriptsize\color{gray} [p. 22]}\ Amnesty International. 2019. “Pushed to the edge: Violence and abuse against refugees and migrants along Balkan Route”. Index Number: EUR 05/9964/2019, March 2019, London, UK. 〔题名或条目译释：被推向边缘：巴尔干路线沿线针对难民与移民的暴力和虐待。〕\par\end{minipage}\hfill\begin{minipage}[t]{\dimexpr.5\textwidth-4mm\relax}{\scriptsize\color{gray} [p. 22]}\ Amnesty International. 2019. “Pushed to the edge: Violence and abuse against refugees and migrants along Balkan Route”. Index Number: EUR 05/9964/2019, March 2019, London, UK.\par\end{minipage}\end{minipage}\par
\noindent\begin{minipage}{\textwidth}\begin{minipage}[t]{\dimexpr.5\textwidth-4mm\relax}{\scriptsize\color{gray} [p. 22]}\ Association for Legal Studies on Immigration. 2023. “Gli interventi di cooperazione di polizia che “inibiscono” l’ingresso in Italia ai cittadini stranieri”. \url{https://www.asgi.it/medea/gli-interventi-di-cooperazione-di-polizia-che-inibiscono-lingresso-in-italia-ai-cittadini-stranieri/} (accessed September 29 , 2025). 〔题名或条目译释：‘阻止’外国公民进入意大利的警务合作行动。〕\par\end{minipage}\hfill\begin{minipage}[t]{\dimexpr.5\textwidth-4mm\relax}{\scriptsize\color{gray} [p. 22]}\ Association for Legal Studies on Immigration. 2023. “Gli interventi di cooperazione di polizia che “inibiscono” l’ingresso in Italia ai cittadini stranieri”. \url{https://www.asgi.it/medea/gli-interventi-di-cooperazione-di-polizia-che-inibiscono-lingresso-in-italia-ai-cittadini-stranieri/} (accessed September 29 , 2025).\par\end{minipage}\end{minipage}\par
\noindent\begin{minipage}{\textwidth}\begin{minipage}[t]{\dimexpr.5\textwidth-4mm\relax}{\scriptsize\color{gray} [p. 23]}\ Association for Legal Studies on Immigration. 2025. “Country Report: Italy.” July 2025. Asylum Information Database. 〔题名或条目译释：国别报告：意大利。〕\par\end{minipage}\hfill\begin{minipage}[t]{\dimexpr.5\textwidth-4mm\relax}{\scriptsize\color{gray} [p. 23]}\ Association for Legal Studies on Immigration. 2025. “Country Report: Italy.” July 2025. Asylum Information Database.\par\end{minipage}\end{minipage}\par
\noindent\begin{minipage}{\textwidth}\begin{minipage}[t]{\dimexpr.5\textwidth-4mm\relax}{\scriptsize\color{gray} [p. 23]}\ Aulsebrook, Gigi, Gruber, Natalie and Pawson, Melissa. 2021. “Pushbacks on the Balkan Route: A Hallmark of EU Border Externalization.” \textit{Forced Migration Rev.} 68: 13–15. 〔题名或条目译释：巴尔干路线上的推回：欧盟边境外部化的标志。〕\par\end{minipage}\hfill\begin{minipage}[t]{\dimexpr.5\textwidth-4mm\relax}{\scriptsize\color{gray} [p. 23]}\ Aulsebrook, Gigi, Gruber, Natalie and Pawson, Melissa. 2021. “Pushbacks on the Balkan Route: A Hallmark of EU Border Externalization.” \textit{Forced Migration Rev.} 68: 13–15.\par\end{minipage}\end{minipage}\par
\noindent\begin{minipage}{\textwidth}\begin{minipage}[t]{\dimexpr.5\textwidth-4mm\relax}{\scriptsize\color{gray} [p. 23]}\ Baranowska, Grażyna, Alpes, Maybritt Jill and Kienzle, Isabel. 2025. “Making Pushback Facts Visible: A Review of Tools in Existing Case Law and the Procedural Framework of the European Court of Human Rights.” \textit{The International J. of Human Rights} 29(10): 1925–50. doi:10.1080/13642987.2025.2531790. 〔题名或条目译释：让推回事实可见：欧洲人权法院既有判例法及程序框架中工具的审视。〕\par\end{minipage}\hfill\begin{minipage}[t]{\dimexpr.5\textwidth-4mm\relax}{\scriptsize\color{gray} [p. 23]}\ Baranowska, Grażyna, Alpes, Maybritt Jill and Kienzle, Isabel. 2025. “Making Pushback Facts Visible: A Review of Tools in Existing Case Law and the Procedural Framework of the European Court of Human Rights.” \textit{The International J. of Human Rights} 29(10): 1925–50. doi:10.1080/13642987.2025.2531790.\par\end{minipage}\end{minipage}\par
\noindent\begin{minipage}{\textwidth}\begin{minipage}[t]{\dimexpr.5\textwidth-4mm\relax}{\scriptsize\color{gray} [p. 23]}\ Bergesio, Noemi and Bialasiewicz, Luiza. 2023. “The Entangled Geographies of Responsibility: Contested Policy Narratives of Migration Governance along the Balkan Route.” \textit{Environment and Planning D: Society and Space} 41(1): 33–55. doi:10.1177/02637758221137345. 〔题名或条目译释：责任交织的地理：巴尔干路线沿线移民治理中争议性的政策叙事。〕\par\end{minipage}\hfill\begin{minipage}[t]{\dimexpr.5\textwidth-4mm\relax}{\scriptsize\color{gray} [p. 23]}\ Bergesio, Noemi and Bialasiewicz, Luiza. 2023. “The Entangled Geographies of Responsibility: Contested Policy Narratives of Migration Governance along the Balkan Route.” \textit{Environment and Planning D: Society and Space} 41(1): 33–55. doi:10.1177/02637758221137345.\par\end{minipage}\end{minipage}\par
\noindent\begin{minipage}{\textwidth}\begin{minipage}[t]{\dimexpr.5\textwidth-4mm\relax}{\scriptsize\color{gray} [p. 23]}\ Billows, Sebastian, Buchter, Lisa and Pélisse, Jérôme. 2019. “Introduction: The Microfoundations of Legal Intermediation in Organizational Contexts.” In \textit{Legal Intermediation: A Processual Approach to Law and Economic Activity}, edited by Sarat, Austin, 1–9. Bingley: Emerald. 〔题名或条目译释：导论：组织情境中法律中介活动的微观基础；所收文集：法律中介活动——法律与经济活动的过程性研究。〕\par\end{minipage}\hfill\begin{minipage}[t]{\dimexpr.5\textwidth-4mm\relax}{\scriptsize\color{gray} [p. 23]}\ Billows, Sebastian, Buchter, Lisa and Pélisse, Jérôme. 2019. “Introduction: The Microfoundations of Legal Intermediation in Organizational Contexts.” In \textit{Legal Intermediation: A Processual Approach to Law and Economic Activity}, edited by Sarat, Austin, 1–9. Bingley: Emerald.\par\end{minipage}\end{minipage}\par
\noindent\begin{minipage}{\textwidth}\begin{minipage}[t]{\dimexpr.5\textwidth-4mm\relax}{\scriptsize\color{gray} [p. 23]}\ Bonizzoni, Paola and Odasso, Laura. 2024. “Legal-Administrative Intermediation in the Migration Field: An Introduction.” \textit{Etnografia e Ricerca Qualitativa} 17(2): 195–212. 〔题名或条目译释：移民领域的法律—行政中介活动：导论。〕\par\end{minipage}\hfill\begin{minipage}[t]{\dimexpr.5\textwidth-4mm\relax}{\scriptsize\color{gray} [p. 23]}\ Bonizzoni, Paola and Odasso, Laura. 2024. “Legal-Administrative Intermediation in the Migration Field: An Introduction.” \textit{Etnografia e Ricerca Qualitativa} 17(2): 195–212.\par\end{minipage}\end{minipage}\par
\noindent\begin{minipage}{\textwidth}\begin{minipage}[t]{\dimexpr.5\textwidth-4mm\relax}{\scriptsize\color{gray} [p. 23]}\ Bonjour, Saskia. 2025. “Gatekeepers or Agents of Change? The Political Agency of Lawyers in the Dutch Family Migration Regime.” \textit{J. of Ethnic and Migration Studies} 52(7): 1762–81. doi:10.1080/1369183X.2025.2504619. 〔题名或条目译释：守门人还是变革推动者？荷兰家庭移民制度中律师的政治能动性。〕\par\end{minipage}\hfill\begin{minipage}[t]{\dimexpr.5\textwidth-4mm\relax}{\scriptsize\color{gray} [p. 23]}\ Bonjour, Saskia. 2025. “Gatekeepers or Agents of Change? The Political Agency of Lawyers in the Dutch Family Migration Regime.” \textit{J. of Ethnic and Migration Studies} 52(7): 1762–81. doi:10.1080/1369183X.2025.2504619.\par\end{minipage}\end{minipage}\par
\noindent\begin{minipage}{\textwidth}\begin{minipage}[t]{\dimexpr.5\textwidth-4mm\relax}{\scriptsize\color{gray} [p. 23]}\ Cabot, Heath. 2012. “The Governance of Things: Documenting Limbo in the Greek Asylum Procedure.” \textit{PoLAR: Political and Legal Anthropology Rev.} 35(1): 11–29. doi:10.1111/j.1555-2934.2012.01177.x. 〔题名或条目译释：物的治理：记录希腊庇护程序中的悬置处境。〕\par\end{minipage}\hfill\begin{minipage}[t]{\dimexpr.5\textwidth-4mm\relax}{\scriptsize\color{gray} [p. 23]}\ Cabot, Heath. 2012. “The Governance of Things: Documenting Limbo in the Greek Asylum Procedure.” \textit{PoLAR: Political and Legal Anthropology Rev.} 35(1): 11–29. doi:10.1111/j.1555-2934.2012.01177.x.\par\end{minipage}\end{minipage}\par
\noindent\begin{minipage}{\textwidth}\begin{minipage}[t]{\dimexpr.5\textwidth-4mm\relax}{\scriptsize\color{gray} [p. 23]}\ Campbell, John R. 2020. “The role of lawyers, judges, country experts and officials in British asylum and immigration law.” \textit{International J. of Law in Context} 16(1): 1–16. doi:10.1017/S1744552320000038. 〔题名或条目译释：律师、法官、国情专家与官员在英国庇护及移民法中的作用。〕\par\end{minipage}\hfill\begin{minipage}[t]{\dimexpr.5\textwidth-4mm\relax}{\scriptsize\color{gray} [p. 23]}\ Campbell, John R. 2020. “The role of lawyers, judges, country experts and officials in British asylum and immigration law.” \textit{International J. of Law in Context} 16(1): 1–16. doi:10.1017/S1744552320000038.\par\end{minipage}\end{minipage}\par
\noindent\begin{minipage}{\textwidth}\begin{minipage}[t]{\dimexpr.5\textwidth-4mm\relax}{\scriptsize\color{gray} [p. 23]}\ Cantat, Céline. 2020. “The Rise and Fall of Migration Solidarity in Belgrade.” \textit{Movements: J. for Critical Migration and Border Regime Studies} 5(1): 97–124. doi:10.64081/mvmnts-5.1-2389. 〔题名或条目译释：贝尔格莱德移民声援行动的兴衰。〕\par\end{minipage}\hfill\begin{minipage}[t]{\dimexpr.5\textwidth-4mm\relax}{\scriptsize\color{gray} [p. 23]}\ Cantat, Céline. 2020. “The Rise and Fall of Migration Solidarity in Belgrade.” \textit{Movements: J. for Critical Migration and Border Regime Studies} 5(1): 97–124. doi:10.64081/mvmnts-5.1-2389.\par\end{minipage}\end{minipage}\par
\noindent\begin{minipage}{\textwidth}\begin{minipage}[t]{\dimexpr.5\textwidth-4mm\relax}{\scriptsize\color{gray} [p. 23]}\ Cantor, David, Tan, Nikolas Feith, Gkliati, Mariana, Mavropoulou, Elizabeth, Allinson, Kathryn, Chakrabarty, Sreetapa, Grundler, Maja, et al. 2022. “Externalization, Access to Territorial Asylum, and International Law.” \textit{International J. of Refugee Law} 34(1): 120–56. doi:10.1093/ijrl/eeac023. 〔题名或条目译释：外部化、领土庇护准入与国际法。〕\par\end{minipage}\hfill\begin{minipage}[t]{\dimexpr.5\textwidth-4mm\relax}{\scriptsize\color{gray} [p. 23]}\ Cantor, David, Tan, Nikolas Feith, Gkliati, Mariana, Mavropoulou, Elizabeth, Allinson, Kathryn, Chakrabarty, Sreetapa, Grundler, Maja, et al. 2022. “Externalization, Access to Territorial Asylum, and International Law.” \textit{International J. of Refugee Law} 34(1): 120–56. doi:10.1093/ijrl/eeac023.\par\end{minipage}\end{minipage}\par
\noindent\begin{minipage}{\textwidth}\begin{minipage}[t]{\dimexpr.5\textwidth-4mm\relax}{\scriptsize\color{gray} [p. 23]}\ Casas-cortes, Maribel, Cobarrubias, Sebastian and Pickles, John. 2016. “‘Good Neighbours Make Good Fences’: Seahorse Operations, Border Externalization and Extra-Territoriality.” \textit{European Urban and Regional Studies} 23(3): 231–51. doi:10.1177/0969776414541136. 〔题名或条目译释：‘好邻居筑好篱笆’：海马行动、边境外部化与域外性。〕\par\end{minipage}\hfill\begin{minipage}[t]{\dimexpr.5\textwidth-4mm\relax}{\scriptsize\color{gray} [p. 23]}\ Casas-cortes, Maribel, Cobarrubias, Sebastian and Pickles, John. 2016. “‘Good Neighbours Make Good Fences’: Seahorse Operations, Border Externalization and Extra-Territoriality.” \textit{European Urban and Regional Studies} 23(3): 231–51. doi:10.1177/0969776414541136.\par\end{minipage}\end{minipage}\par
\noindent\begin{minipage}{\textwidth}\begin{minipage}[t]{\dimexpr.5\textwidth-4mm\relax}{\scriptsize\color{gray} [p. 23]}\ Conant, Lisa. 2016. “Who Files Suit? Legal Mobilization and Torture Violations in Europe,” \textit{Law \& Policy} 38(4): 280–303. doi:10.1111/lapo.12055. 〔题名或条目译释：谁提起诉讼？欧洲的法律动员与酷刑侵权。〕\par\end{minipage}\hfill\begin{minipage}[t]{\dimexpr.5\textwidth-4mm\relax}{\scriptsize\color{gray} [p. 23]}\ Conant, Lisa. 2016. “Who Files Suit? Legal Mobilization and Torture Violations in Europe,” \textit{Law \& Policy} 38(4): 280–303. doi:10.1111/lapo.12055.\par\end{minipage}\end{minipage}\par
\noindent\begin{minipage}{\textwidth}\begin{minipage}[t]{\dimexpr.5\textwidth-4mm\relax}{\scriptsize\color{gray} [p. 23]}\ Costello, Cathryn and Mann, Itamar. 2020. “Border Justice: Migration and Accountability for Human Rights Violations.” \textit{German Law J.} 21(3): 311–34. doi:10.1017/glj.2020.27. 〔题名或条目译释：边境正义：移民与人权侵害问责。〕\par\end{minipage}\hfill\begin{minipage}[t]{\dimexpr.5\textwidth-4mm\relax}{\scriptsize\color{gray} [p. 23]}\ Costello, Cathryn and Mann, Itamar. 2020. “Border Justice: Migration and Accountability for Human Rights Violations.” \textit{German Law J.} 21(3): 311–34. doi:10.1017/glj.2020.27.\par\end{minipage}\end{minipage}\par
\noindent\begin{minipage}{\textwidth}\begin{minipage}[t]{\dimexpr.5\textwidth-4mm\relax}{\scriptsize\color{gray} [p. 23]}\ Coutin, Susan Bibler. 2025. \textit{On the Record: Papers, Immigration, and Legal Advocacy}. Oakland: University of California Press. 〔题名或条目译释：记入案卷：文书、移民与法律倡导。〕\par\end{minipage}\hfill\begin{minipage}[t]{\dimexpr.5\textwidth-4mm\relax}{\scriptsize\color{gray} [p. 23]}\ Coutin, Susan Bibler. 2025. \textit{On the Record: Papers, Immigration, and Legal Advocacy}. Oakland: University of California Press.\par\end{minipage}\end{minipage}\par
\noindent\begin{minipage}{\textwidth}\begin{minipage}[t]{\dimexpr.5\textwidth-4mm\relax}{\scriptsize\color{gray} [p. 23]}\ Coutin, Susan Bibler and Fortin, Véronique. 2023. “The Craft of Translation: Documentary Practices within Immigration Advocacy in the United States.” \textit{PoLAR: Political and Legal Anthropology Rev.} 46(1): 24–38. doi:10.1111/plar.12531. 〔题名或条目译释：转译的技艺：美国移民权益倡导中的文书实践。〕\par\end{minipage}\hfill\begin{minipage}[t]{\dimexpr.5\textwidth-4mm\relax}{\scriptsize\color{gray} [p. 23]}\ Coutin, Susan Bibler and Fortin, Véronique. 2023. “The Craft of Translation: Documentary Practices within Immigration Advocacy in the United States.” \textit{PoLAR: Political and Legal Anthropology Rev.} 46(1): 24–38. doi:10.1111/plar.12531.\par\end{minipage}\end{minipage}\par
\noindent\begin{minipage}{\textwidth}\begin{minipage}[t]{\dimexpr.5\textwidth-4mm\relax}{\scriptsize\color{gray} [p. 23]}\ Crooke, Catherine 2024. “Frustration and Fidelity: How Public Interest Lawyers Navigate Procedure in the Direct Representation of Asylum Seekers.” \textit{Law \& Society Rev.} 58(2): 270–93. doi:10.1017/lsr.2024.20. 〔题名或条目译释：挫折与忠实：公益律师如何在直接代理寻求庇护者时应对程序。〕\par\end{minipage}\hfill\begin{minipage}[t]{\dimexpr.5\textwidth-4mm\relax}{\scriptsize\color{gray} [p. 23]}\ Crooke, Catherine 2024. “Frustration and Fidelity: How Public Interest Lawyers Navigate Procedure in the Direct Representation of Asylum Seekers.” \textit{Law \& Society Rev.} 58(2): 270–93. doi:10.1017/lsr.2024.20.\par\end{minipage}\end{minipage}\par
\noindent\begin{minipage}{\textwidth}\begin{minipage}[t]{\dimexpr.5\textwidth-4mm\relax}{\scriptsize\color{gray} [p. 23]}\ Cummings, Scott L. 2018. “The Social Movement Turn in Law.” \textit{Law \& Social Inquiry} 43(2): 360–416. doi:10.1111/lsi.12308. 〔题名或条目译释：法律中的社会运动转向。〕\par\end{minipage}\hfill\begin{minipage}[t]{\dimexpr.5\textwidth-4mm\relax}{\scriptsize\color{gray} [p. 23]}\ Cummings, Scott L. 2018. “The Social Movement Turn in Law.” \textit{Law \& Social Inquiry} 43(2): 360–416. doi:10.1111/lsi.12308.\par\end{minipage}\end{minipage}\par
\noindent\begin{minipage}{\textwidth}\begin{minipage}[t]{\dimexpr.5\textwidth-4mm\relax}{\scriptsize\color{gray} [p. 23]}\ Cummings, Scott L. 2020. “Movement Lawyering.” \textit{Indiana J. of Global Legal Studies} 27(1): 87–130. doi:10.2979/indjglolegstu.27.1.0087. 〔题名或条目译释：社会运动律师工作。〕\par\end{minipage}\hfill\begin{minipage}[t]{\dimexpr.5\textwidth-4mm\relax}{\scriptsize\color{gray} [p. 23]}\ Cummings, Scott L. 2020. “Movement Lawyering.” \textit{Indiana J. of Global Legal Studies} 27(1): 87–130. doi:10.2979/indjglolegstu.27.1.0087.\par\end{minipage}\end{minipage}\par
\noindent\begin{minipage}{\textwidth}\begin{minipage}[t]{\dimexpr.5\textwidth-4mm\relax}{\scriptsize\color{gray} [p. 23]}\ Cunningham, Clark D. 1992. “Lawyer as Translator, Representation as Text: Towards an Ethnography of Legal Discourse.” \textit{Cornell Law Rev.} 77(6): 1298–387. 〔题名或条目译释：作为转译者的律师，作为文本的代理：迈向法律话语民族志。〕\par\end{minipage}\hfill\begin{minipage}[t]{\dimexpr.5\textwidth-4mm\relax}{\scriptsize\color{gray} [p. 23]}\ Cunningham, Clark D. 1992. “Lawyer as Translator, Representation as Text: Towards an Ethnography of Legal Discourse.” \textit{Cornell Law Rev.} 77(6): 1298–387.\par\end{minipage}\end{minipage}\par
\noindent\begin{minipage}{\textwidth}\begin{minipage}[t]{\dimexpr.5\textwidth-4mm\relax}{\scriptsize\color{gray} [p. 23]}\ Cusumano, Eugenio and Villa, Matteo. 2021. “From ‘Angels’ to ‘Vice Smugglers’: The Criminalization of Sea Rescue NGOs in Italy.” \textit{European J. on Criminal Policy and Research} 27(1): 23–40. doi:10.1007/s10610-020-09464-1. 〔题名或条目译释：从‘天使’到‘走私帮凶’：意大利海上救援非政府组织的刑事化。〕\par\end{minipage}\hfill\begin{minipage}[t]{\dimexpr.5\textwidth-4mm\relax}{\scriptsize\color{gray} [p. 23]}\ Cusumano, Eugenio and Villa, Matteo. 2021. “From ‘Angels’ to ‘Vice Smugglers’: The Criminalization of Sea Rescue NGOs in Italy.” \textit{European J. on Criminal Policy and Research} 27(1): 23–40. doi:10.1007/s10610-020-09464-1.\par\end{minipage}\end{minipage}\par
\noindent\begin{minipage}{\textwidth}\begin{minipage}[t]{\dimexpr.5\textwidth-4mm\relax}{\scriptsize\color{gray} [p. 23]}\ Davies, Thom, Isakjee, Arshad and Obradović-wochnik, Jelena. 2023. “Epistemic Borderwork: Violent Pushbacks, Refugees, and the Politics of Knowledge at the EU Border.” \textit{Annals of the American Association of Geographers} 113(1): 169–88. doi:10.1080/24694452.2022.2077167. 〔题名或条目译释：知识层面的边境工作：欧盟边境的暴力推回、难民与知识政治。〕\par\end{minipage}\hfill\begin{minipage}[t]{\dimexpr.5\textwidth-4mm\relax}{\scriptsize\color{gray} [p. 23]}\ Davies, Thom, Isakjee, Arshad and Obradović-wochnik, Jelena. 2023. “Epistemic Borderwork: Violent Pushbacks, Refugees, and the Politics of Knowledge at the EU Border.” \textit{Annals of the American Association of Geographers} 113(1): 169–88. doi:10.1080/24694452.2022.2077167.\par\end{minipage}\end{minipage}\par
\noindent\begin{minipage}{\textwidth}\begin{minipage}[t]{\dimexpr.5\textwidth-4mm\relax}{\scriptsize\color{gray} [p. 24]}\ El-Shaarawi, Nadine and Razsa, Maple. 2019. “Movements upon Movements: Refugee and Activist Struggles to Open the Balkan Route to Europe.” \textit{History and Anthropology} 30(1): 91–112. doi:10.1080/02757206.2018.1530668. 〔题名或条目译释：运动之上的运动：难民与行动者为打通通往欧洲的巴尔干路线而进行的斗争。〕\par\end{minipage}\hfill\begin{minipage}[t]{\dimexpr.5\textwidth-4mm\relax}{\scriptsize\color{gray} [p. 24]}\ El-Shaarawi, Nadine and Razsa, Maple. 2019. “Movements upon Movements: Refugee and Activist Struggles to Open the Balkan Route to Europe.” \textit{History and Anthropology} 30(1): 91–112. doi:10.1080/02757206.2018.1530668.\par\end{minipage}\end{minipage}\par
\noindent\begin{minipage}{\textwidth}\begin{minipage}[t]{\dimexpr.5\textwidth-4mm\relax}{\scriptsize\color{gray} [p. 24]}\ Epp, Charles R. 1998. \textit{The Rights Revolution: Lawyers, Activists, and Supreme Courts in Comparative Perspective}. Chicago: University of Chicago Press. 〔题名或条目译释：权利革命：比较视野中的律师、行动者与最高法院。〕\par\end{minipage}\hfill\begin{minipage}[t]{\dimexpr.5\textwidth-4mm\relax}{\scriptsize\color{gray} [p. 24]}\ Epp, Charles R. 1998. \textit{The Rights Revolution: Lawyers, Activists, and Supreme Courts in Comparative Perspective}. Chicago: University of Chicago Press.\par\end{minipage}\end{minipage}\par
\noindent\begin{minipage}{\textwidth}\begin{minipage}[t]{\dimexpr.5\textwidth-4mm\relax}{\scriptsize\color{gray} [p. 24]}\ Epp, Charles R. 2011. “The Support Structure as a Necessary Condition for Sustained Judicial Attention to Rights: A Response.” \textit{J. of Politics} 73(2): 406–09.doi: 10.1017/S0022381611000089. 〔题名或条目译释：支持结构作为司法持续关注权利的必要条件：一项回应。〕\par\end{minipage}\hfill\begin{minipage}[t]{\dimexpr.5\textwidth-4mm\relax}{\scriptsize\color{gray} [p. 24]}\ Epp, Charles R. 2011. “The Support Structure as a Necessary Condition for Sustained Judicial Attention to Rights: A Response.” \textit{J. of Politics} 73(2): 406–09.doi: 10.1017/S0022381611000089.\par\end{minipage}\end{minipage}\par
\noindent\begin{minipage}{\textwidth}\begin{minipage}[t]{\dimexpr.5\textwidth-4mm\relax}{\scriptsize\color{gray} [p. 24]}\ Facchini, Duccio. 2020. “Rotta balcanica: Nel 2020 record di respingimenti dall’Italia verso la Slovenia.” Altreconomia, December 28 . 〔题名或条目译释：巴尔干路线：2020年从意大利向斯洛文尼亚的推回创纪录。〕\par\end{minipage}\hfill\begin{minipage}[t]{\dimexpr.5\textwidth-4mm\relax}{\scriptsize\color{gray} [p. 24]}\ Facchini, Duccio. 2020. “Rotta balcanica: Nel 2020 record di respingimenti dall’Italia verso la Slovenia.” Altreconomia, December 28 .\par\end{minipage}\end{minipage}\par
\noindent\begin{minipage}{\textwidth}\begin{minipage}[t]{\dimexpr.5\textwidth-4mm\relax}{\scriptsize\color{gray} [p. 24]}\ Fassin, Didier and D’halluin, Estelle. 2005. “The Truth from the Body: Medical Certificates as Ultimate Evidence for Asylum Seekers.” \textit{American Anthropologist} 107(4): 597–608. doi:10.1525/aa.2005.107.4.597. 〔题名或条目译释：来自身体的真相：医疗证明作为寻求庇护者的终极证据。〕\par\end{minipage}\hfill\begin{minipage}[t]{\dimexpr.5\textwidth-4mm\relax}{\scriptsize\color{gray} [p. 24]}\ Fassin, Didier and D’halluin, Estelle. 2005. “The Truth from the Body: Medical Certificates as Ultimate Evidence for Asylum Seekers.” \textit{American Anthropologist} 107(4): 597–608. doi:10.1525/aa.2005.107.4.597.\par\end{minipage}\end{minipage}\par
\noindent\begin{minipage}{\textwidth}\begin{minipage}[t]{\dimexpr.5\textwidth-4mm\relax}{\scriptsize\color{gray} [p. 24]}\ Felstiner, William L. F., Abel, Richard L. and Sarat, Austin. 1981. “The Emergence and Transformation of Disputes: Naming, Blaming, Claiming.” \textit{Law \& Society Rev.} 15: 631–54. doi:10.2307/3053505. 〔题名或条目译释：纠纷的产生与转变：命名、归责与索求。〕\par\end{minipage}\hfill\begin{minipage}[t]{\dimexpr.5\textwidth-4mm\relax}{\scriptsize\color{gray} [p. 24]}\ Felstiner, William L. F., Abel, Richard L. and Sarat, Austin. 1981. “The Emergence and Transformation of Disputes: Naming, Blaming, Claiming.” \textit{Law \& Society Rev.} 15: 631–54. doi:10.2307/3053505.\par\end{minipage}\end{minipage}\par
\noindent\begin{minipage}{\textwidth}\begin{minipage}[t]{\dimexpr.5\textwidth-4mm\relax}{\scriptsize\color{gray} [p. 24]}\ Fortarezza, Francesca. 2026. “Testing the Liberal Borders of the EU: (De)Constructing the Right of Asylum through Informality.” \textit{Political Geography} 126: 103499. doi:10.1016/j.polgeo.2026.103499. 〔题名或条目译释：检验欧盟自由主义边界：通过非正式性建构与解构庇护权。〕\par\end{minipage}\hfill\begin{minipage}[t]{\dimexpr.5\textwidth-4mm\relax}{\scriptsize\color{gray} [p. 24]}\ Fortarezza, Francesca. 2026. “Testing the Liberal Borders of the EU: (De)Constructing the Right of Asylum through Informality.” \textit{Political Geography} 126: 103499. doi:10.1016/j.polgeo.2026.103499.\par\end{minipage}\end{minipage}\par
\noindent\begin{minipage}{\textwidth}\begin{minipage}[t]{\dimexpr.5\textwidth-4mm\relax}{\scriptsize\color{gray} [p. 24]}\ Gahleitner-gertz, Lukas. 2025. “Country Report: Austria. Update on 2024.” Asylum Information Database (AIDA), European Council on Refugees and Exiles, [\url{https://asylumineurope.org/wp-content/uploads/2025/07/AIDA_AT_2024-Update.pdf]} (accessed October 21 , 2025). 〔题名或条目译释：国别报告：奥地利——2024年度更新。〕\par\end{minipage}\hfill\begin{minipage}[t]{\dimexpr.5\textwidth-4mm\relax}{\scriptsize\color{gray} [p. 24]}\ Gahleitner-gertz, Lukas. 2025. “Country Report: Austria. Update on 2024.” Asylum Information Database (AIDA), European Council on Refugees and Exiles, [\url{https://asylumineurope.org/wp-content/uploads/2025/07/AIDA_AT_2024-Update.pdf]} (accessed October 21 , 2025).\par\end{minipage}\end{minipage}\par
\noindent\begin{minipage}{\textwidth}\begin{minipage}[t]{\dimexpr.5\textwidth-4mm\relax}{\scriptsize\color{gray} [p. 24]}\ Gammeltoft-Hansen, Thomas. 2011. \textit{Access to Asylum: International Refugee Law and the Globalisation of Migration Control}. Cambridge: Cambridge University Press. DOI: 10.1017/CBO9780511763403 〔题名或条目译释：庇护准入：国际难民法与移民管控的全球化。〕\par\end{minipage}\hfill\begin{minipage}[t]{\dimexpr.5\textwidth-4mm\relax}{\scriptsize\color{gray} [p. 24]}\ Gammeltoft-Hansen, Thomas. 2011. \textit{Access to Asylum: International Refugee Law and the Globalisation of Migration Control}. Cambridge: Cambridge University Press. DOI: 10.1017/CBO9780511763403\par\end{minipage}\end{minipage}\par
\noindent\begin{minipage}{\textwidth}\begin{minipage}[t]{\dimexpr.5\textwidth-4mm\relax}{\scriptsize\color{gray} [p. 24]}\ Gammeltoft-hansen, Thomas and Tan, Nikolas Feith. 2021. “Extraterritorial Migration Control and Deterrence.” In \textit{The Oxford Handbook of International Refugee Law}, edited by Costello, Cathryn, Foster, Michelle and McAdam, Jane, 502–17. Oxford: Oxford University Press. DOI: 10.1093/law/9780198848639.003.0028 〔题名或条目译释：域外移民管控与威慑；所收文集：牛津国际难民法手册。〕\par\end{minipage}\hfill\begin{minipage}[t]{\dimexpr.5\textwidth-4mm\relax}{\scriptsize\color{gray} [p. 24]}\ Gammeltoft-hansen, Thomas and Tan, Nikolas Feith. 2021. “Extraterritorial Migration Control and Deterrence.” In \textit{The Oxford Handbook of International Refugee Law}, edited by Costello, Cathryn, Foster, Michelle and McAdam, Jane, 502–17. Oxford: Oxford University Press. DOI: 10.1093/law/9780198848639.003.0028\par\end{minipage}\end{minipage}\par
\noindent\begin{minipage}{\textwidth}\begin{minipage}[t]{\dimexpr.5\textwidth-4mm\relax}{\scriptsize\color{gray} [p. 24]}\ Giuffré, Mariagiulia. 2020. \textit{The Readmission of Asylum Seekers under International Law}. Oxford: Hart Publishing. DOI: 10.5040/9781509902514 〔题名或条目译释：国际法下寻求庇护者的重新接收。〕\par\end{minipage}\hfill\begin{minipage}[t]{\dimexpr.5\textwidth-4mm\relax}{\scriptsize\color{gray} [p. 24]}\ Giuffré, Mariagiulia. 2020. \textit{The Readmission of Asylum Seekers under International Law}. Oxford: Hart Publishing. DOI: 10.5040/9781509902514\par\end{minipage}\end{minipage}\par
\noindent\begin{minipage}{\textwidth}\begin{minipage}[t]{\dimexpr.5\textwidth-4mm\relax}{\scriptsize\color{gray} [p. 24]}\ Hameršak, Marijana. 2021. “Between Supporting and Reporting: Grassroots Textual Responses to the Pushbacks at the Fringes of EU.” \textit{Intersections. East European J. of Society and Politics} 7(2): 49–72. 〔题名或条目译释：支持与报告之间：欧盟边缘针对推回的基层文本回应。〕\par\end{minipage}\hfill\begin{minipage}[t]{\dimexpr.5\textwidth-4mm\relax}{\scriptsize\color{gray} [p. 24]}\ Hameršak, Marijana. 2021. “Between Supporting and Reporting: Grassroots Textual Responses to the Pushbacks at the Fringes of EU.” \textit{Intersections. East European J. of Society and Politics} 7(2): 49–72.\par\end{minipage}\end{minipage}\par
\noindent\begin{minipage}{\textwidth}\begin{minipage}[t]{\dimexpr.5\textwidth-4mm\relax}{\scriptsize\color{gray} [p. 24]}\ Hameršak, Marijana, Hess, Sabine, Speer, Marc and Stojić Mitrović, Marta. 2020. “The Forging of the Balkan Route: Contextualizing the Border Regime in the EU Periphery.” \textit{Movements: J. for Critical Migration and Border Regime Studies} 5(1): 9–29. doi:10.64081/mvmnts-5.1-2385. 〔题名或条目译释：巴尔干路线的形成：将欧盟外围边境制度置于背景之中。〕\par\end{minipage}\hfill\begin{minipage}[t]{\dimexpr.5\textwidth-4mm\relax}{\scriptsize\color{gray} [p. 24]}\ Hameršak, Marijana, Hess, Sabine, Speer, Marc and Stojić Mitrović, Marta. 2020. “The Forging of the Balkan Route: Contextualizing the Border Regime in the EU Periphery.” \textit{Movements: J. for Critical Migration and Border Regime Studies} 5(1): 9–29. doi:10.64081/mvmnts-5.1-2385.\par\end{minipage}\end{minipage}\par
\noindent\begin{minipage}{\textwidth}\begin{minipage}[t]{\dimexpr.5\textwidth-4mm\relax}{\scriptsize\color{gray} [p. 24]}\ Italian Ministry of the Interior. 2020. “Question time alla Camera con il ministro Lamorgese”, [\url{https://www.interno.gov.it/it/notizie/question-time-alla-camera-ministro-lamorgese-0]} (accessed October 20 , 2025). 〔题名或条目译释：众议院与拉莫尔杰塞部长的质询时间。〕\par\end{minipage}\hfill\begin{minipage}[t]{\dimexpr.5\textwidth-4mm\relax}{\scriptsize\color{gray} [p. 24]}\ Italian Ministry of the Interior. 2020. “Question time alla Camera con il ministro Lamorgese”, [\url{https://www.interno.gov.it/it/notizie/question-time-alla-camera-ministro-lamorgese-0]} (accessed October 20 , 2025).\par\end{minipage}\end{minipage}\par
\noindent\begin{minipage}{\textwidth}\begin{minipage}[t]{\dimexpr.5\textwidth-4mm\relax}{\scriptsize\color{gray} [p. 24]}\ Italian Ministry of the Interior. 2021. Reply to the parliament inquiry. June 2020, \url{https://www.asgi.it/wp-content/uploads/2020/08/Risposta-interpellanza-rotta-balcanica.pdf} (accessed August 15 , 2025) 〔题名或条目译释：对议会质询的答复。〕\par\end{minipage}\hfill\begin{minipage}[t]{\dimexpr.5\textwidth-4mm\relax}{\scriptsize\color{gray} [p. 24]}\ Italian Ministry of the Interior. 2021. Reply to the parliament inquiry. June 2020, \url{https://www.asgi.it/wp-content/uploads/2020/08/Risposta-interpellanza-rotta-balcanica.pdf} (accessed August 15 , 2025)\par\end{minipage}\end{minipage}\par
\noindent\begin{minipage}{\textwidth}\begin{minipage}[t]{\dimexpr.5\textwidth-4mm\relax}{\scriptsize\color{gray} [p. 24]}\ Jacobs, Marie and Maryns, Katrijn. 2022. “Managing narratives, managing identities: Language and credibility in legal consultations with asylum seekers.” \textit{Language in Society} 51(3): 375–402. doi:10.1017/S0047404521000117. 〔题名或条目译释：管理叙事，管理身份：寻求庇护者法律咨询中的语言与可信度。〕\par\end{minipage}\hfill\begin{minipage}[t]{\dimexpr.5\textwidth-4mm\relax}{\scriptsize\color{gray} [p. 24]}\ Jacobs, Marie and Maryns, Katrijn. 2022. “Managing narratives, managing identities: Language and credibility in legal consultations with asylum seekers.” \textit{Language in Society} 51(3): 375–402. doi:10.1017/S0047404521000117.\par\end{minipage}\end{minipage}\par
\noindent\begin{minipage}{\textwidth}\begin{minipage}[t]{\dimexpr.5\textwidth-4mm\relax}{\scriptsize\color{gray} [p. 24]}\ Johan, Lindquist, Xiang, Biao and Yeoh, Brenda S. A.. 2012. “Opening the Black Box of Migration: Brokers, the Organization of Transnational Mobility and the Changing Political Economy in Asia.” \textit{Pacific Affairs} 85(1): 7–19. doi:10.5509/20128517. 〔题名或条目译释：打开移民黑箱：经纪者、跨国流动的组织与亚洲变动中的政治经济。〕\par\end{minipage}\hfill\begin{minipage}[t]{\dimexpr.5\textwidth-4mm\relax}{\scriptsize\color{gray} [p. 24]}\ Johan, Lindquist, Xiang, Biao and Yeoh, Brenda S. A.. 2012. “Opening the Black Box of Migration: Brokers, the Organization of Transnational Mobility and the Changing Political Economy in Asia.” \textit{Pacific Affairs} 85(1): 7–19. doi:10.5509/20128517.\par\end{minipage}\end{minipage}\par
\noindent\begin{minipage}{\textwidth}\begin{minipage}[t]{\dimexpr.5\textwidth-4mm\relax}{\scriptsize\color{gray} [p. 24]}\ Jovanović, Teodora. 2020. “Transformations of Humanitarian Aid and Response Modes to Migration Movements: A Case Study of the Miksalište Center in Belgrade.” \textit{Movements: J. for Critical Migration and Border Regime Studies} 5(1): 125–47. doi:10.64081/mvmnts-5.1-2390. 〔题名或条目译释：人道援助与迁徙运动应对方式的转变：贝尔格莱德Miksalište中心案例研究。〕\par\end{minipage}\hfill\begin{minipage}[t]{\dimexpr.5\textwidth-4mm\relax}{\scriptsize\color{gray} [p. 24]}\ Jovanović, Teodora. 2020. “Transformations of Humanitarian Aid and Response Modes to Migration Movements: A Case Study of the Miksalište Center in Belgrade.” \textit{Movements: J. for Critical Migration and Border Regime Studies} 5(1): 125–47. doi:10.64081/mvmnts-5.1-2390.\par\end{minipage}\end{minipage}\par
\noindent\begin{minipage}{\textwidth}\begin{minipage}[t]{\dimexpr.5\textwidth-4mm\relax}{\scriptsize\color{gray} [p. 24]}\ Kawar, Leila. 2015. \textit{Contesting Immigration Policy in Court: Legal Activism and Its Radiating Effects in the United States and France}. Cambridge: Cambridge University Press. DOI: 10.1017/CBO9781107741980 〔题名或条目译释：在法院挑战移民政策：美国与法国的法律行动及其辐射效应。〕\par\end{minipage}\hfill\begin{minipage}[t]{\dimexpr.5\textwidth-4mm\relax}{\scriptsize\color{gray} [p. 24]}\ Kawar, Leila. 2015. \textit{Contesting Immigration Policy in Court: Legal Activism and Its Radiating Effects in the United States and France}. Cambridge: Cambridge University Press. DOI: 10.1017/CBO9781107741980\par\end{minipage}\end{minipage}\par
\noindent\begin{minipage}{\textwidth}\begin{minipage}[t]{\dimexpr.5\textwidth-4mm\relax}{\scriptsize\color{gray} [p. 24]}\ Kawar, Leila. 2021. “Contesting Migration Governance through Legal Mobilization.” In \textit{Handbook on the Governance and Politics of Migration}, edited by Carmel, Emma, Lenner, Katharina and Paul, Regine, 380–90. Cheltenham: Edward Elgar. 〔题名或条目译释：通过法律动员挑战移民治理；所收文集：移民治理与政治手册。〕\par\end{minipage}\hfill\begin{minipage}[t]{\dimexpr.5\textwidth-4mm\relax}{\scriptsize\color{gray} [p. 24]}\ Kawar, Leila. 2021. “Contesting Migration Governance through Legal Mobilization.” In \textit{Handbook on the Governance and Politics of Migration}, edited by Carmel, Emma, Lenner, Katharina and Paul, Regine, 380–90. Cheltenham: Edward Elgar.\par\end{minipage}\end{minipage}\par
\noindent\begin{minipage}{\textwidth}\begin{minipage}[t]{\dimexpr.5\textwidth-4mm\relax}{\scriptsize\color{gray} [p. 24]}\ Kogovšek Šalamon, Neža, ed. 2020. \textit{Causes and Consequences of Migrant Criminalization}. Cham: Springer. DOI: 10.1007/978-3-030-43732-9 〔题名或条目译释：移民刑事化的原因与后果。〕\par\end{minipage}\hfill\begin{minipage}[t]{\dimexpr.5\textwidth-4mm\relax}{\scriptsize\color{gray} [p. 24]}\ Kogovšek Šalamon, Neža, ed. 2020. \textit{Causes and Consequences of Migrant Criminalization}. Cham: Springer. DOI: 10.1007/978-3-030-43732-9\par\end{minipage}\end{minipage}\par
\noindent\begin{minipage}{\textwidth}\begin{minipage}[t]{\dimexpr.5\textwidth-4mm\relax}{\scriptsize\color{gray} [p. 24]}\ Könönen, Jukka. 2018. “Border Struggles within the State: Administrative Bordering of Non-Citizens in Finland.” \textit{Nordic J. of Migration Research} 8(3): 143–50. doi:10.2478/njmr-2018-0018. 〔题名或条目译释：国家内部的边界斗争：芬兰以行政手段对非公民划界。〕\par\end{minipage}\hfill\begin{minipage}[t]{\dimexpr.5\textwidth-4mm\relax}{\scriptsize\color{gray} [p. 24]}\ Könönen, Jukka. 2018. “Border Struggles within the State: Administrative Bordering of Non-Citizens in Finland.” \textit{Nordic J. of Migration Research} 8(3): 143–50. doi:10.2478/njmr-2018-0018.\par\end{minipage}\end{minipage}\par
\noindent\begin{minipage}{\textwidth}\begin{minipage}[t]{\dimexpr.5\textwidth-4mm\relax}{\scriptsize\color{gray} [p. 24]}\ Lacchei, Alice and Moraru, Madalina. 2026. “Asylum Access Adjudication: A Multi-Level Framework for Socio-Legal Research.” \textit{J. of Refugee Studies} 39 (1) : 41–61. doi:10.1093/jrs/feaf069. 〔题名或条目译释：庇护准入裁判：法律与社会研究的多层次框架。〕\par\end{minipage}\hfill\begin{minipage}[t]{\dimexpr.5\textwidth-4mm\relax}{\scriptsize\color{gray} [p. 24]}\ Lacchei, Alice and Moraru, Madalina. 2026. “Asylum Access Adjudication: A Multi-Level Framework for Socio-Legal Research.” \textit{J. of Refugee Studies} 39 (1) : 41–61. doi:10.1093/jrs/feaf069.\par\end{minipage}\end{minipage}\par
\noindent\begin{minipage}{\textwidth}\begin{minipage}[t]{\dimexpr.5\textwidth-4mm\relax}{\scriptsize\color{gray} [p. 25]}\ Lakhani, Sarah Morando. 2014. “From Problems of Living to Problems of Law: The Legal Translation and Documentation of Immigrant Abuse and Helpfulness.” \textit{Law \& Social Inquiry} 39(3): 643–65. doi:10.1111/lsi.12082. 〔题名或条目译释：从生活问题到法律问题：移民受虐及提供协助的法律转译与文书化。〕\par\end{minipage}\hfill\begin{minipage}[t]{\dimexpr.5\textwidth-4mm\relax}{\scriptsize\color{gray} [p. 25]}\ Lakhani, Sarah Morando. 2014. “From Problems of Living to Problems of Law: The Legal Translation and Documentation of Immigrant Abuse and Helpfulness.” \textit{Law \& Social Inquiry} 39(3): 643–65. doi:10.1111/lsi.12082.\par\end{minipage}\end{minipage}\par
\noindent\begin{minipage}{\textwidth}\begin{minipage}[t]{\dimexpr.5\textwidth-4mm\relax}{\scriptsize\color{gray} [p. 25]}\ Latour, Bruno. 2010. \textit{The Making of Law: An Ethnography of the Conseil D’état}. Translated by Marina Brilman and Alain Pottage. Cambridge: Polity. 〔题名或条目译释：法律的制作：法国最高行政法院民族志。〕\par\end{minipage}\hfill\begin{minipage}[t]{\dimexpr.5\textwidth-4mm\relax}{\scriptsize\color{gray} [p. 25]}\ Latour, Bruno. 2010. \textit{The Making of Law: An Ethnography of the Conseil D’état}. Translated by Marina Brilman and Alain Pottage. Cambridge: Polity.\par\end{minipage}\end{minipage}\par
\noindent\begin{minipage}{\textwidth}\begin{minipage}[t]{\dimexpr.5\textwidth-4mm\relax}{\scriptsize\color{gray} [p. 25]}\ Lavenex, Sandra. 2022. The cat and mouse game of refugee externalisation policies. Between law and politics \textit{Refugee Externalisation Policies. Responsibility, Legitimacy and Accountability} Azadeh Dastyari, Amy Nethery and Asher Hirsch 27–44. London: Routledge. DOI: 10.4324/9781003167273-3 〔题名或条目译释：难民外部化政策的猫鼠游戏：法律与政治之间；所收文集：难民外部化政策——责任、正当性与问责。〕\par\end{minipage}\hfill\begin{minipage}[t]{\dimexpr.5\textwidth-4mm\relax}{\scriptsize\color{gray} [p. 25]}\ Lavenex, Sandra. 2022. The cat and mouse game of refugee externalisation policies. Between law and politics \textit{Refugee Externalisation Policies. Responsibility, Legitimacy and Accountability} Azadeh Dastyari, Amy Nethery and Asher Hirsch 27–44. London: Routledge. DOI: 10.4324/9781003167273-3\par\end{minipage}\end{minipage}\par
\noindent\begin{minipage}{\textwidth}\begin{minipage}[t]{\dimexpr.5\textwidth-4mm\relax}{\scriptsize\color{gray} [p. 25]}\ Legomsky, Stephen H. 2003. “Secondary Refugee Movements and the Return of Asylum Seekers to Third Countries: The Meaning of Effective Protection.” \textit{International J. of Refugee Law} 15(4): 567–677. doi:10.1093/ijrl/15.4.567. 〔题名或条目译释：难民的二次流动与将寻求庇护者遣返第三国：有效保护的含义。〕\par\end{minipage}\hfill\begin{minipage}[t]{\dimexpr.5\textwidth-4mm\relax}{\scriptsize\color{gray} [p. 25]}\ Legomsky, Stephen H. 2003. “Secondary Refugee Movements and the Return of Asylum Seekers to Third Countries: The Meaning of Effective Protection.” \textit{International J. of Refugee Law} 15(4): 567–677. doi:10.1093/ijrl/15.4.567.\par\end{minipage}\end{minipage}\par
\noindent\begin{minipage}{\textwidth}\begin{minipage}[t]{\dimexpr.5\textwidth-4mm\relax}{\scriptsize\color{gray} [p. 25]}\ Lejeune, Aude and Spire, Alexis. 2021. “The Role of Legal Intermediaries in the Dispute Pyramid: Inequalities before the French Legal System.” \textit{International J. of Law in Context} 17(4): 455–72. doi:10.1017/S1744552321000513. 〔题名或条目译释：法律中介者在纠纷金字塔中的作用：法国法律体系面前的不平等。〕\par\end{minipage}\hfill\begin{minipage}[t]{\dimexpr.5\textwidth-4mm\relax}{\scriptsize\color{gray} [p. 25]}\ Lejeune, Aude and Spire, Alexis. 2021. “The Role of Legal Intermediaries in the Dispute Pyramid: Inequalities before the French Legal System.” \textit{International J. of Law in Context} 17(4): 455–72. doi:10.1017/S1744552321000513.\par\end{minipage}\end{minipage}\par
\noindent\begin{minipage}{\textwidth}\begin{minipage}[t]{\dimexpr.5\textwidth-4mm\relax}{\scriptsize\color{gray} [p. 25]}\ Mann, Itamar. 2013. “Dialectic of Transnationalism: Unauthorized Migration and Human Rights, 1993–2013.” \textit{Harvard International Law J.} 54(2): 315–91. 〔题名或条目译释：跨国主义的辩证法：未经授权的迁移与人权，1993—2013。〕\par\end{minipage}\hfill\begin{minipage}[t]{\dimexpr.5\textwidth-4mm\relax}{\scriptsize\color{gray} [p. 25]}\ Mann, Itamar. 2013. “Dialectic of Transnationalism: Unauthorized Migration and Human Rights, 1993–2013.” \textit{Harvard International Law J.} 54(2): 315–91.\par\end{minipage}\end{minipage}\par
\noindent\begin{minipage}{\textwidth}\begin{minipage}[t]{\dimexpr.5\textwidth-4mm\relax}{\scriptsize\color{gray} [p. 25]}\ Massimi, Adelaide. 2026. Guida pratica sul contenzioso strategico avverso i respingimenti nel Mediterraneo e per il diritto di ingresso. [\url{https://www.asgi.it/wp-content/uploads/2026/01/guida-respingimenti-mediterraneo_DEF.pdf]} (accessed March 25 , 2026). 〔题名或条目译释：针对地中海推回及争取入境权的战略诉讼实务指南。〕\par\end{minipage}\hfill\begin{minipage}[t]{\dimexpr.5\textwidth-4mm\relax}{\scriptsize\color{gray} [p. 25]}\ Massimi, Adelaide. 2026. Guida pratica sul contenzioso strategico avverso i respingimenti nel Mediterraneo e per il diritto di ingresso. [\url{https://www.asgi.it/wp-content/uploads/2026/01/guida-respingimenti-mediterraneo_DEF.pdf]} (accessed March 25 , 2026).\par\end{minipage}\end{minipage}\par
\noindent\begin{minipage}{\textwidth}\begin{minipage}[t]{\dimexpr.5\textwidth-4mm\relax}{\scriptsize\color{gray} [p. 25]}\ Mather, Lynn and Yngvesson, Barbara. 1981. “Language, Audience, and the Transformation of Disputes.” \textit{Law \& Society Rev.} 15(3–4): 775–821. doi:10.2307/3053512. 〔题名或条目译释：语言、受众与纠纷的转变。〕\par\end{minipage}\hfill\begin{minipage}[t]{\dimexpr.5\textwidth-4mm\relax}{\scriptsize\color{gray} [p. 25]}\ Mather, Lynn and Yngvesson, Barbara. 1981. “Language, Audience, and the Transformation of Disputes.” \textit{Law \& Society Rev.} 15(3–4): 775–821. doi:10.2307/3053512.\par\end{minipage}\end{minipage}\par
\noindent\begin{minipage}{\textwidth}\begin{minipage}[t]{\dimexpr.5\textwidth-4mm\relax}{\scriptsize\color{gray} [p. 25]}\ Miaz, Jonathan, Odasso, Laura and Sabrié, Romane. 2021. “Le droit de la migration et ses intermédiaires: Usages sociopolitiques du droit et production des politiques migratoires. Présentation du dossier.” \textit{Droit et Société} 107(1): 7–15. doi:10.3917/drs1.107.0007. 〔题名或条目译释：移民法及其中介者：法律的社会政治运用与移民政策生产——专题导言。〕\par\end{minipage}\hfill\begin{minipage}[t]{\dimexpr.5\textwidth-4mm\relax}{\scriptsize\color{gray} [p. 25]}\ Miaz, Jonathan, Odasso, Laura and Sabrié, Romane. 2021. “Le droit de la migration et ses intermédiaires: Usages sociopolitiques du droit et production des politiques migratoires. Présentation du dossier.” \textit{Droit et Société} 107(1): 7–15. doi:10.3917/drs1.107.0007.\par\end{minipage}\end{minipage}\par
\noindent\begin{minipage}{\textwidth}\begin{minipage}[t]{\dimexpr.5\textwidth-4mm\relax}{\scriptsize\color{gray} [p. 25]}\ Michelson, Ethan. 2006. “The Practice of Law as an Obstacle to Justice: Chinese Lawyers at Work.” \textit{Law \& Society Rev.} 40(1):1–38. doi:10.1111/j.1540-5893.2006.00257.x. 〔题名或条目译释：作为正义障碍的法律执业：工作中的中国律师。〕\par\end{minipage}\hfill\begin{minipage}[t]{\dimexpr.5\textwidth-4mm\relax}{\scriptsize\color{gray} [p. 25]}\ Michelson, Ethan. 2006. “The Practice of Law as an Obstacle to Justice: Chinese Lawyers at Work.” \textit{Law \& Society Rev.} 40(1):1–38. doi:10.1111/j.1540-5893.2006.00257.x.\par\end{minipage}\end{minipage}\par
\noindent\begin{minipage}{\textwidth}\begin{minipage}[t]{\dimexpr.5\textwidth-4mm\relax}{\scriptsize\color{gray} [p. 25]}\ Miller, Richard E. and Sarat, Austin. 1980–1981. “Grievances, Claims, and Disputes: Assessing the Adversary Culture.” \textit{Law \& Society Rev}. 15: 525–66. doi:10.2307/3053502. 〔题名或条目译释：不满、索求与纠纷：评估对抗文化。〕\par\end{minipage}\hfill\begin{minipage}[t]{\dimexpr.5\textwidth-4mm\relax}{\scriptsize\color{gray} [p. 25]}\ Miller, Richard E. and Sarat, Austin. 1980–1981. “Grievances, Claims, and Disputes: Assessing the Adversary Culture.” \textit{Law \& Society Rev}. 15: 525–66. doi:10.2307/3053502.\par\end{minipage}\end{minipage}\par
\noindent\begin{minipage}{\textwidth}\begin{minipage}[t]{\dimexpr.5\textwidth-4mm\relax}{\scriptsize\color{gray} [p. 25]}\ Mitrović, Marta Stojić and Vilenica, Ana. 2019. “Enforcing and Disrupting Circular Movement in an EU Borderscape: Housingscaping in Serbia.” \textit{Citizenship Studies} 23(6): 540–58. doi:10.1080/13621025.2019.1634368. 〔题名或条目译释：欧盟边境景观中循环移动的强制与扰乱：塞尔维亚的居住空间塑造。〕\par\end{minipage}\hfill\begin{minipage}[t]{\dimexpr.5\textwidth-4mm\relax}{\scriptsize\color{gray} [p. 25]}\ Mitrović, Marta Stojić and Vilenica, Ana. 2019. “Enforcing and Disrupting Circular Movement in an EU Borderscape: Housingscaping in Serbia.” \textit{Citizenship Studies} 23(6): 540–58. doi:10.1080/13621025.2019.1634368.\par\end{minipage}\end{minipage}\par
\noindent\begin{minipage}{\textwidth}\begin{minipage}[t]{\dimexpr.5\textwidth-4mm\relax}{\scriptsize\color{gray} [p. 25]}\ Moraru, Madalina and Reyhani, Adel-Naim. 2026. “Evading EU Law through Summary Returns at Internal Borders: Practice, Legality and the Role of Courts.” \textit{Common Market Law Rev.} 63(2): 413–54. doi:10.54648/COLA2026001. 〔题名或条目译释：通过内部边界的即决送返规避欧盟法：实践、合法性与法院作用。〕\par\end{minipage}\hfill\begin{minipage}[t]{\dimexpr.5\textwidth-4mm\relax}{\scriptsize\color{gray} [p. 25]}\ Moraru, Madalina and Reyhani, Adel-Naim. 2026. “Evading EU Law through Summary Returns at Internal Borders: Practice, Legality and the Role of Courts.” \textit{Common Market Law Rev.} 63(2): 413–54. doi:10.54648/COLA2026001.\par\end{minipage}\end{minipage}\par
\noindent\begin{minipage}{\textwidth}\begin{minipage}[t]{\dimexpr.5\textwidth-4mm\relax}{\scriptsize\color{gray} [p. 25]}\ Moreno-lax, Violeta. 2024. “Meta-Borders and the Rule of Law: From Externalisation to ‘Responsibilisation’ in Systems of Contactless Control.” \textit{Netherlands International Law Rev.} 71: 21–57. doi:10.1007/s40802-024-00257-5. 〔题名或条目译释：元边界与法治：无接触管控体系中从外部化到‘责任化’。〕\par\end{minipage}\hfill\begin{minipage}[t]{\dimexpr.5\textwidth-4mm\relax}{\scriptsize\color{gray} [p. 25]}\ Moreno-lax, Violeta. 2024. “Meta-Borders and the Rule of Law: From Externalisation to ‘Responsibilisation’ in Systems of Contactless Control.” \textit{Netherlands International Law Rev.} 71: 21–57. doi:10.1007/s40802-024-00257-5.\par\end{minipage}\end{minipage}\par
\noindent\begin{minipage}{\textwidth}\begin{minipage}[t]{\dimexpr.5\textwidth-4mm\relax}{\scriptsize\color{gray} [p. 25]}\ Border Violence Monitoring Network. 2024. Annual Report 2024, \url{https://borderviolence.eu/uploads/publication/document/8/Annual_Report_2024.pdf} (accessed March 24 , 2026). 〔题名或条目译释：2024年度报告。〕\par\end{minipage}\hfill\begin{minipage}[t]{\dimexpr.5\textwidth-4mm\relax}{\scriptsize\color{gray} [p. 25]}\ Border Violence Monitoring Network. 2024. Annual Report 2024, \url{https://borderviolence.eu/uploads/publication/document/8/Annual_Report_2024.pdf} (accessed March 24 , 2026).\par\end{minipage}\end{minipage}\par
\noindent\begin{minipage}{\textwidth}\begin{minipage}[t]{\dimexpr.5\textwidth-4mm\relax}{\scriptsize\color{gray} [p. 25]}\ Olesen, Annette and Hammerslev, Ole. 2023. “The Dynamic and Iterative Pre-Dispute Phase: The Transformation from a Justiciable Problem into a Legal Dispute.” \textit{J. of Law and Society} 50(1): 120–38. doi:10.1111/jols.12411. 〔题名或条目译释：动态且反复的纠纷前阶段：从可诉问题转变为法律纠纷。〕\par\end{minipage}\hfill\begin{minipage}[t]{\dimexpr.5\textwidth-4mm\relax}{\scriptsize\color{gray} [p. 25]}\ Olesen, Annette and Hammerslev, Ole. 2023. “The Dynamic and Iterative Pre-Dispute Phase: The Transformation from a Justiciable Problem into a Legal Dispute.” \textit{J. of Law and Society} 50(1): 120–38. doi:10.1111/jols.12411.\par\end{minipage}\end{minipage}\par
\noindent\begin{minipage}{\textwidth}\begin{minipage}[t]{\dimexpr.5\textwidth-4mm\relax}{\scriptsize\color{gray} [p. 25]}\ Passalacqua, Virginia. 2022. “Who Mobilizes the Court? Migrant Rights Defenders Before the Court of Justice of the EU.” \textit{Law and Development Rev.} 15(2): 381–405. doi:10.1515/ldr-2021-0102. 〔题名或条目译释：谁动员法院？欧盟法院面前的移民权利捍卫者。〕\par\end{minipage}\hfill\begin{minipage}[t]{\dimexpr.5\textwidth-4mm\relax}{\scriptsize\color{gray} [p. 25]}\ Passalacqua, Virginia. 2022. “Who Mobilizes the Court? Migrant Rights Defenders Before the Court of Justice of the EU.” \textit{Law and Development Rev.} 15(2): 381–405. doi:10.1515/ldr-2021-0102.\par\end{minipage}\end{minipage}\par
\noindent\begin{minipage}{\textwidth}\begin{minipage}[t]{\dimexpr.5\textwidth-4mm\relax}{\scriptsize\color{gray} [p. 25]}\ Pélisse, Jérôme. 2019. “Varieties of Legal Intermediaries: When Nonlegal Professionals Act as Legal Intermediaries.” In \textit{Legal Intermediation: A Processual Approach to Law and Economic Activity}, edited by Sarat, Austin, 101–28. Bingley: Emerald. DOI: 10.1108/S1059-433720190000081005 〔题名或条目译释：法律中介者的多种形式：当非法律专业人员担任法律中介者；所收文集：法律中介活动——法律与经济活动的过程性研究。〕\par\end{minipage}\hfill\begin{minipage}[t]{\dimexpr.5\textwidth-4mm\relax}{\scriptsize\color{gray} [p. 25]}\ Pélisse, Jérôme. 2019. “Varieties of Legal Intermediaries: When Nonlegal Professionals Act as Legal Intermediaries.” In \textit{Legal Intermediation: A Processual Approach to Law and Economic Activity}, edited by Sarat, Austin, 101–28. Bingley: Emerald. DOI: 10.1108/S1059-433720190000081005\par\end{minipage}\end{minipage}\par
\noindent\begin{minipage}{\textwidth}\begin{minipage}[t]{\dimexpr.5\textwidth-4mm\relax}{\scriptsize\color{gray} [p. 25]}\ Pijnenburg, Annick and Van der Pas, Kris. 2022. “Strategic Litigation against European Migration Control Policies: The Legal Battleground of the Central Mediterranean Migration Route.” \textit{European J. of Migration and Law} 24(3): 401–29. DOI: 10.1163/15718166-12340135 〔题名或条目译释：针对欧洲移民管控政策的战略诉讼：地中海中部迁徙路线的法律战场。〕\par\end{minipage}\hfill\begin{minipage}[t]{\dimexpr.5\textwidth-4mm\relax}{\scriptsize\color{gray} [p. 25]}\ Pijnenburg, Annick and Van der Pas, Kris. 2022. “Strategic Litigation against European Migration Control Policies: The Legal Battleground of the Central Mediterranean Migration Route.” \textit{European J. of Migration and Law} 24(3): 401–29. DOI: 10.1163/15718166-12340135\par\end{minipage}\end{minipage}\par
\noindent\begin{minipage}{\textwidth}\begin{minipage}[t]{\dimexpr.5\textwidth-4mm\relax}{\scriptsize\color{gray} [p. 25]}\ Pörtner, Ephraim. 2021. \textit{Re-Cording Lives: Governing Asylum in Switzerland and the Need to Resolve}. Bielefeld: transcript Verlag. DOI: 10.14361/9783839453490 〔题名或条目译释：重新记录生命：瑞士的庇护治理与作出裁断的需要。〕\par\end{minipage}\hfill\begin{minipage}[t]{\dimexpr.5\textwidth-4mm\relax}{\scriptsize\color{gray} [p. 25]}\ Pörtner, Ephraim. 2021. \textit{Re-Cording Lives: Governing Asylum in Switzerland and the Need to Resolve}. Bielefeld: transcript Verlag. DOI: 10.14361/9783839453490\par\end{minipage}\end{minipage}\par
\noindent\begin{minipage}{\textwidth}\begin{minipage}[t]{\dimexpr.5\textwidth-4mm\relax}{\scriptsize\color{gray} [p. 25]}\ Reyhani, Adel-Naim. 2021. “Anomaly upon Anomaly: The 1951 Convention and State Disintegration.” \textit{International J. of Refugee Law} 33(2): 277–99. doi:10.1093/ijrl/eeab042. 〔题名或条目译释：反常之上的反常：1951年公约与国家解体。〕\par\end{minipage}\hfill\begin{minipage}[t]{\dimexpr.5\textwidth-4mm\relax}{\scriptsize\color{gray} [p. 25]}\ Reyhani, Adel-Naim. 2021. “Anomaly upon Anomaly: The 1951 Convention and State Disintegration.” \textit{International J. of Refugee Law} 33(2): 277–99. doi:10.1093/ijrl/eeab042.\par\end{minipage}\end{minipage}\par
\noindent\begin{minipage}{\textwidth}\begin{minipage}[t]{\dimexpr.5\textwidth-4mm\relax}{\scriptsize\color{gray} [p. 26]}\ Riva, Sara and Hoffstaedter, Gerhard. 2021. “The Aporia of Refugee Rights in a Time of Crises: The Role of Brokers in Accessing Refugee Protection in Transit and at the Border.” \textit{Comparative Migration Studies} 9(1): 1–17. DOI: 10.1186/s40878-020-00212-2 〔题名或条目译释：危机时期难民权利的困境：经纪者在过境途中及边境获取难民保护中的作用。〕\par\end{minipage}\hfill\begin{minipage}[t]{\dimexpr.5\textwidth-4mm\relax}{\scriptsize\color{gray} [p. 26]}\ Riva, Sara and Hoffstaedter, Gerhard. 2021. “The Aporia of Refugee Rights in a Time of Crises: The Role of Brokers in Accessing Refugee Protection in Transit and at the Border.” \textit{Comparative Migration Studies} 9(1): 1–17. DOI: 10.1186/s40878-020-00212-2\par\end{minipage}\end{minipage}\par
\noindent\begin{minipage}{\textwidth}\begin{minipage}[t]{\dimexpr.5\textwidth-4mm\relax}{\scriptsize\color{gray} [p. 26]}\ Shachar, Ayelet. 2020. \textit{The Shifting Border: Legal Cartographies of Migration and Mobility}. Manchester: Manchester University Press. DOI: 10.7765/9781526145321 〔题名或条目译释：移动的边界：迁移与流动的法律制图。〕\par\end{minipage}\hfill\begin{minipage}[t]{\dimexpr.5\textwidth-4mm\relax}{\scriptsize\color{gray} [p. 26]}\ Shachar, Ayelet. 2020. \textit{The Shifting Border: Legal Cartographies of Migration and Mobility}. Manchester: Manchester University Press. DOI: 10.7765/9781526145321\par\end{minipage}\end{minipage}\par
\noindent\begin{minipage}{\textwidth}\begin{minipage}[t]{\dimexpr.5\textwidth-4mm\relax}{\scriptsize\color{gray} [p. 26]}\ Southworth, Ann. 2000. “The Rights Revolution and Support Structures for Rights Advocacy.” \textit{Law \& Society Rev.} 34(4): 1203–19. DOI: 10.2307/3115134 〔题名或条目译释：权利革命与权利倡导的支持结构。〕\par\end{minipage}\hfill\begin{minipage}[t]{\dimexpr.5\textwidth-4mm\relax}{\scriptsize\color{gray} [p. 26]}\ Southworth, Ann. 2000. “The Rights Revolution and Support Structures for Rights Advocacy.” \textit{Law \& Society Rev.} 34(4): 1203–19. DOI: 10.2307/3115134\par\end{minipage}\end{minipage}\par
\noindent\begin{minipage}{\textwidth}\begin{minipage}[t]{\dimexpr.5\textwidth-4mm\relax}{\scriptsize\color{gray} [p. 26]}\ Sterett, Susan M. 2016. “Legal Mobilization and Juridification: Migration as a Central Case.” \textit{Law \& Policy} 38(4): 273–79. DOI: 10.1111/lapo.12063 〔题名或条目译释：法律动员与司法化：作为核心案例的移民。〕\par\end{minipage}\hfill\begin{minipage}[t]{\dimexpr.5\textwidth-4mm\relax}{\scriptsize\color{gray} [p. 26]}\ Sterett, Susan M. 2016. “Legal Mobilization and Juridification: Migration as a Central Case.” \textit{Law \& Policy} 38(4): 273–79. DOI: 10.1111/lapo.12063\par\end{minipage}\end{minipage}\par
\noindent\begin{minipage}{\textwidth}\begin{minipage}[t]{\dimexpr.5\textwidth-4mm\relax}{\scriptsize\color{gray} [p. 26]}\ Talesh, Shauhin A. 2015. “Legal Intermediaries: How Insurance Companies Construct the Meaning of Compliance with Anti-Discrimination Laws.” \textit{Law \& Policy} 37(3): 209–39. DOI: 10.1111/lapo.12037 〔题名或条目译释：法律中介者：保险公司如何建构遵守反歧视法的意义。〕\par\end{minipage}\hfill\begin{minipage}[t]{\dimexpr.5\textwidth-4mm\relax}{\scriptsize\color{gray} [p. 26]}\ Talesh, Shauhin A. 2015. “Legal Intermediaries: How Insurance Companies Construct the Meaning of Compliance with Anti-Discrimination Laws.” \textit{Law \& Policy} 37(3): 209–39. DOI: 10.1111/lapo.12037\par\end{minipage}\end{minipage}\par
\noindent\begin{minipage}{\textwidth}\begin{minipage}[t]{\dimexpr.5\textwidth-4mm\relax}{\scriptsize\color{gray} [p. 26]}\ Tan, Nikolas Feith and Gammeltoft-Hansen, Thomas. 2020. “A Topographical Approach to Accountability for Human Rights Violations in Migration Control.” \textit{German Law J.} 21(3): 335–54. doi:10.1017/glj.2020.31. 〔题名或条目译释：移民管控中人权侵害问责的地形学研究路径。〕\par\end{minipage}\hfill\begin{minipage}[t]{\dimexpr.5\textwidth-4mm\relax}{\scriptsize\color{gray} [p. 26]}\ Tan, Nikolas Feith and Gammeltoft-Hansen, Thomas. 2020. “A Topographical Approach to Accountability for Human Rights Violations in Migration Control.” \textit{German Law J.} 21(3): 335–54. doi:10.1017/glj.2020.31.\par\end{minipage}\end{minipage}\par
\noindent\begin{minipage}{\textwidth}\begin{minipage}[t]{\dimexpr.5\textwidth-4mm\relax}{\scriptsize\color{gray} [p. 26]}\ Tazzioli, Martina. 2021. “‘Choking without killing’: Opacity and the grey area of migration governmentality.” \textit{Political Geography} 89: 102412. doi:10.1016/j.polgeo.2021.102412. 〔题名或条目译释：‘使人窒息而不杀人’：移民治理术的不透明性与灰色地带。〕\par\end{minipage}\hfill\begin{minipage}[t]{\dimexpr.5\textwidth-4mm\relax}{\scriptsize\color{gray} [p. 26]}\ Tazzioli, Martina. 2021. “‘Choking without killing’: Opacity and the grey area of migration governmentality.” \textit{Political Geography} 89: 102412. doi:10.1016/j.polgeo.2021.102412.\par\end{minipage}\end{minipage}\par
\noindent\begin{minipage}{\textwidth}\begin{minipage}[t]{\dimexpr.5\textwidth-4mm\relax}{\scriptsize\color{gray} [p. 26]}\ Thomas, Gammeltoft-Hansen and Hathaway, James C.. 2014. “Non-Refoulement in a World of Cooperative Deterrence.” \textit{Columbia J. of Transnational Law} 53(2): 235–84. 〔题名或条目译释：合作威慑世界中的不推回原则。〕\par\end{minipage}\hfill\begin{minipage}[t]{\dimexpr.5\textwidth-4mm\relax}{\scriptsize\color{gray} [p. 26]}\ Thomas, Gammeltoft-Hansen and Hathaway, James C.. 2014. “Non-Refoulement in a World of Cooperative Deterrence.” \textit{Columbia J. of Transnational Law} 53(2): 235–84.\par\end{minipage}\end{minipage}\par
\noindent\begin{minipage}{\textwidth}\begin{minipage}[t]{\dimexpr.5\textwidth-4mm\relax}{\scriptsize\color{gray} [p. 26]}\ Timmermans, Stefan and Tavory, Iddo. 2012. “Theory Construction in Qualitative Research: From Grounded Theory to Abductive Analysis.” \textit{Sociological Theory} 30(3): 167–86. doi:10.1177/0735275112457914. 〔题名或条目译释：定性研究中的理论建构：从扎根理论到溯因分析。〕\par\end{minipage}\hfill\begin{minipage}[t]{\dimexpr.5\textwidth-4mm\relax}{\scriptsize\color{gray} [p. 26]}\ Timmermans, Stefan and Tavory, Iddo. 2012. “Theory Construction in Qualitative Research: From Grounded Theory to Abductive Analysis.” \textit{Sociological Theory} 30(3): 167–86. doi:10.1177/0735275112457914.\par\end{minipage}\end{minipage}\par
\noindent\begin{minipage}{\textwidth}\begin{minipage}[t]{\dimexpr.5\textwidth-4mm\relax}{\scriptsize\color{gray} [p. 26]}\ Trucco, Daniela. 2023. “The Criminalization of Border Solidarity and the Law as a Double-Edged Weapon.” \textit{Partecipazione \& Conflitto} 16(3): 470–89. 〔题名或条目译释：边境声援的刑事化与作为双刃剑的法律。〕\par\end{minipage}\hfill\begin{minipage}[t]{\dimexpr.5\textwidth-4mm\relax}{\scriptsize\color{gray} [p. 26]}\ Trucco, Daniela. 2023. “The Criminalization of Border Solidarity and the Law as a Double-Edged Weapon.” \textit{Partecipazione \& Conflitto} 16(3): 470–89.\par\end{minipage}\end{minipage}\par
\noindent\begin{minipage}{\textwidth}\begin{minipage}[t]{\dimexpr.5\textwidth-4mm\relax}{\scriptsize\color{gray} [p. 26]}\ Trucco, Daniela, Lamarche, Karine and Philippe, Oriana. 2023. “À la frontière du droit: Usages du droit en solidarité avec les migrants aux frontières.” \textit{Droit et Société} 113(1): 181–202. doi:10.3917/drs1.113.0181. 〔题名或条目译释：在法律的边界：边境移民声援中的法律运用。〕\par\end{minipage}\hfill\begin{minipage}[t]{\dimexpr.5\textwidth-4mm\relax}{\scriptsize\color{gray} [p. 26]}\ Trucco, Daniela, Lamarche, Karine and Philippe, Oriana. 2023. “À la frontière du droit: Usages du droit en solidarité avec les migrants aux frontières.” \textit{Droit et Société} 113(1): 181–202. doi:10.3917/drs1.113.0181.\par\end{minipage}\end{minipage}\par
\noindent\begin{minipage}{\textwidth}\begin{minipage}[t]{\dimexpr.5\textwidth-4mm\relax}{\scriptsize\color{gray} [p. 26]}\ Tuckett, Anna. 2018. “Ethical Brokerage and Self-Fashioning in Italian Immigration Bureaucracy.” \textit{Critique of Anthropology} 38(3): 245–64. doi:10.1177/0308275X18775199. 〔题名或条目译释：意大利移民官僚体系中的伦理经纪与自我塑造。〕\par\end{minipage}\hfill\begin{minipage}[t]{\dimexpr.5\textwidth-4mm\relax}{\scriptsize\color{gray} [p. 26]}\ Tuckett, Anna. 2018. “Ethical Brokerage and Self-Fashioning in Italian Immigration Bureaucracy.” \textit{Critique of Anthropology} 38(3): 245–64. doi:10.1177/0308275X18775199.\par\end{minipage}\end{minipage}\par
\noindent\begin{minipage}{\textwidth}\begin{minipage}[t]{\dimexpr.5\textwidth-4mm\relax}{\scriptsize\color{gray} [p. 26]}\ UN Special Rapporteur on the human rights of migrants. 2021. “Report on means to address the human rights impact of pushbacks of migrants on land and at sea. Report of the Special Rapporteur on the human rights of migrants.” A/HRC/47/30, 12 May 2021. 〔题名或条目译释：应对陆上及海上推回移民对人权影响的手段报告——移民人权问题特别报告员的报告。〕\par\end{minipage}\hfill\begin{minipage}[t]{\dimexpr.5\textwidth-4mm\relax}{\scriptsize\color{gray} [p. 26]}\ UN Special Rapporteur on the human rights of migrants. 2021. “Report on means to address the human rights impact of pushbacks of migrants on land and at sea. Report of the Special Rapporteur on the human rights of migrants.” A/HRC/47/30, 12 May 2021.\par\end{minipage}\end{minipage}\par
\noindent\begin{minipage}{\textwidth}\begin{minipage}[t]{\dimexpr.5\textwidth-4mm\relax}{\scriptsize\color{gray} [p. 26]}\ Van der Pas, Kris. 2023. “Legal Mobilization in the Field of Asylum Law: A Revival of Political Opportunity Structures?” \textit{Recht der Werkelijkheid} 44(3): 14–31. doi:10.5553/RdW/138064242023044003002. 〔题名或条目译释：庇护法领域的法律动员：政治机会结构的复兴？。〕\par\end{minipage}\hfill\begin{minipage}[t]{\dimexpr.5\textwidth-4mm\relax}{\scriptsize\color{gray} [p. 26]}\ Van der Pas, Kris. 2023. “Legal Mobilization in the Field of Asylum Law: A Revival of Political Opportunity Structures?” \textit{Recht der Werkelijkheid} 44(3): 14–31. doi:10.5553/RdW/138064242023044003002.\par\end{minipage}\end{minipage}\par
\noindent\begin{minipage}{\textwidth}\begin{minipage}[t]{\dimexpr.5\textwidth-4mm\relax}{\scriptsize\color{gray} [p. 26]}\ Vanhala, Lisa. 2012. “Legal Opportunity Structures and the Paradox of Legal Mobilization by the Environmental Movement in the UK.” \textit{Law \& Society Rev}. 46(3): 523–56. doi:10.1111/j.1540-5893.2012.00505.x. 〔题名或条目译释：法律机会结构与英国环境运动法律动员的悖论。〕\par\end{minipage}\hfill\begin{minipage}[t]{\dimexpr.5\textwidth-4mm\relax}{\scriptsize\color{gray} [p. 26]}\ Vanhala, Lisa. 2012. “Legal Opportunity Structures and the Paradox of Legal Mobilization by the Environmental Movement in the UK.” \textit{Law \& Society Rev}. 46(3): 523–56. doi:10.1111/j.1540-5893.2012.00505.x.\par\end{minipage}\end{minipage}\par
\noindent\begin{minipage}{\textwidth}\begin{minipage}[t]{\dimexpr.5\textwidth-4mm\relax}{\scriptsize\color{gray} [p. 26]}\ Vanhala, Lisa. 2018. “Is Legal Mobilization for the Birds? Legal Opportunity Structures and Environmental Nongovernmental Organizations in the United Kingdom, France, Finland, and Italy.” \textit{Comparative Political Studies} 51(3): 380–412. doi:10.1177/0010414017710257. 〔题名或条目译释：法律动员有用吗？英国、法国、芬兰与意大利的法律机会结构及环境非政府组织（题名含 for the birds 的双关）。〕\par\end{minipage}\hfill\begin{minipage}[t]{\dimexpr.5\textwidth-4mm\relax}{\scriptsize\color{gray} [p. 26]}\ Vanhala, Lisa. 2018. “Is Legal Mobilization for the Birds? Legal Opportunity Structures and Environmental Nongovernmental Organizations in the United Kingdom, France, Finland, and Italy.” \textit{Comparative Political Studies} 51(3): 380–412. doi:10.1177/0010414017710257.\par\end{minipage}\end{minipage}\par
\noindent\begin{minipage}{\textwidth}\begin{minipage}[t]{\dimexpr.5\textwidth-4mm\relax}{\scriptsize\color{gray} [p. 26]}\ Zocchi, Benedetta. 2024. “The Game: Ritualized Exhaustion and Subversion on the Western Balkan Route.” \textit{J. of Borderlands Studies} 39(5): 875–95. doi:10.1080/08865655.2023.2229846. 〔题名或条目译释：游戏：西巴尔干路线上的仪式化耗竭与颠覆。〕\par\end{minipage}\hfill\begin{minipage}[t]{\dimexpr.5\textwidth-4mm\relax}{\scriptsize\color{gray} [p. 26]}\ Zocchi, Benedetta. 2024. “The Game: Ritualized Exhaustion and Subversion on the Western Balkan Route.” \textit{J. of Borderlands Studies} 39(5): 875–95. doi:10.1080/08865655.2023.2229846.\par\end{minipage}\end{minipage}\par
\noindent\begin{minipage}{\textwidth}\begin{minipage}[t]{\dimexpr.5\textwidth-4mm\relax}{\scriptsize\color{gray} [p. 26]}\ Zoppi, Marco and Puleri, Massimiliano. 2022. “The Balkan Route (and Its Afterlife): The New Normal in the European Politics of Migration.” \textit{J. of Balkan and Near Eastern Studies} 24(3): 576–93 doi:10.1080/19448953.2021.2015658. 〔题名或条目译释：巴尔干路线及其后续生命：欧洲移民政治的新常态。〕\par\end{minipage}\hfill\begin{minipage}[t]{\dimexpr.5\textwidth-4mm\relax}{\scriptsize\color{gray} [p. 26]}\ Zoppi, Marco and Puleri, Massimiliano. 2022. “The Balkan Route (and Its Afterlife): The New Normal in the European Politics of Migration.” \textit{J. of Balkan and Near Eastern Studies} 24(3): 576–93 doi:10.1080/19448953.2021.2015658.\par\end{minipage}\end{minipage}\par
\Needspace{10\baselineskip}\section*{所引判例 / Cases cited}
\Needspace{10\baselineskip}\section*{意大利： / Italy:.}
\noindent\begin{minipage}{\textwidth}\begin{minipage}[t]{\dimexpr.5\textwidth-4mm\relax}{\scriptsize\color{gray} [p. 26]}\ Tribunale di Roma, N.R.G. 56420/2020, 18 January 2021a. 〔题名或条目译释：罗马法院，案号N.R.G. 56420/2020，2021年1月18日（文内标识2021a）。〕\par\end{minipage}\hfill\begin{minipage}[t]{\dimexpr.5\textwidth-4mm\relax}{\scriptsize\color{gray} [p. 26]}\ Tribunale di Roma, N.R.G. 56420/2020, 18 January 2021a.\par\end{minipage}\end{minipage}\par
\noindent\begin{minipage}{\textwidth}\begin{minipage}[t]{\dimexpr.5\textwidth-4mm\relax}{\scriptsize\color{gray} [p. 26]}\ Tribunale di Roma, N.R.G. 7045/2021, 27th April 2021b. 〔题名或条目译释：罗马法院，案号N.R.G. 7045/2021，2021年4月27日（文内标识2021b）。〕\par\end{minipage}\hfill\begin{minipage}[t]{\dimexpr.5\textwidth-4mm\relax}{\scriptsize\color{gray} [p. 26]}\ Tribunale di Roma, N.R.G. 7045/2021, 27th April 2021b.\par\end{minipage}\end{minipage}\par
\noindent\begin{minipage}{\textwidth}\begin{minipage}[t]{\dimexpr.5\textwidth-4mm\relax}{\scriptsize\color{gray} [p. 26]}\ Tribunale di Roma, N.R.G. 3938/2022, 9th May 2023. 〔题名或条目译释：罗马法院，案号N.R.G. 3938/2022，2023年5月9日。〕\par\end{minipage}\hfill\begin{minipage}[t]{\dimexpr.5\textwidth-4mm\relax}{\scriptsize\color{gray} [p. 26]}\ Tribunale di Roma, N.R.G. 3938/2022, 9th May 2023.\par\end{minipage}\end{minipage}\par
\Needspace{10\baselineskip}\section*{奥地利： / Austria:.}
\noindent\begin{minipage}{\textwidth}\begin{minipage}[t]{\dimexpr.5\textwidth-4mm\relax}{\scriptsize\color{gray} [p. 26]}\ Verwaltungsgerichtshof (VwGH), Ra 2021/21/0274, 5 May 2022a. 〔题名或条目译释：奥地利行政法院（VwGH），案号Ra 2021/21/0274，2022年5月5日（文内标识2022a）。〕\par\end{minipage}\hfill\begin{minipage}[t]{\dimexpr.5\textwidth-4mm\relax}{\scriptsize\color{gray} [p. 26]}\ Verwaltungsgerichtshof (VwGH), Ra 2021/21/0274, 5 May 2022a.\par\end{minipage}\end{minipage}\par
\noindent\begin{minipage}{\textwidth}\begin{minipage}[t]{\dimexpr.5\textwidth-4mm\relax}{\scriptsize\color{gray} [p. 27]}\ Verwaltungsgerichtshof (VwGH), Ra 2022/21/0074, 19 May 2022b. 〔题名或条目译释：奥地利行政法院（VwGH），案号Ra 2022/21/0074，2022年5月19日（文内标识2022b）。〕\par\end{minipage}\hfill\begin{minipage}[t]{\dimexpr.5\textwidth-4mm\relax}{\scriptsize\color{gray} [p. 27]}\ Verwaltungsgerichtshof (VwGH), Ra 2022/21/0074, 19 May 2022b.\par\end{minipage}\end{minipage}\par
\noindent\begin{minipage}{\textwidth}\begin{minipage}[t]{\dimexpr.5\textwidth-4mm\relax}{\scriptsize\color{gray} [p. 27]}\ Landesverwaltungsgericht Steiermark (LVwG), 20.3-2725/2020-86, 1 July 2021. 〔题名或条目译释：施泰尔马克州行政法院（LVwG），案号20.3-2725/2020-86，2021年7月1日。〕\par\end{minipage}\hfill\begin{minipage}[t]{\dimexpr.5\textwidth-4mm\relax}{\scriptsize\color{gray} [p. 27]}\ Landesverwaltungsgericht Steiermark (LVwG), 20.3-2725/2020-86, 1 July 2021.\par\end{minipage}\end{minipage}\par
\noindent\begin{minipage}{\textwidth}\begin{minipage}[t]{\dimexpr.5\textwidth-4mm\relax}{\scriptsize\color{gray} [p. 27]}\ Landesverwaltungsgericht Steiermark (LVwG), 20.3-2621/2021-49, 16 February 2022. 〔题名或条目译释：施泰尔马克州行政法院（LVwG），案号20.3-2621/2021-49，2022年2月16日。〕\par\end{minipage}\hfill\begin{minipage}[t]{\dimexpr.5\textwidth-4mm\relax}{\scriptsize\color{gray} [p. 27]}\ Landesverwaltungsgericht Steiermark (LVwG), 20.3-2621/2021-49, 16 February 2022.\par\end{minipage}\end{minipage}\par
\Needspace{10\baselineskip}\section*{所引法规 / Statutes cited}
\noindent\begin{minipage}{\textwidth}\begin{minipage}[t]{\dimexpr.5\textwidth-4mm\relax}{\scriptsize\color{gray} [p. 27]}\ Austria. 2005. Federal Act Concerning the Granting of Asylum (Asylgesetz 2005; AsylG 2005). Federal Law Gazette I No. 100/2005, section 15. 〔题名或条目译释：奥地利，《授予庇护联邦法》（2005年庇护法；AsylG 2005），《联邦法律公报》第一部分第100/2005号，第15条。〕\par\end{minipage}\hfill\begin{minipage}[t]{\dimexpr.5\textwidth-4mm\relax}{\scriptsize\color{gray} [p. 27]}\ Austria. 2005. Federal Act Concerning the Granting of Asylum (Asylgesetz 2005; AsylG 2005). Federal Law Gazette I No. 100/2005, section 15.\par\end{minipage}\end{minipage}\par
\noindent\begin{minipage}{\textwidth}\begin{minipage}[t]{\dimexpr.5\textwidth-4mm\relax}{\scriptsize\color{gray} [p. 27]}\ Italy. 1947. Constitution of the Italian Republic. Gazzetta Ufficiale della Repubblica Italiana, no. 298, 27 December 1947 (edizione straordinaria). 〔题名或条目译释：意大利，《意大利共和国宪法》，《意大利共和国官方公报》第298号，1947年12月27日（特别版）。〕\par\end{minipage}\hfill\begin{minipage}[t]{\dimexpr.5\textwidth-4mm\relax}{\scriptsize\color{gray} [p. 27]}\ Italy. 1947. Constitution of the Italian Republic. Gazzetta Ufficiale della Repubblica Italiana, no. 298, 27 December 1947 (edizione straordinaria).\par\end{minipage}\end{minipage}\par
\noindent\begin{minipage}{\textwidth}\begin{minipage}[t]{\dimexpr.5\textwidth-4mm\relax}{\scriptsize\color{gray} [p. 27]}\ Italy. 1940. Codice di procedura civile. Regio Decreto 28 October 1940, no. 1443, art. 700. 〔题名或条目译释：意大利，《民事诉讼法典》，1940年10月28日第1443号王室法令，第700条。〕\par\end{minipage}\hfill\begin{minipage}[t]{\dimexpr.5\textwidth-4mm\relax}{\scriptsize\color{gray} [p. 27]}\ Italy. 1940. Codice di procedura civile. Regio Decreto 28 October 1940, no. 1443, art. 700.\par\end{minipage}\end{minipage}\par
\Needspace{10\baselineskip}\section*{附录（访谈清单） / Appendix (list of interviews)}
\begin{paracol}{2}
{\scriptsize\color{gray} [p. 27]}\ 访谈1。\textit{活跃于全国性法律协会、在移民领域开展战略诉讼，并担任当地非政府组织顾问的律师}，意大利，2025年4月7日。\par
\switchcolumn
{\scriptsize\color{gray} [p. 27]}\ Interview 1. \textit{Lawyer active in a national legal association, pursuing strategic litigation in the migration field, and a consultant for a local NGO}, Italy, 7 April 2025.\par
\switchcolumn*
{\scriptsize\color{gray} [p. 27]}\ 访谈2。\textit{当地非政府组织行动者，参与巴尔干路线沿线的权益倡导及监测}，意大利，2025年4月14日。\par
\switchcolumn
{\scriptsize\color{gray} [p. 27]}\ Interview 2. \textit{Activist in a local NGO involved in advocacy and monitoring along the Balkan route}, Italy, 14 April 2025.\par
\switchcolumn*
{\scriptsize\color{gray} [p. 27]}\ 访谈3。\textit{活跃于全国性法律协会、在移民领域开展战略诉讼的律师}，意大利，2025年4月28日。\par
\switchcolumn
{\scriptsize\color{gray} [p. 27]}\ Interview 3. \textit{Lawyer active in a national legal association, pursuing strategic litigation in the migration field.} Italy, 28 April 2025.\par
\switchcolumn*
{\scriptsize\color{gray} [p. 27]}\ 访谈4。\textit{为数家国际媒体工作的记者兼摄影记者}，意大利，5月5日。\par
\switchcolumn
{\scriptsize\color{gray} [p. 27]}\ Interview 4. \textit{Journalist and photo-reporter working for several international media}, Italy, 5 May.\par
\switchcolumn*
{\scriptsize\color{gray} [p. 27]}\ 访谈5。两名\textit{建立全国热线、为遭遇推回者提供首次联系的行动者}，奥地利，2025年4月29日。\par
\switchcolumn
{\scriptsize\color{gray} [p. 27]}\ Interview 5. Two \textit{activists who established a national hotline for first contact in case of pushbacks}, Austria, 29 April 2025.\par
\switchcolumn*
{\scriptsize\color{gray} [p. 27]}\ 访谈6。\textit{专长为庇护与移民事务的律师兼全国性非政府组织协调员}，奥地利，2025年7月23日。\par
\switchcolumn
{\scriptsize\color{gray} [p. 27]}\ Interview 6. \textit{Lawyer and national NGO coordinator, specialized in asylum and migration}, Austria, 23 July 2025.\par
\switchcolumn*
{\scriptsize\color{gray} [p. 27]}\ 访谈7。\textit{地区行政法院前法官}，奥地利，2024年6月22日。\par
\switchcolumn
{\scriptsize\color{gray} [p. 27]}\ Interview 7. \textit{Former judge in a regional administrative court}, Austria, 22 June 2024.\par
\switchcolumn*
{\scriptsize\color{gray} [p. 27]}\ 访谈8。\textit{专长为庇护、移民及警察行为的律师与诉讼代理律师}，奥地利，2024年6月24日。\par
\switchcolumn
{\scriptsize\color{gray} [p. 27]}\ Interview 8. \textit{Lawyer and attorney specialized in asylum, migration, and police conduct}, Austria, 24 June 2024.\par
\switchcolumn*
{\scriptsize\color{gray} [p. 27]}\ 访谈9。\textit{非政府组织内专长为庇护与移民事务的律师}，斯洛文尼亚，2025年5月23日。\par
\switchcolumn
{\scriptsize\color{gray} [p. 27]}\ Interview 9. \textit{Lawyer, within an NGO, specialized in asylum and migration}, Slovenia, 23 May 2025.\par
\switchcolumn*
\end{paracol}
\Needspace{10\baselineskip}\section*{作者简介 / About the authors}
\begin{paracol}{2}
{\scriptsize\color{gray} [p. 27]}\ Alice Lacchei（通讯作者）是博洛尼亚大学政治与社会科学系博士后研究员，在欧洲研究理事会启动基金项目“国际难民法的守门人？——法院在塑造庇护准入中的作用”中，研究法院在塑造希腊庇护准入方面的作用。其研究兴趣包括移民与法院的法律与社会研究、移民政策执行，以及移民政策的司法化。alice.lacchei2@unibo.it。\par
\switchcolumn
{\scriptsize\color{gray} [p. 27]}\ Alice Lacchei (corresponding author) is a postdoctoral researcher at the Department of Political and Social Sciences of the University of Bologna within the ERC Starting Grant “Gatekeepers to international refugee law? – The role of courts in shaping access to asylum” project, where she conducts research on the role of courts in shaping access to asylum in Greece. Her research interests include socio-legal research on migration and courts, migration policy implementation, and judicialization of migration policy. alice.lacchei2@unibo.it.\par
\switchcolumn*
{\scriptsize\color{gray} [p. 27]}\ Adel-Naim Reyhani是路德维希·玻尔兹曼基本与人权研究所（LBI-GMR）博士后研究员、“外部化与无权利状态”研究项目负责人，以及该研究所参与欧洲研究理事会启动基金ACCESS项目部分的科研负责人。adel-naim.reyhani@univie.ac.at。\par
\switchcolumn
{\scriptsize\color{gray} [p. 27]}\ Adel-Naim Reyhani is a postdoctoral researcher at the Ludwig Boltzmann Institute of Fundamental and Human Rights (LBI-GMR), PI of the research project “Externalization and Rightlessness,” and scientific lead for the contribution of the LBI-GMR to the ERC Starting Grant ACCESS project. adel-naim.reyhani@univie.ac.at.\par
\switchcolumn*
\end{paracol}
\section*{编者说明}
pushback 译为推回，指原文界定的未经事先评估风险或保护需求而强制移出领土；non-refoulement 单独译为不推回原则。readmission 译为重新接收；case misalignment 译为案件错位；case assembly 译为案件组装；legal intermediaries 译为法律中介者。hostile setting 译为充满阻碍的环境，以表达文中对诉讼不利的制度与实际条件。\par
claimant 采用诉请人，具体引文中的 complainant 译为申诉人，applicant 译为申请人。原文在不同国内程序中使用 appeal，中文依语境译为上诉或提出诉请，避免把所有首次向法院提出的程序都误当作上诉。mandate 与 power of attorney 按语境译为授权或授权委托书；apostille 译为附加证明书。\par
英文以出版方 HTML 为底本，斜体照录。书目区的 Cases cited、Statutes cited 分组标题已恢复，Italy 与 Austria 按国家小标题处理；共91条文献、7项判例、3项法规。附录按照原刊 PDF 的顺序放在书目、判例和法规之后。两位作者简介从原刊第27页补回；‘作者简介’为本版添加的组织性标题。页码指原刊段落或条目的起始页，跨页内容不拆分。\par
书目保留原始检索信息，并附题名或条目译释；译释不意味着存在正式中译本。HTML 中与前文粘连的 DOI 补加分隔符。HTML 与 PDF 中部分合著者姓名次序不同，本版书目统一按 HTML 照录，不据此更改文内引文。\par
原文引言有重复的 Case studies、Pélisse 2019, 2019 及多余引号。英文保留；中文结构介绍按句意表达，重复年份仍保留。正文中的 Lavenex 2023 与书目中的2022未自行统一；Italian Ministry of the Interior 2021 条目内又写 June 2020，均照录。\par
访谈3引文出现 readmitted [the claimant] times 的异常占位形式，本译文保留为‘重新接收[诉请人]次’，不自行推定次数。访谈2引文将地点写作 Biaz，其他段落多为 Bihać；引文提及在的里雅斯特的月份为9月，前文个案叙述另有7月及10月，均保留原文写法。\par
附录访谈4仅列5月5日、未列年份，访谈7和8列2024年，其余多为2025年；本版不替作者补年份或统一日期。作者未访谈诉请人，因此有关诉请人能动性的材料为间接叙述；译文完整保留这一方法限制。数据可用性声明中的 on paper 按‘以纸面形式提供’译出，不推定为公开数据链接。\par
\end{document}
